\documentclass[11pt,a4paper]{article}

% =============================================================================
% Packages
% =============================================================================
\usepackage[margin=2.5cm]{geometry}
\usepackage{tabularx}
\usepackage{longtable}
\usepackage{booktabs}
\usepackage{xcolor}
\usepackage{listings}
\usepackage{hyperref}
\usepackage{amsmath}
\usepackage{enumitem}
\usepackage{parskip}
\usepackage[skins,breakable]{tcolorbox}

\definecolor{codebg}{rgb}{0.95,0.95,0.95}
\definecolor{notebg}{rgb}{0.93,0.96,0.99}
\definecolor{notebar}{rgb}{0.20,0.45,0.75}

% -----------------------------------------------------------------------------
% Code / inline macros
% -----------------------------------------------------------------------------
\newcommand{\code}[1]{\texttt{#1}}

% -----------------------------------------------------------------------------
% MATLAB code listing style
% -----------------------------------------------------------------------------
\lstdefinestyle{matlabstyle}{
  backgroundcolor=\color{codebg},
  basicstyle=\ttfamily\small,
  breaklines=true,
  showstringspaces=false,
  frame=single,
  rulecolor=\color{gray!40},
  xleftmargin=4pt, xrightmargin=4pt,
  aboveskip=6pt, belowskip=6pt
}
\lstset{style=matlabstyle}

% -----------------------------------------------------------------------------
% Option/field table environments (plain tabular + booktabs)
% NOTE: tabularx's X column cannot have its \end{tabularx} hidden inside a
% wrapper \newenvironment (its width-measurement pre-pass cannot see through
% the macro expansion). We therefore use plain `tabular` with an
% automatically computed width for the last column instead.
% -----------------------------------------------------------------------------
% 3-column version: {Header1}{Header2}{Header3}{width1}{width2}{lastcoltype (ignored)}
\newenvironment{argtable}[6]{%
  \begin{center}
  \renewcommand{\arraystretch}{1.25}
  \setlength{\dimen2}{\dimexpr\textwidth-#4-#5-6\tabcolsep-2\arrayrulewidth\relax}
  \begin{tabular}{@{}p{#4}p{#5}p{\dimen2}@{}}
  \toprule
  \textbf{#1} & \textbf{#2} & \textbf{#3} \\
  \midrule
}{%
  \bottomrule
  \end{tabular}
  \end{center}
}

% 2-column version: {Header1}{Header2}{width1}{lastcoltype (ignored)}
\newenvironment{argtabletwo}[4]{%
  \begin{center}
  \renewcommand{\arraystretch}{1.25}
  \setlength{\dimen2}{\dimexpr\textwidth-#3-4\tabcolsep-2\arrayrulewidth\relax}
  \begin{tabular}{@{}p{#3}p{\dimen2}@{}}
  \toprule
  \textbf{#1} & \textbf{#2} \\
  \midrule
}{%
  \bottomrule
  \end{tabular}
  \end{center}
}

% -----------------------------------------------------------------------------
% Note / remark box
% -----------------------------------------------------------------------------
\newtcolorbox{tbnote}{
  enhanced, breakable,
  colback=notebg, colframe=notebar,
  boxrule=0pt, leftrule=2.5pt,
  arc=0pt, outer arc=0pt,
  left=8pt, right=8pt, top=6pt, bottom=6pt,
  before skip=8pt, after skip=8pt
}

\hypersetup{
  colorlinks=true,
  linkcolor=notebar,
  urlcolor=notebar,
  citecolor=notebar
}

% =============================================================================
\title{\textbf{G09/G16 MATLAB Toolbox}\\[4pt]
       \large A parsing and visualisation toolbox for Gaussian~09/16 output files}
\author{Sebastiano Trusso \\ CNR -- Istituto per i Processi Chimico-Fisici (IPCF), Messina, Italy}
\date{Version 1.0 --- 2026}

% =============================================================================
\begin{document}

\maketitle
\tableofcontents
\newpage

% =============================================================================
\section{Introduction}
% =============================================================================
This manual documents the \textbf{G09/G16 Toolbox}, a set of MATLAB functions
for parsing, analysing, and visualising Gaussian~09 and Gaussian~16 output
files (\code{.out}/\code{.log}) and formatted checkpoint files
(\code{.fchk}). The toolbox covers:

\begin{itemize}[nosep]
  \item extraction of molecular geometry, energies, thermochemistry, dipole
        moment, polarisability, and hyperpolarisability;
  \item vibrational analysis: normal modes, IR/Raman spectra, and interactive
        mode visualisation;
  \item molecular-orbital energy diagrams and TD-DFT excited-state analysis;
  \item 3D molecular structure rendering (CPK ball-and-stick models) with
        atomic charge and dipole overlays;
  \item utility functions for bond-length tables, Gaussian version detection,
        restart-file generation, and toolbox self-documentation;
  \item real-space molecular-orbital, electron-density, and electrostatic-
        potential isosurfaces evaluated directly from \code{.fchk} basis-set
        data (\code{MOSurface}, \S\ref{sec:gutil-mo-surface});
  \item an interactive click-to-select Raman/IR stick-spectrum browser
        (\code{RamanBrowser}, \S\ref{sec:gutil-raman-browser});
  \item atomic-unit \(\leftrightarrow\) SI conversions for the physical
        quantities computational chemistry results are commonly reported in
        (\code{au2SI}, \S\ref{sec:gutil-au2si}).
  \end{itemize}

Function names follow the \code{G09\_xxx}/\code{G16\_xxx} convention,
indicating which Gaussian version's output format the function targets.
Where the two versions behave identically, this manual documents them in a
single combined subsection (\code{G09\_xxx / G16\_xxx}); where behaviour
differs, the differences are called out explicitly.

\begin{tbnote}
Every parsing function that reads a \code{.out}/\code{.log} file accepts an
optional \code{'Lines'} Name-Value parameter: a cell array of pre-read file
lines. This lets \code{G09\_read\_all}/\code{G16\_read\_all} read the file
from disk once and reuse the same lines across every sub-parser, instead of
each function re-reading and re-splitting the file independently. Passing
\code{'Lines'} is entirely optional --- omitting it falls back to reading
the file normally.
\end{tbnote}

\begin{tbnote}
\textbf{Wrong-toolbox warning.} Because Gaussian~09 and Gaussian~16 output
formats differ, using the wrong toolbox on a file (e.g.\ \code{G16\_energy}
on a Gaussian~09 file) can silently misparse it rather than raising an
obvious error --- for example, \code{G16\_dipole\_polar} on a Gaussian~09
file returns \code{NaN} for the polarisability instead of an error. To
catch this, every function that reads a file from disk checks the
Gaussian-version citation line (\code{"Gaussian NN, Revision X.YY,"}) it
just read and prints a non-blocking \code{warning(...)} if \code{NN} does
not match the toolbox being used, suggesting the other one instead. The
check runs only once per actual disk read (see
Section~\ref{ssub:checkmatch}), so \code{G09\_read\_all}/\code{G16\_read\_all}
produce a single warning, not one per sub-parser.
\end{tbnote}

\newpage

% =============================================================================
\section{Which \code{.in} keywords generate which \code{.out} section}
\label{sec:keyword-section-map}
% =============================================================================
Each extraction function reads one specific section of the \code{.out}/
\code{.log} file, identified by a fixed header string or line pattern.
Most of those sections only appear if the corresponding Gaussian route
keyword was present in the \code{.in}/\code{.gjf} job that produced the
file --- otherwise the function raises an error (or, for optional
response-property fields, returns \code{NaN}/empty rather than failing).
The table below maps each function to the \code{.out} section it reads
and the \code{.in} keyword(s) that generate it. Function names omit the
\code{G09\_}/\code{G16\_} prefix (identical for both versions unless
noted).

\begin{longtable}{@{}p{2.3cm}p{4.5cm}p{3.3cm}p{\dimexpr\textwidth-2.3cm-4.5cm-3.3cm-8\tabcolsep-2\arrayrulewidth\relax}@{}}
\toprule
\textbf{Function} & \textbf{\code{.out} section read} & \textbf{\code{.in} keyword} & \textbf{Notes} \\
\midrule
\endhead
\code{structure} & \code{"Standard orientation"} / \code{"Input orientation"} & (default) & Always printed. \code{nosymm}/\code{nosym} suppresses "Standard orientation", leaving only "Input orientation" \\
\code{energy} (SCF) & \code{"SCF Done:"} & (default) & Printed by any HF/DFT energy evaluation (single point, opt, freq, ...) \\
\code{energy} (thermochemistry) & \code{"Zero-point correction="}, \code{"Sum of electronic and ... Energies="} & \code{freq} & ZPE, H, G, S \\
\code{charges} (Mulliken) & \code{"Mulliken charges:"} & (default) & \\
\code{charges} (APT) & \code{"APT charges:"} & \code{freq} & Printed automatically with any frequency job (dipole-derivative charges) --- no separate \code{pop=apt} needed \\
\code{dipole\_polar} (dipole) & \code{"Dipole moment (...)"} & (default) & \\
\code{dipole\_polar} (static Alpha) & \code{"Exact polarizability:"} & \code{polar} & \\
\code{dipole\_polar} (dynamic Alpha) & \code{"Alpha(-w,w) frequency N"} & \code{polar cphf=rdfreq} & Requires an explicit frequency list (a.u.) after the route/title/charge-multiplicity block in the \code{.in} file \\
\code{nmodes} & \code{"and normal coordinates:"} & \code{freq} & \\
\code{spectra} (IR) & \code{"Harmonic frequencies ... IR intensities"} & \code{freq} & \\
\code{spectra} (Raman) & \code{"Harmonic frequencies ... Raman scattering"} & \code{freq=raman} & \\
\code{orbital\_\discretionary{}{}{}energies} & \code{"Alpha occ./virt. eigenvalues"} (+ \code{Beta ...} for unrestricted) & (default) & Full eigenvalue list is printed unconditionally, not gated by \code{pop=full} \\
\code{convergence} & \code{"Item               Value     Threshold  Converged?"} & \code{opt} & \\
\code{charge\_mult} & charge/multiplicity echo near the top of the file & (default) & Mirrors the \code{.in} charge/multiplicity line \\
\code{route} & route-section echo near the top of the file & (default) & \\
\code{get\_\discretionary{}{}{}bond\_\discretionary{}{}{}length} & derived from the structure section & (default) & Purely geometric (covalent radii); no dedicated Gaussian section \\
\code{nbo\_bonds} & \code{"Natural Bond Orbitals (Summary)"} & \code{pop=nbo} & See Section~\ref{sec:nbo-bonds-ref} \\
\code{gaussian\_\discretionary{}{}{}version} & \code{"Gaussian NN, Revision X.YY,"} citation line & (default) & \\
\code{hyperpolar} (vibrational Beta, G16 only) & \code{"Diagonal vibrational hyperpolarizability:"} & \code{polar freq} & Present even without \code{cphf=rdfreq} \\
\code{hyperpolar} (electronic Beta, G16 only) & \code{"First dipole hyperpolarizability, Beta ..."}, \code{"Beta(0;0,0)"}, \code{"Beta(-w;w,0)"} & \code{polar field=... cphf=rdfreq} & Needs the field/frequency list in the \code{.in} file; \code{polar} alone gives only the vibrational Beta above \\
\code{tddft} (G16 only) & \code{"Excited State N: ..."} & \code{td} (e.g.\ \code{td=(nstates=N)}) & \\
\code{restart} & reuses the structure/route/charge\_mult sections above & (any completed step) & \\
\bottomrule
\end{longtable}

\begin{tbnote}
\textbf{\code{.fchk} functions.} \code{fchk\_read} and
\code{charges\_fchk} (provided in both \code{G09/} and \code{G16/}) do
not read \code{.out}/\code{.log} sections at all --- they read a
formatted checkpoint file, produced by running \code{formchk} on the
binary \code{.chk} file written by the job's \code{\%chk=...} link0
command. No specific route keyword is required beyond \code{\%chk=...}
itself.
\end{tbnote}

\begin{tbnote}
\textbf{Evidence, not assumption.} The keyword requirements above were
verified directly against real Gaussian output files rather than
inferred from general documentation --- e.g.\ APT charges and the full
orbital-eigenvalue list turned out to be printed unconditionally
(no \code{pop=apt}/\code{pop=full} needed), and the vibrational vs.\
electronic hyperpolarisability distinction was confirmed by comparing a
file run with plain \code{polar} against one run with
\code{field=... cphf=rdfreq}, only the latter containing the full
electronic Beta tensor.
\end{tbnote}

\newpage

% =============================================================================
\section{File I/O utilities}
% =============================================================================
\subsection{G09\_read\_lines}
\label{ssub:readlines}

\begin{lstlisting}
lines = G09_read_lines(filename)
\end{lstlisting}

Low-level reader used internally by every other \code{G09\_XXX} parser.
Reads a Gaussian~09 output file and returns a cell array of lines, with line
endings stripped. Transparently handles CRLF line endings (as produced by
Windows G09W builds) and ISO-8859-1 (Latin-1) encoding.

\begin{tbnote}
G16 parsers read files directly with the platform default encoding rather
than calling a separate \code{G16\_read\_lines} helper; the \code{'Lines'}
parameter interface is identical across both toolboxes.
\end{tbnote}

% -----------------------------------------------------------------------------
\subsection{G09\_check\_gaussian\_match / G16\_check\_gaussian\_match}
\label{ssub:checkmatch}

\begin{lstlisting}
ok = G09_check_gaussian_match(lines, filename)
ok = G16_check_gaussian_match(lines, filename)
\end{lstlisting}

Internal helper, called automatically by every function that reads a
\code{.out}/\code{.log} file from disk (not part of the typical user-facing
API, but documented here for completeness). Scans the already-read
\code{lines} (no extra disk I/O) for the \code{"Gaussian NN, Revision
X.YY,"} citation line printed near the top of every output file. If found
and \code{NN} does not match the toolbox (9 for \code{G09\_*}, 16 for
\code{G16\_*}), prints a non-blocking warning suggesting the other
toolbox; otherwise does nothing. Returns \code{true} if no mismatch was
detected (or the version is unknown), \code{false} if a warning was
issued.

\begin{tbnote}
Because this check lives inside each function's own "read the file fresh
from disk" branch, it only fires when a file is actually opened --- if a
function receives pre-read \code{lines} via the \code{'Lines'} parameter
(e.g.\ every sub-parser called by \code{G09\_read\_all}/\code{G16\_read\_all}),
it skips both the disk read and the check, so a single call to
\code{read\_all} produces exactly one warning, not one per sub-parser.
\end{tbnote}

\newpage

% =============================================================================
\section{Core data extraction}
% =============================================================================
\subsection{G09\_route / G16\_route}

\begin{lstlisting}
route = G09_route(filename)
route = G16_route(filename)
route = G16_route(filename, 'Lines', lines)
\end{lstlisting}

Extracts the complete route section (the \code{\#} line(s) between the two
\code{-{}-{}-{}-} separators) as a single string.

\begin{argtable}{Option}{Default}{Description}{2.6cm}{2.4cm}{X}
\code{'Lines'} & \code{\{\}} & Pre-read file lines (Section~\ref{ssub:readlines}); read normally if omitted \\
\end{argtable}

\begin{lstlisting}
r = G09_route('indaco.log');
% '# opt freq=raman b3lyp/6-311++g(d,p) nosymm geom=connectivity field=x+50'
\end{lstlisting}

\begin{tbnote}
Both \code{G09\_route}/\code{G16\_route} and \code{G09\_charge\_mult}/
\code{G16\_charge\_mult} now accept the \code{'Lines'} option, consistent
with the rest of the toolbox.
\end{tbnote}

% -----------------------------------------------------------------------------
\subsection{G09\_structure / G16\_structure}

\begin{lstlisting}
mol = G09_structure(filename)
mol = G09_structure(filename, 'step', 'last')
mol = G09_structure(filename, 'step', 'first')
mol = G09_structure(filename, 'step', N)

mol = G16_structure(filename)
mol = G16_structure(filename, 'orientation', 'standard')
mol = G16_structure(filename, 'orientation', 'input')
\end{lstlisting}

Extracts the molecular geometry from a Gaussian output file.

\begin{tbnote}
\textbf{Difference:} Gaussian~09 writes only \code{"Input orientation:"}
blocks (never \code{"Standard orientation:"}), so \code{G09\_structure} has
no \code{'orientation'} option. Gaussian~16 can write either, so
\code{G16\_structure} exposes \code{'orientation'} = \code{'standard'}
(default, falls back to input if absent) \textbar\ \code{'input'} \textbar\
\code{'auto'}.
\end{tbnote}

\begin{argtable}{Option}{Default}{Description}{2.8cm}{2.6cm}{X}
\code{'step'}        & \code{'last'} & Which geometry block to extract: \code{'first'}, \code{'last'}, or an integer index \\
\code{'orientation'} & \code{'auto'} (G16 only) & \code{'standard'} \textbar\ \code{'input'} \textbar\ \code{'auto'} \\
\code{'Lines'}       & \code{\{\}}   & Pre-read file lines; read normally if omitted \\
\end{argtable}

\begin{argtable}{Field}{Type}{Description}{2.8cm}{2.2cm}{X}
\code{.symbols}     & cell   & Atomic symbols \\
\code{.xyz}         & double & Coordinates in \AA\ (X Y Z) \\
\code{.Z}           & int    & Atomic numbers \\
\code{.Natoms}      & int    & Number of atoms \\
\code{.step}        & int    & Index of the extracted step (1-based) \\
\code{.n\_steps}    & int    & Total geometry blocks found in the file \\
\code{.orientation} & char   & \code{'Input orientation'} (always, G09) or \code{'standard'}/\code{'input'} (G16) \\
\code{.filename}    & char   & Source file name \\
\end{argtable}

\begin{lstlisting}
mol = G09_structure('indaco.log');
mol = G16_structure('calc.out', 'orientation', 'input');
\end{lstlisting}

% -----------------------------------------------------------------------------
\subsection{G09\_energy / G16\_energy}

\begin{lstlisting}
en = G09_energy(filename)
en = G09_energy(filename, 'step', 'last')
en = G09_energy(filename, 'step', N)
\end{lstlisting}

Extracts SCF energy and thermochemistry data.

\begin{argtable}{Field}{Type}{Description}{3.0cm}{1.8cm}{X}
\code{.SCF}        & double & SCF energy (Hartree) \\
\code{.method}     & char   & Method string, e.g.\ \code{'RB3LYP'} \\
\code{.ZPE\_corr}  & double & Zero-point correction (Hartree) \\
\code{.U\_corr}    & double & Thermal correction to energy (Hartree) \\
\code{.H\_corr}    & double & Thermal correction to enthalpy (Hartree) \\
\code{.G\_corr}    & double & Thermal correction to Gibbs free energy (Hartree) \\
\code{.E0}         & double & SCF~+~ZPE (Hartree) \\
\code{.U}          & double & SCF~+~U\_corr (Hartree, G16 only) \\
\code{.H}          & double & SCF~+~H\_corr (Hartree) \\
\code{.G}          & double & SCF~+~G\_corr (Hartree) \\
\code{.S\_JmolK}   & double & Entropy $S$ (J\,mol$^{-1}$K$^{-1}$, G09) \\
\code{.ZPE\_kJ}    & double & Zero-point energy (kJ/mol) \\
\code{.T}, \code{.P} & double & Temperature (K), pressure (atm) \\
\code{.has\_thermo}& logical& \code{false} for opt-only jobs (no freq); \code{*\_corr}, \code{E0}, \code{U}, \code{H}, \code{G}, \code{T}, \code{P} are \code{NaN} in that case \\
\code{.G\_kJ}, \code{.H\_kJ} & double & $G$, $H$ in kJ/mol (G09) \\
\code{.filename}   & char   & \\
\end{argtable}

\begin{argtable}{Option}{Default}{Description}{2.6cm}{2.4cm}{X}
\code{'step'}  & \code{'last'} & Which SCF block to extract: \code{'last'}, \code{'first'}, or integer $N$ \\
\code{'Lines'} & \code{\{\}}   & Pre-read file lines \\
\end{argtable}

% -----------------------------------------------------------------------------
\subsection{G09\_charge\_mult / G16\_charge\_mult}

\begin{lstlisting}
[charge, mult] = G09_charge_mult(filename)
[charge, mult] = G16_charge_mult(filename)
[charge, mult] = G16_charge_mult(filename, 'Lines', lines)
\end{lstlisting}

Extracts molecular charge and spin multiplicity from the
\code{"Charge = ... Multiplicity = ..."} line. Format identical across G09
and G16.

\begin{lstlisting}
[q, m] = G09_charge_mult('indaco.log');   % q = 0, m = 1
\end{lstlisting}

% -----------------------------------------------------------------------------
\subsection{G09\_convergence / G16\_convergence}

\begin{lstlisting}
cv = G09_convergence(filename)
cv = G09_convergence(filename, 'plot', true)

cv = G16_convergence(filename)
cv = G16_convergence(filename, 'plot', true)
\end{lstlisting}

Extracts geometry-optimisation convergence criteria (Max/RMS Force,
Max/RMS Displacement) at every optimisation step, together with their
thresholds and whether/when the job converged. Block format identical
across G09 and G16.

\begin{argtable}{Option}{Default}{Description}{2.6cm}{2.4cm}{X}
\code{'plot'} & \code{false} & Plot the four convergence series against their thresholds \\
\end{argtable}

\begin{argtable}{Field}{Type}{Description}{3.0cm}{2.0cm}{X}
\code{.MaxForce}, \code{.RMSForce} & vector & Max / RMS force per step (a.u.) \\
\code{.MaxDisp}, \code{.RMSDisp}   & vector & Max / RMS displacement per step (a.u.) \\
\code{.thr\_MaxForce}, \code{.thr\_RMSForce}, \code{.thr\_MaxDisp}, \code{.thr\_RMSDisp} & double & Corresponding convergence thresholds \\
\code{.converged} & logical & Whether all four criteria were satisfied \\
\code{.conv\_step} & int/NaN & Step at which convergence was reached \\
\code{.Nsteps}     & int     & Number of optimisation steps found \\
\code{.filename}   & char    & \\
\end{argtable}

\begin{lstlisting}
cv = G16_convergence('V_E00t.out', 'plot', true);
cv.converged
cv.conv_step
\end{lstlisting}

% -----------------------------------------------------------------------------
\subsection{G09\_dipole\_polar / G16\_dipole\_polar}

\begin{lstlisting}
dp = G09_dipole_polar(filename)
dp = G09_dipole_polar(filename, 'units', 'Debye')

dp = G16_dipole_polar(filename)
dp = G16_dipole_polar(filename, 'units', 'SI')
\end{lstlisting}

Extracts the dipole moment and polarisability tensor.

\begin{tbnote}
\textbf{Differences between G09 and G16:} G09 prints \code{"Exact
polarizability:"} as six values on a single line (upper triangle, row
order), rather than G16's formatted \code{Alpha(0;0)}/\code{Alpha(-w;w)}
blocks. G09 has no dynamic (laser-frequency-dependent) polarisability
unless requested via \code{Polar=OptRot} or similar keywords, and
additionally prints a \code{"Diagonal vibrational polarizability:"} line
(3 values, returned in \code{.alpha\_vib}). An approximate CPHF tensor, if
printed by G09 (\code{"Approx polarizability:"}), is returned in
\code{.alpha\_approx}. G16 instead reports full dynamic polarisability
entries per laser frequency (\code{.alpha\_dyn}) and modal derivatives
\code{.dmu\_dQ}/\code{.dalpha\_dQ} (related to IR intensities and Raman
activities; for exact scaled values use \code{G16\_spectra} instead).
\end{tbnote}

\begin{argtable}{Option}{Default}{Description}{2.6cm}{2.4cm}{X}
\code{'units'} & \code{'au'} & \code{'au'} \textbar\ \code{'Debye'} (dipole only, G09) / \code{'Debye'} (full, G16) \textbar\ \code{'SI'} \\
\code{'Lines'} & \code{\{\}} & Pre-read file lines \\
\end{argtable}

\begin{argtabletwo}{Field}{Description}{3.2cm}{X}
\code{.mu\_x}, \code{.mu\_y}, \code{.mu\_z}, \code{.mu\_tot}, \code{.mu\_units} & Dipole components, magnitude, unit string \\
\code{.alpha\_iso}, \code{.alpha\_aniso} & Isotropic polarisability, anisotropy \\
\code{.alpha\_tensor}, \code{.alpha\_units} & $3\times3$ static tensor (a.u.), unit string \\
\code{.alpha\_vib}, \code{.has\_alpha\_vib} & Diagonal vibrational polarisability (a.u.), logical (G09) \\
\code{.alpha\_approx} & CPHF approximate tensor (a.u.), \code{NaN} if absent (G09) \\
\code{.alpha\_dyn}, \code{.N\_dyn} & Struct array of dynamic polarisability entries (per laser frequency), count (G16) \\
\code{.dmu\_dQ}, \code{.dalpha\_dQ}, \code{.has\_deriv} & Dipole/polarisability derivatives w.r.t.\ normal modes, logical (G16) \\
\code{.filename} & \\
\end{argtabletwo}

% -----------------------------------------------------------------------------
\subsection{G09\_nmodes / G16\_nmodes}
\label{sec:nmodes-ref}

\begin{lstlisting}
nm = G09_nmodes(filename)
nm = G09_nmodes(filename, 'modes', [5 10 20])

nm = G16_nmodes(filename)
nm = G16_nmodes(filename, 'section', 'last')
\end{lstlisting}

Extracts normal-mode displacement vectors.

\begin{tbnote}
G09 writes only one \emph{Harmonic frequencies} section per file (even for
opt+freq jobs), so \code{G09\_nmodes} has no \code{'section'} option --- it
is implicit. G16 may contain more than one such section (e.g.\ a
preliminary and a final one) and defaults to the last.
\end{tbnote}

\begin{argtable}{Option}{Default}{Description}{2.6cm}{2.4cm}{X}
\code{'modes'}  & \code{[]}   & Restrict output to specific mode indices \\
\code{'section'}& \code{'last'} (G16 only) & \code{'last'} \textbar\ \code{'first'} \\
\code{'Lines'}  & \code{\{\}} & Pre-read file lines \\
\end{argtable}

\begin{argtabletwo}{Field}{Description}{2.8cm}{X}
\code{.freq}, \code{.IR}, \code{.Raman} & Frequencies (cm$^{-1}$), IR intensities (KM/Mole), Raman activities (\AA$^4$/AMU, \code{[]} if absent) \\
\code{.redmass}, \code{.frcconst} & Reduced masses (AMU), force constants (mDyne/\AA) \\
\code{.symmetry} & Symmetry labels \\
\code{.disp} & $[N_\text{atoms}\times3\times N_\text{modes}]$ Cartesian displacement vectors \\
\code{.Nmodes}, \code{.Natoms}, \code{.has\_Raman}, \code{.filename} & \\
\code{.RamanFr} & $[N_\text{modes}\times K]$ per-incident-light-frequency Raman activities (\AA$^4$/AMU), \code{[]} if the file has no \code{CPHF=RdFreq} data. Column 1 is always the static (0~cm$^{-1}$) value --- numerically identical to \code{.Raman} --- columns 2..$K$ are the dynamic/pre-resonance values, one per incident wavelength actually requested \\
\code{.IncidentLight} & $[1\times K]$ incident light frequencies in cm$^{-1}$, as printed on Gaussian's own \emph{``Incident light (cm**-1): ...''} line (\code{[]} if absent); \code{IncidentLight(1)} is always \code{0.00} \\
\code{.has\_RamanFr} & Logical, true if \code{.RamanFr} is non-empty \\
\end{argtabletwo}

\begin{tbnote}
\textbf{Static vs.\ dynamic (CPHF=RdFreq) Raman activities.} With
\code{Freq=Raman} plus \code{CPHF=RdFreq} (e.g.\ a \code{"785 nm"} line in
the input stream), Gaussian's own output prints, for every mode, both a
generic \code{Raman Activ --} row (numerically the STATIC, 0~cm$^{-1}$
value) and explicit \code{RamAct Fr=~1--}/\code{RamAct Fr=~2--} rows ---
Fr=1 repeating the same static value, Fr=2 being the genuinely
frequency-dependent one at the requested wavelength. Earlier versions of
\code{G09\_nmodes}/\code{G16\_nmodes} read only the generic row (per their
own inline comment, \emph{``RamAct Fr=... lines are skipped''}), so
\code{.Raman} was always the static value even for a \code{CPHF=RdFreq}
job, with the dynamic value silently unavailable. Verified directly on a
785~nm job (\code{MPBA\_SH.out}): mode~1 has \code{.Raman} $=1.7384$
(static) but \code{.RamanFr(1,2)} $=1.8922$ (785~nm, from the same file)
--- a real, non-negligible difference. \code{.IR} has no such ambiguity:
Gaussian prints only one \code{IR Inten --} row (IR intensity comes from
the dipole derivative, not the frequency-dependent polarizability, so it
is not itself computed at more than one incident frequency).
\end{tbnote}

% -----------------------------------------------------------------------------
\subsection{G09\_spectra / G16\_spectra}
\label{sec:spectra-ref}

\begin{lstlisting}
sp = G09_spectra(filename)
sp = G09_spectra(filename, Name, Value, ...)

sp = G16_spectra(filename)
sp = G16_spectra(filename, 'section', 'last')
\end{lstlisting}

Extracts IR/Raman line spectra and generates Lorentzian-broadened
continuous spectra.

\begin{tbnote}
G09 writes only one Harmonic-frequencies section per file, so
\code{G09\_spectra} needs no \code{'section'} parameter; \code{G16\_spectra}
exposes it to select which section to use when more than one is present.
\end{tbnote}

\begin{argtable}{Option}{Default}{Description}{2.6cm}{2.4cm}{X}
\code{'FWHM'}      & \code{10}   & Lorentzian FWHM (cm$^{-1}$) \\
\code{'xmin'}      & \code{0}    & Lower wavenumber limit (cm$^{-1}$) \\
\code{'xmax'}      & \code{4000} & Upper wavenumber limit (cm$^{-1}$) \\
\code{'dx'}        & \code{1}    & Grid step (cm$^{-1}$) \\
\code{'normalize'} & \code{false}& Normalise continua to maximum $=1$ \\
\code{'plot'}      & \code{false}& Generate figure \\
\code{'section'}   & \code{'last'} (G16 only) & \code{'last'} \textbar\ \code{'first'} \\
\code{'Lines'}     & \code{\{\}} & Pre-read file lines \\
\end{argtable}

\begin{argtabletwo}{Field}{Description}{2.8cm}{X}
\code{.freq}, \code{.IR}, \code{.Raman} & Frequencies (cm$^{-1}$), IR intensities, Raman activities (\code{[]} if absent) \\
\code{.has\_Raman}, \code{.Nmodes} & Logical, number of modes \\
\code{.x}, \code{.IR\_cont}, \code{.Raman\_cont} & Wavenumber grid, broadened IR/Raman continua \\
\code{.FWHM}, \code{.filename} & \\
\end{argtabletwo}

\begin{lstlisting}
sp = G16_spectra('calc.out', 'FWHM', 15, 'xmin', 400, 'plot', true);
\end{lstlisting}

% -----------------------------------------------------------------------------
\subsection{G\_spectra\_nm \normalfont{(shared, no G09\_/G16\_ prefix)}}
\label{sec:spectra-nm-ref}

\begin{lstlisting}
sp = G_spectra_nm(nm)
sp = G_spectra_nm(nm, Name, Value, ...)
\end{lstlisting}

Builds the same Lorentzian-broadened IR/Raman continuum spectra as
\code{G09\_spectra}/\code{G16\_spectra}, but directly from an
already-parsed \code{nm} (normal modes) struct instead of reading a
\code{.log}/\code{.out} file. Accepts the \code{nm} struct returned by
\code{G09\_nmodes}/\code{G16\_nmodes}, or the \code{.nm} sub-struct
returned by \code{G09\_fchk\_read}/\code{G16\_fchk\_read}
(Section~\ref{sec:fchk-read-ref}) --- all four share the same field
layout (\code{.freq}, \code{.IR}, \code{.Raman}, \code{.has\_Raman},
\code{.Nmodes}, \code{.filename}), so a single function works with any
of them regardless of Gaussian version or data source.

\begin{tbnote}
This fills a gap left by \code{G09\_fchk\_read}/\code{G16\_fchk\_read}:
their output struct's \code{.nm} field has frequencies and IR/Raman
intensities, but --- unlike \code{G09\_spectra}/\code{G16\_spectra} ---
does not itself build the broadened continuum spectrum. \code{G\_spectra\_nm}
does exactly that, directly from \code{data.nm}, with no \code{.out}/\code{.log}
file needed at all.
\end{tbnote}

\begin{argtable}{Option}{Default}{Description}{2.6cm}{2.4cm}{X}
\code{'FWHM'}      & \code{10}   & Lorentzian FWHM (cm$^{-1}$) \\
\code{'xmin'}      & \code{0}    & Lower wavenumber limit (cm$^{-1}$) \\
\code{'xmax'}      & \code{4000} & Upper wavenumber limit (cm$^{-1}$) \\
\code{'dx'}        & \code{1}    & Grid step (cm$^{-1}$) \\
\code{'normalize'} & \code{false}& Normalise continua to maximum $=1$ \\
\code{'plot'}      & \code{false}& Generate figure \\
\end{argtable}

Output fields are identical to \code{G09\_spectra}/\code{G16\_spectra}'s
(Section~\ref{sec:spectra-ref} above): \code{.freq}, \code{.IR},
\code{.Raman}, \code{.has\_Raman}, \code{.Nmodes}, \code{.x},
\code{.IR\_cont}, \code{.Raman\_cont}, \code{.FWHM}, \code{.filename}.

\begin{lstlisting}
data = G16_fchk_read('molecule.fchk');
sp   = G_spectra_nm(data.nm, 'FWHM', 15, 'plot', true);   % from a .fchk, no .out needed

nm = G16_nmodes('molecule.out');
sp = G_spectra_nm(nm, 'xmin', 400, 'normalize', true);    % from an already-read .out
\end{lstlisting}

\newpage

% =============================================================================
\section{Atomic charges}
% =============================================================================
\subsection{G09\_charges / G16\_charges}

\begin{lstlisting}
ch = G09_charges(filename)
ch = G09_charges(filename, 'type', 'APT')
ch = G09_charges(filename, 'mode', 'heavy')
ch = G09_charges(filename, 'plot', true, 'threshold', 0.05)
\end{lstlisting}

Extracts Mulliken or APT atomic charges and optionally renders them on the
3D structure, with an optional dipole-moment arrow overlay.

\begin{tbnote}
\textbf{Robustness:} handles all known Mulliken label variants across
Gaussian revisions/OSes (e.g.\ \code{"Mulliken atomic charges:"} on G09
Rev.\ A.02/Windows and some Linux builds; \code{"Mulliken charges:"} on G09
Rev.\ B.01/C.01/D.01 and on all G16 revisions), by matching any line
containing the charge-type keyword and \code{"charges:"}, excluding
\code{"Sum of"} and H-summed lines. The exact label found is reported in
\code{ch.label}.
\end{tbnote}

\begin{argtable}{Option}{Default}{Description}{2.8cm}{3.0cm}{X}
\code{'type'}            & \code{'Mulliken'} & \code{'Mulliken'} \textbar\ \code{'APT'} \\
\code{'mode'}            & \code{'atom'}      & \code{'atom'} \textbar\ \code{'heavy'} (H charges summed onto heavy atoms) \\
\code{'plot'}            & \code{true}        & Render labels on the 3D structure \\
\code{'AtomScale'}       & \code{0.35}        & CPK sphere scale \\
\code{'BondTol'}         & \code{1.30}        & Bond detection tolerance \\
\code{'BondList'}        & \code{[]}          & Passed straight through to \code{draw\_molecule}: explicit atom-index pairs (optionally with a pre-computed bond order), bypassing \code{BondTol} --- see the note below \\
\code{'FontSize'}        & \code{8}           & Charge-label font size \\
\code{'ColorScale'}      & \code{'RdBu'}      & \code{'RdBu'} (blue=neg/red=pos) \textbar\ \code{'none'} \\
\code{'threshold'}       & \code{0}           & Hide labels with $|q| <$ threshold \\
\code{'ShowDipole'}      & \code{false}       & Overlay the total dipole moment vector \emph{in the 3D plot}. The dipole moment itself (\code{.dipole} and related output fields, below) is always computed regardless of this option --- see the note below \\
\code{'DipoleOrigin'}    & \code{'negcharge'} & Arrow anchor: \code{'negcharge'} (centroid of negative charges) \textbar\ \code{'poscharge'} \textbar\ \code{'centroid'} \textbar\ \code{[x y z]} \\
\code{'DipoleScale'}     & \code{1.0}         & Arrow length, \AA\ per Debye \\
\code{'DipoleColor'}     & \code{[0\ 0.6\ 0]} & Arrow colour \\
\code{'DipoleLineWidth'} & \code{2.5}         & Arrow line width \\
\code{'ShowDipoleLabel'} & \code{true}        & Annotate the arrow with $|\mu|$ \\
\code{'DipoleFontSize'}  & \code{11}          & Dipole label font size \\
\code{'DipoleUnits'}     & \code{'Debye'}     & Units for $|\mu|$: \code{'Debye'} \textbar\ \code{'au'} \\
\code{'Lines'}           & \code{\{\}}        & Pre-read file lines \\
\end{argtable}

\begin{argtabletwo}{Field}{Description}{2.8cm}{X}
\code{.symbols}, \code{.charges} & Atomic symbols, per-atom charges (e) \\
\code{.charges\_H} & H-summed charges on heavy atoms \\
\code{.sum\_q} & Sum of all charges ($\approx 0$ for neutral molecules) \\
\code{.type}, \code{.label} & \code{'Mulliken'}/\code{'APT'}, exact header label found \\
\code{.Natoms} & \\
\code{.dipole} & \code{[1x3]} dipole moment (Debye) --- always computed (see the note below); \code{[]} if the file has no readable dipole moment \\
\code{.dipole\_origin} & \code{[1x3]} anchor point used for the arrow (\AA), computed alongside \code{.dipole} \\
\code{.dipole\_Debye} & $|\mu|$ in Debye, computed alongside \code{.dipole} \\
\code{.dipole\_au} & $|\mu|$ in atomic units, computed alongside \code{.dipole} \\
\end{argtabletwo}

\begin{tbnote}
\textbf{The dipole moment is always computed.} \code{'ShowDipole'} only
controls whether the arrow is actually drawn in the 3D plot (and, when
true, whether a dipole read failure is reported as a warning) ---
\code{.dipole}/\code{.dipole\_origin}/\code{.dipole\_Debye}/
\code{.dipole\_au} are populated in the returned struct regardless, since
reading the dipole moment is cheap. The underlying
\code{G09/G16\_dipole\_polar} call's own console summary is suppressed
unless \code{'ShowDipole'} is true, so this does not add unrequested
console output to the common (\code{'ShowDipole'} left at its default
\code{false}) case.
\end{tbnote}

\begin{tbnote}
\textbf{Rendering NBO-derived bond orders.} \code{'BondList'} accepts
the same \code{[Nbonds x 3]} \code{Atom1}, \code{Atom2}, \code{BondOrder}
matrix that \code{draw\_molecule} does (Section~\ref{sec:nbo-bonds-ref}),
so the charge overlay renders on top of the real NBO-derived bond order
instead of the geometric estimate:
\begin{lstlisting}
bt = G16_nbo_bonds('molecule_nbo.out');
bond_list = [bt.Atom1, bt.Atom2, double(bt.BondOrder)];
G16_charges('molecule_nbo.out', 'BondList', bond_list, 'ShowDipole', true);
\end{lstlisting}
\end{tbnote}

% -----------------------------------------------------------------------------
\subsection{G09\_charges\_fchk / G16\_charges\_fchk}

\begin{lstlisting}
ch_out = G09_charges_fchk(mol, ch)
ch_out = G09_charges_fchk(mol, ch, Name, Value, ...)

ch_out = G16_charges_fchk(mol, ch)
\end{lstlisting}

Visualises atomic charges directly from a \code{G09\_fchk\_read}/
\code{G16\_fchk\_read} struct, without requiring the corresponding
\code{.log}/\code{.out} file. Mirrors the interface of \code{G09\_charges}/
\code{G16\_charges} exactly, so the two can be used interchangeably once
\code{data.mol} and \code{data.ch} are available.

\begin{argtable}{Option}{Default}{Description}{2.6cm}{2.6cm}{X}
\code{'mode'}       & \code{'atom'}  & \code{'atom'} \textbar\ \code{'heavy'} (requires \code{ch.charges\_H} non-empty --- \code{.fchk} files have no native H-summed data; use \code{G09\_charges}/\code{G16\_charges} on the \code{.log} for that) \\
\code{'plot'}       & \code{true}    & \\
\code{'AtomScale'}  & \code{0.35}    & \\
\code{'BondTol'}    & \code{1.30}    & \\
\code{'FontSize'}   & \code{8}       & \\
\code{'ColorScale'} & \code{'RdBu'}  & \code{'RdBu'} \textbar\ \code{'none'} \\
\code{'threshold'}  & \code{0}       & \\
\end{argtable}

\begin{lstlisting}
data = G16_fchk_read('3typ.fchk');
ch   = G16_charges_fchk(data.mol, data.ch);                      % default
ch   = G16_charges_fchk(data.mol, data.ch, 'plot', false, 'threshold', 0.05);
ch   = G16_charges_fchk(data.mol, data.ch, 'ColorScale', 'none'); % drop-in for G16_charges
\end{lstlisting}

Requires \code{mol} with fields \code{.symbols}, \code{.xyz}, \code{.Natoms},
\code{.filename}, and \code{ch} with fields \code{.charges}, \code{.symbols},
\code{.type}, \code{.Natoms} (\code{.charges\_H} optional).

\begin{tbnote}
Since the \code{.fchk} format is identical for G09 and G16,
\code{G09\_charges\_fchk} and \code{G16\_charges\_fchk} are functionally
identical (the latter simply calls \code{G16\_draw\_molecule} internally
instead of \code{G09\_draw\_molecule}) --- both are provided so that either
toolbox folder alone is enough.
\end{tbnote}

\newpage

% =============================================================================
\section{Formatted checkpoint (.fchk) support}
% =============================================================================
\subsection{G09\_fchk\_read / G16\_fchk\_read}
\label{sec:fchk-read-ref}

\begin{lstlisting}
data = G09_fchk_read(filename)
data = G09_fchk_read(filename, 'verbose', false)

data = G16_fchk_read(filename)
\end{lstlisting}

Reads a Gaussian 09/16 formatted checkpoint file (\code{.fchk}), generated
with \code{formchk jobname.chk jobname.fchk}. Uses a single generic section
parser (no hard-coded line counts) and is compatible with both G09 and G16
\code{.fchk} files.

\begin{tbnote}
Supersedes the old \code{.fck}/newzmat-format reader: reads the modern
\code{.fchk} ASCII format, returns SI-ready fields (XYZ in \AA, dipole in
Debye), and needs no external helper functions. The formatted-checkpoint
layout does not depend on the Gaussian version, so \code{G09\_fchk\_read}
and \code{G16\_fchk\_read} are identical and each works on \code{.fchk}
files produced by \emph{either} G09 or G16 --- both are provided purely so
that either toolbox folder alone (\code{G09/} or \code{G16/}) is enough to
read any \code{.fchk} file, without needing the other one on the path too.
\end{tbnote}

\begin{argtable}{Option}{Default}{Description}{2.6cm}{2.4cm}{X}
\code{'verbose'} & \code{true} & Print progress messages while parsing \\
\end{argtable}

\textbf{Output struct \code{data}} (grouped by section):

\begin{argtabletwo}{Field}{Description}{3.2cm}{X}
\code{.title}, \code{.method}, \code{.basis} & First line, method (e.g.\ \code{'RB3LYP'}), basis set (e.g.\ \code{'def2SVPP'}) \\
\code{.Nat}, \code{.charge}, \code{.mult} & Number of atoms, molecular charge, spin multiplicity \\
\code{.Nelec}, \code{.Nalpha}, \code{.Nbeta} & Total, alpha, beta electron counts \\
\code{.Nbasis}, \code{.Nbasis\_indep} & Basis functions, independent basis functions \\
\code{.symbols}, \code{.AN}, \code{.masses} & Atomic symbols, atomic numbers, real atomic masses (AMU) \\
\code{.xyz}, \code{.xyz\_bohr} & Coordinates (\AA), coordinates (Bohr) \\
\code{.SCF\_energy}, \code{.total\_energy}, \code{.virial\_ratio} & SCF energy, total energy (Hartree), virial ratio ($\approx 2$) \\
\code{.alpha\_\discretionary{}{}{}orb\_\discretionary{}{}{}energies}, \code{.beta\_orb\_energies} & MO energies (Hartree); beta is \code{[]} if RHF \\
\code{.alpha\_MO\_coeff} & Alpha MO coefficients (column-major) \\
\code{.mulliken\_charges} & Mulliken charges (e) \\
\code{.HOMO\_idx}, \code{.HOMO\_eV}, \code{.LUMO\_eV}, \code{.gap\_eV} & Frontier-orbital summary \\
\code{.gradient}, \code{.rms\_force} & Cartesian gradient (Hartree/Bohr), RMS force \\
\end{argtabletwo}

\newpage

% =============================================================================
\section{Molecule and normal-mode visualisation}
% =============================================================================
\subsection{G09\_draw\_molecule / G16\_draw\_molecule}
\label{sec:draw-molecule}

\begin{lstlisting}
G09_draw_molecule(mol)
G09_draw_molecule(mol, 'Name', Value, ...)
\end{lstlisting}

Renders a 3D CPK ball-and-stick model from a \code{mol} struct (as returned
by \code{G09\_structure}/\code{G16\_structure}), with bonds detected
automatically from a covalent-radius criterion and drawn as 1, 2, or 3
parallel lines according to an estimated bond order.

\begin{argtable}{Option}{Default}{Description}{2.6cm}{2.4cm}{X}
\code{'AtomScale'}   & \code{0.35}             & Sphere radius as a fraction of the covalent radius \\
\code{'BondTol'}     & \code{1.30}             & Bond tolerance: bonded if $d < (r_1+r_2)\times$\code{BondTol} \\
\code{'ShowLabels'}  & \code{true}             & Show atom labels (\code{"C1"}, \code{"H2"}, ...) \\
\code{'ShowLegend'}  & \code{true}             & Show the element-colour legend \\
\code{'Title'}       & filename                & Figure title \\
\code{'BgColor'}     & \code{[0.95 0.95 0.95]} & Axes background colour \\
\code{'Ax'}          & (new figure)            & Draw into an existing axes handle \\
\code{'ShowAxes'}    & \code{false}            & Draw a small Cartesian X/Y/Z reference triad, anchored at the lower-left-front corner of the plotted molecule \\
\code{'AxesLength'}  & \code{[]}               & Length of the X/Y/Z arrows, in \AA\ (\code{[]} = auto, 20\% of the molecule's bounding-box diagonal) \\
\code{'BondList'}    & \code{[]}               & $[N_\text{bonds}\times2]$ or $[N_\text{bonds}\times3]$ explicit atom-index pairs (optionally with a 3rd column giving a pre-computed bond order 1/2/3) to draw as bonds, bypassing the \code{BondTol} distance criterion (\code{[]} = auto-detect). Used by \code{G09\_animate\_mode}/\code{G16\_animate\_mode} to keep a fixed bond topology (and, with the 3rd column, a fixed bond order) across all frames of an animation \\
\end{argtable}

\begin{lstlisting}
mol = G09_structure('zeatin.out');
G09_draw_molecule(mol);
G09_draw_molecule(mol, 'ShowLabels', false, 'AtomScale', 0.4);
G09_draw_molecule(mol, 'ShowAxes', true);
\end{lstlisting}

\begin{tbnote}
\textbf{Bond order (single/double/triple).} For C-C and C-O pairs, bond
order is estimated purely from bond length (any other element pair,
including C-N, is always drawn as a single bond) and rendered as 1/2/3
parallel lines in the usual chemical-drawing convention --- a geometric
estimate, not Gaussian's own bond-order analysis (e.g.\ Wiberg/NBO
indices). For an actual bond order derived from Gaussian's own NBO
analysis, see \code{G09\_nbo\_bonds}/\code{G16\_nbo\_bonds}
(Section~\ref{sec:nbo-bonds-ref}). The C-C
double/single threshold is deliberately set to
1.36~\AA\ rather than the generic reference-length midpoint
($\approx$1.44~\AA): symmetric aromatic rings have essentially equal C-C
bond lengths ($\approx$1.39--1.40~\AA, with no real single/double
alternation to recover from geometry alone), so the generic midpoint
would classify every ring bond as double. With the adjusted threshold,
aromatic rings are instead rendered as all-single --- a more honest
depiction than an arbitrary alternating pattern, since the true bond
order is $\approx$1.5 all the way around the ring.
\end{tbnote}

\begin{tbnote}
\textbf{Rendering NBO-derived bond orders.} \code{nbo\_bonds}'s output
table can be fed straight into \code{draw\_molecule}'s \code{'BondList'}
parameter (Section~\ref{sec:nbo-bonds-ref}), replacing the geometric
estimate above with the real bond order from Gaussian's own NBO
analysis:
\begin{lstlisting}
mol = G16_structure('molecule_nbo.out');
bt  = G16_nbo_bonds('molecule_nbo.out');   % requires 'pop=nbo' in the source job

bond_list = [bt.Atom1, bt.Atom2, double(bt.BondOrder)];
G16_draw_molecule(mol, 'BondList', bond_list);
\end{lstlisting}
\end{tbnote}

\begin{tbnote}
\textbf{Why C-N is excluded.} C-N was classified by length in an earlier
version of the toolbox, using the same kind of threshold as C-C/C-O.
Unlike C-C, a single threshold cannot separate aromatic C-N (e.g.\
pyridine rings, $\approx$1.34~\AA) from a genuine C=N: on asymmetric
pyridine-like rings this produced one ring C-N bond rendered double and
its chemically-equivalent neighbour rendered single, an inconsistent and
misleading picture. C-N is therefore always drawn as a single bond,
regardless of length.
\end{tbnote}

\begin{tbnote}
\textbf{Bug fixed (2026-08-05): lighting not following interactive
rotation.} The plot supports interactive rotation via \code{rotate3d}
(drag with the mouse). \code{CAMLIGHT} only positions a light relative
to the camera at the moment it is called --- MATLAB's own documentation
is explicit about this: ``In order for a light created with CAMLIGHT to
stay in a constant position relative to the camera, CAMLIGHT must be
called whenever the camera is moved.'' \code{ROTATE3D} does not do this
on its own, so previously the two lights stayed fixed in the axes'
3-D world frame while the molecule rotated underneath them: after a
large enough rotation (e.g.\ 180 degrees) the CPK spheres ended up lit
from behind instead of from the viewer's side, losing the shading
gradient that gives them their 3-D appearance and making them look
flat/uniformly lit. Fixed by keeping a handle to both lights and
assigning \code{rotate3d}'s \code{ActionPostCallback} property (not an
event --- assigning it again simply replaces the previous callback, so
repeated calls into the same figure, e.g.\ one \code{draw\_molecule}
call per \code{animate\_mode} frame, do not accumulate callbacks) to
reposition both lights via \code{CAMLIGHT} after every interactive
rotation.
\end{tbnote}

No differences between the G09 and G16 versions beyond the input struct's
origin.

% -----------------------------------------------------------------------------
\subsection{G09\_draw\_mode / G16\_draw\_mode}

\begin{lstlisting}
G09_draw_mode(mol, nm, mode_idx)
G09_draw_mode(mol, nm, mode_idx, 'Name', Value, ...)
\end{lstlisting}

Renders a single vibrational normal mode as displacement arrows over the
CPK structure.

\begin{argtable}{Option}{Default}{Description}{2.6cm}{2.4cm}{X}
\code{'Scale'}      & \code{1.5}           & Arrow length scale \\
\code{'ArrowColor'} & \code{[1 0.4 0.1]}   & Arrow colour \\
\code{'AtomScale'}  & \code{0.35}          & CPK sphere scale factor \\
\code{'BondTol'}    & \code{1.30}          & Bond detection tolerance \\
\code{'ShowLabels'} & \code{false}         & Show atom index labels \\
\code{'FlipSign'}   & \code{false}         & Invert all arrow directions \\
\end{argtable}

\begin{tbnote}
\textbf{Why \code{'FlipSign'} exists.} Normal-mode eigenvectors have an
arbitrary overall sign --- a mode and its 180\textdegree-phase-shifted twin
are physically identical. If the arrows point opposite to GaussView's
rendering for a given mode, set \code{'FlipSign', true} to flip them. This
is purely cosmetic and has no effect on frequencies or intensities.
\end{tbnote}

% -----------------------------------------------------------------------------
\subsection{G09\_modeViewer / G16\_modeViewer}

\begin{lstlisting}
G09_modeViewer(filename)
G09_modeViewer(filename, 'Name', Value, ...)
\end{lstlisting}

Interactive selector GUI for browsing vibrational modes. Reads the
structure and all normal modes via \code{G09\_structure}/\code{G09\_nmodes},
then opens a window listing every mode (index, frequency, symmetry).
Selecting an entry calls \code{G09\_draw\_mode} to render that mode in a new
figure window; previously drawn mode figures stay open, so different modes
(or render settings) can be compared side by side.

\begin{tbnote}
\textbf{Window controls:}
\begin{itemize}[nosep]
  \item \emph{Order by:} sort the mode list by mode number (default), IR
        intensity, or Raman intensity (descending; only offered when Raman
        data is present). The current selection is preserved across
        reordering.
  \item A text box to set a custom title on the current mode figure.
  \item A control for every \code{G09\_draw\_mode} option (\code{Scale},
        \code{ArrowColor}, \code{AtomScale}, \code{BondTol},
        \code{ShowLabels}, \code{FlipSign}); changing any of them
        immediately redraws the selected mode in a new figure. Arrow colour
        can be set from a named preset (Orange/Red/Blue/Green/Black/
        Magenta) or a custom colour.
  \item A \emph{Save figure...} button exporting the target mode figure to
        JPEG, EPS, or PDF (EPS/PDF as true vector graphics).
  \item An \emph{Animate mode (MP4)...} button exporting an MP4 animation
        of the currently selected mode via \code{G09\_animate\_mode}/
        \code{G16\_animate\_mode}, using the current Scale/AtomScale/
        BondTol/ShowLabels/FlipSign settings and, automatically, the
        camera orientation of the currently displayed mode figure (a
        manual rotation carries over into the animation).
\end{itemize}
Name-Value arguments passed to \code{G09\_modeViewer} pre-populate the
corresponding option control(s), e.g.\ \code{G09\_modeViewer('file.log',
'Scale', 2, 'ShowLabels', true)}.
\end{tbnote}

\begin{lstlisting}
G09_modeViewer('V_E00t.out');
G16_modeViewer('INDACO_E025_ppDP.LOG', 'Scale', 2, 'ShowLabels', true);
\end{lstlisting}

See also \code{G09\_structure}/\code{G16\_structure},
\code{G09\_nmodes}/\code{G16\_nmodes}, \code{G09\_draw\_mode}/\code{G16\_draw\_mode},
\code{G09\_animate\_mode}/\code{G16\_animate\_mode}.

% -----------------------------------------------------------------------------
\subsection{G09\_animate\_mode / G16\_animate\_mode}

\begin{lstlisting}
outfile = G09_animate_mode(mol, nm, mode_idx)
outfile = G09_animate_mode(mol, nm, mode_idx, 'Name', Value, ...)

outfile = G16_animate_mode(mol, nm, mode_idx)
\end{lstlisting}

Exports an MP4 animation of a single vibrational mode: the molecule
oscillates along the mode's displacement vector (equilibrium
$\pm$~amplitude), in the style of GaussView's mode animations, and the
result is saved via \code{VideoWriter} (\code{'MPEG-4'} profile). A figure
window opens and animates live while the video is rendered; avoid
interacting with it until the function returns.

\begin{argtable}{Option}{Default}{Description}{2.8cm}{2.6cm}{X}
\code{'Filename'}       & \code{[]} (auto) & Output path; default \code{<source>\_mode<N>.mp4}, \code{.mp4} appended if missing \\
\code{'Scale'}          & \code{1.5}    & Displacement amplitude scale, same meaning as \code{G09\_draw\_mode}/\code{G16\_draw\_mode}'s \code{'Scale'} \\
\code{'FlipSign'}       & \code{false}  & Invert the displacement direction \\
\code{'AtomScale'}      & \code{0.35}   & See \code{G09\_draw\_molecule}/\code{G16\_draw\_molecule} \\
\code{'BondTol'}        & \code{1.30}   & See \code{G09\_draw\_molecule}/\code{G16\_draw\_molecule}; also used to compute the fixed bond list (see below) \\
\code{'ShowLabels'}     & \code{false}  & See \code{G09\_draw\_molecule}/\code{G16\_draw\_molecule} \\
\code{'FramesPerCycle'} & \code{30}     & Frames per oscillation period \\
\code{'NCycles'}        & \code{2}      & Number of periods rendered \\
\code{'FPS'}            & \code{20}     & Video frame rate \\
\code{'View'}           & \code{[]}     & \code{[azimuth elevation]} in degrees, the starting camera orientation (\code{[]} = MATLAB's default 3D view, azimuth$=-37.5$, elevation$=30$). Pass the output of MATLAB's own \code{view(ax)} to match an already-rotated figure \\
\code{'ShowWaitbar'}    & \code{true}   & Print rendering progress to the command window \\
\code{'Crop'}           & \code{true}   & Tightly crop every frame around the molecule instead of writing the whole (mostly blank) figure window (see below) \\
\code{'CropMargin'}     & \code{12}     & Blank margin kept around the molecule when \code{'Crop'} is true, in pixels \\
\code{'ShowTitle'}      & \code{false}  & Draw the filename/mode/frequency title text in every frame. Left off by default so the title's own bounding box does not force a taller crop than the molecule itself needs \\
\end{argtable}

\begin{lstlisting}
mol = G16_structure('V_E00t.out');
nm  = G16_nmodes('V_E00t.out');
G16_animate_mode(mol, nm, 70, 'Scale', 2, 'FPS', 25);
\end{lstlisting}

\begin{tbnote}
\textbf{Fixed bond list and bond order.} Bonds (and their estimated order,
1/2/3) are computed once from the \emph{equilibrium} geometry (via
\code{G09\_get\_bond\_length}/\code{G16\_get\_bond\_length}, using
\code{'BondTol'} as the tolerance) and passed to
\code{G09\_draw\_molecule}/\code{G16\_draw\_molecule} through its
\code{'BondList'} parameter (atom-index pairs plus a 3rd bond-order
column), instead of being re-detected from instantaneous inter-atomic
distances on every frame. Without this, bonds would flicker in and out
--- and double/triple bonds would flicker to/from single --- as
oscillating atoms cross the relevant distance thresholds.

\textbf{No graphical waitbar.} Rendering progress is printed to the
command window rather than shown with MATLAB's \code{waitbar}: on some
MATLAB versions, \code{waitbar} itself was found to error intermittently
across many rapid updates (\code{"Not enough input arguments"} inside its
own internal \code{title()} call).

\textbf{Axes handling.} Each frame's axes are deleted and recreated
(rather than cleared with \code{cla}) before redrawing: in this
3D-plus-lighting configuration, \code{cla(ax)} was found to leave stale
graphics objects behind, causing the exported video to show all frames
superimposed instead of one at a time.

\textbf{Frame cropping.} Without \code{'Crop'}, every frame is the full
figure window (default MATLAB figure size), leaving most of it blank
around the molecule. The crop rectangle is computed \emph{once}, from the
two oscillation extremes (phase $+1$ and $-1$, which bound the whole
animation), by finding the tight bounding box of non-background pixels
and padding it by \code{'CropMargin'}; the same rectangle is then applied
to every frame. Because the axis limits are already fixed across the
whole animation (see above), this rectangle stays valid throughout the
video without ever clipping an atom at any phase, typically reducing the
frame area by 70--80\%.
\end{tbnote}

No differences between the G09 and G16 versions beyond the input structs'
origin.

See also \code{G09\_draw\_molecule}\slash\code{G16\_draw\_molecule},
\code{G09\_draw\_mode}\slash\code{G16\_draw\_mode},
\code{G09\_modeViewer}\slash\code{G16\_modeViewer},
\code{G09\_get\_bond\_length}\slash\code{G16\_get\_bond\_length}.

% -----------------------------------------------------------------------------
\subsection{G09\_draw\_deformation / G16\_draw\_deformation}

\begin{lstlisting}
G09_draw_deformation(mol1, mol2)
G09_draw_deformation(mol1, mol2, 'Name', Value, ...)
\end{lstlisting}

Visualises the geometric deformation between two related structures of
the \emph{same} molecule --- typically with vs without a small static
electric field (finite-field NLO workflows), two conformers, or
before/after optimisation --- by drawing an arrow from each atom's
position in \code{mol1} to its position in \code{mol2}
(\code{mol2.xyz - mol1.xyz}), the same displacement-arrow style
\code{G09\_draw\_mode} uses for a normal mode, but here the actual
geometric shift between two structures rather than a vibrational
eigenvector. \code{mol1}/\code{mol2} can come from
\code{G09\_structure}/\code{G16\_structure}, \code{G\_read\_structure\_file},
or \code{G09\_fchk\_read}/\code{G16\_fchk\_read}'s \code{.mol} sub-struct
--- same atom count and ordering required.

\begin{argtable}{Option}{Default}{Description}{2.8cm}{2.8cm}{X}
\code{'Overlay'}        & \code{false}         & See "Two display modes" below \\
\code{'Scale'}          & \code{1}             & Arrow length AND overlay-position multiplier, applied identically to both (see below) \\
\code{'ArrowColor'}     & \code{[1 0.4 0.1]}   & Arrow colour (same default as \code{G09\_draw\_mode}) \\
\code{'Overlay2Color'}  & \code{[0.75 0.1 0.1]}& Colour of \code{mol2}'s skeletal overlay when \code{'Overlay'} is true \\
\code{'AtomScale'}      & \code{0.35}          & CPK sphere scale for \code{mol1} \\
\code{'BondTol'}        & \code{1.30}          & Bond-detection tolerance, both structures \\
\code{'ShowLabels'}     & \code{false}         & Show atom index labels on \code{mol1} \\
\code{'ArrowThreshold'} & \code{0.05}          & Skip the arrow for atoms whose displacement is below this fraction of the single largest per-atom displacement \\
\end{argtable}

\subsubsection*{Two display modes}

\begin{itemize}[nosep]
  \item \textbf{\code{'Overlay', false}} (default): draws \code{mol1}
        only, as a full CPK ball-and-stick model, with displacement
        arrows.
  \item \textbf{\code{'Overlay', true}}: ALSO draws \code{mol2} as a
        simplified skeletal overlay (thin dashed lines + small dots, no
        spheres) in \code{'Overlay2Color'}, so both ends of every arrow
        are visually anchored to a real atom position instead of just an
        arrowhead in empty space.
\end{itemize}

\begin{tbnote}
\textbf{Console report.} Every call prints the largest single-atom
displacement (with which atom) and the RMSD over all atoms, in
\AA{} --- the numbers to look at before picking \code{'Scale'}: real
deformations (e.g.\ field-induced) are often only $10^{-3}$--$10^{-2}$
\AA, far too small to see at \code{Scale=1}.

\textbf{Why \code{'Scale'} affects the overlay too.} The overlay is drawn
at \code{mol1.xyz + Scale*(mol2.xyz - mol1.xyz)}, i.e.\ the SAME
exaggerated position the arrows point to --- \emph{not} \code{mol2}'s true
coordinates. At \code{Scale} $\neq 1$ (the normal case for these
typically-tiny deformations), drawing the overlay at the true,
unscaled position would leave it visually detached from the arrow tips,
since the arrows themselves are drawn exaggerated. Bond connectivity for
the overlay is still detected from \code{mol2}'s true, unscaled geometry
(an exaggerated geometry at a large \code{Scale} is not a real molecular
geometry and would misjudge bonding distances).
\end{tbnote}

\begin{lstlisting}
mol1 = G16_structure('nofield.out');
mol2 = G16_structure('field_x025.out');
G16_draw_deformation(mol1, mol2, 'Scale', 200);
G16_draw_deformation(mol1, mol2, 'Overlay', true, 'Scale', 200);
\end{lstlisting}

No differences between the G09 and G16 versions beyond which
\code{draw\_molecule}/\code{get\_bond\_length} pair each calls internally.

See also \code{G09\_draw\_molecule}\slash\code{G16\_draw\_molecule},
\code{G09\_draw\_mode}\slash\code{G16\_draw\_mode},
\code{G09\_get\_bond\_length}\slash\code{G16\_get\_bond\_length}.

\newpage

% =============================================================================
\section{au2SI: atomic units \texorpdfstring{$\leftrightarrow$}{<->} SI conversions}
\label{sec:gutil-au2si}
% =============================================================================
\code{AU\_convert}/\code{AU\_table} convert physical quantities between
atomic units (au) and SI (or other common units used in computational
chemistry); \code{G\_gaussian\_field\_convert} additionally converts
Gaussian's own \code{Field=X+N} route-keyword integer to/from a physical
field strength. All conversion factors are based on CODATA 2022 (NIST,
November 2024 release). The toolbox covers 16 physical-quantity
categories, all three functions are also ported to Python directly
inside \code{G16parser} (\S\ref{sec:gutil-au2si-python}):

\begin{argtabletwo}{Quantity}{au unit}{4.2cm}{X}
Energy               & Hartree (\(E_h\)) \\
Length               & Bohr (\(a_0\)) \\
Dipole moment        & \(e\cdot a_0\) \\
Polarisability \(\alpha\)        & \(e^2a_0^2/E_h\) \\
1st hyperpolarisability \(\beta\)& \(e^3a_0^3/E_h^2\) \\
2nd hyperpolarisability \(\gamma\)& \(e^4a_0^4/E_h^3\) \\
Force                & \(E_h/a_0\) \\
Force constant       & \(E_h/a_0^2\) \\
Frequency            & \(E_h/\hbar\) \\
Time                 & \(\hbar/E_h\) \\
Pressure             & \(E_h/a_0^3\) \\
Electric field       & \(E_h/(e\,a_0)\) \\
Mass                 & \(m_e\) \\
Charge               & \(e\) \\
Magnetic moment      & \(\mu_B\) \\
Temperature          & \(k_BT \leftrightarrow K\) \\
\end{argtabletwo}

\begin{tbnote}
No additional toolboxes or external dependencies are required. Both
functions use only MATLAB built-ins.
\end{tbnote}

% -----------------------------------------------------------------------------
\subsection{AU\_convert}
\label{sec:gutil-au-convert}

\begin{lstlisting}
result = AU_convert(value, quantity, direction)
result = AU_convert(value, quantity, direction, 'target', unit)
AU_convert([], 'help', '')               % print full reference table
\end{lstlisting}

\begin{argtabletwo}{Argument}{Description}{2.4cm}{X}
\code{value}     & Scalar or array of values to convert (double). Arrays are converted element-wise. \\
\code{quantity}  & Physical quantity string (case-insensitive) --- see the reference table below. \\
\code{direction} & \code{'au2si'} --- from atomic units to SI (or target unit); \code{'si2au'} --- from SI (or target unit) to atomic units \\
\code{'target'}  & Name-Value: target unit when multiple options exist for the SI side. Optional --- a default is used when omitted. \\
\end{argtabletwo}

\code{AU\_convert} returns a \code{double} array of the same shape as
\code{value}. When \code{value} is a scalar, the conversion is also
printed to the console.

\begin{tbnote}
\textbf{Quantity strings and available targets.} The table below lists
every supported quantity string, the atomic unit on the left side, and
the available target units on the SI side. The default target (used when
\code{'target'} is omitted) is the first row of each quantity block.
\end{tbnote}

\begin{longtable}{@{}p{2.4cm}p{2.6cm}p{2.4cm}p{\dimexpr\textwidth-2.4cm-2.6cm-2.4cm-8\tabcolsep-2\arrayrulewidth\relax}@{}}
\toprule
\textbf{quantity string} & \textbf{au unit} & \textbf{'target'} & \textbf{SI/target unit (1 au =)} \\
\midrule
\endhead
\code{'energy'}     & Hartree ($E_h$) & (default) & J \quad $4.3597447\times10^{-18}$ \\
                     &                 & \code{'kJ/mol'}   & $2625.4996$ kJ/mol \\
                     &                 & \code{'kcal/mol'} & $627.5095$ kcal/mol \\
                     &                 & \code{'eV'}       & $27.211386$ eV \\
                     &                 & \code{'cm-1'}     & $219474.631$ cm$^{-1}$ \\
\code{'length'}     & Bohr ($a_0$)    & (default) & m \quad $5.2917721\times10^{-11}$ \\
                     &                 & \code{'Angstrom'} & $0.5291772$ \AA \\
                     &                 & \code{'pm'}       & $52.9177$ pm \\
\code{'dipole'}     & $e\cdot a_0$    & (default) & C$\cdot$m \quad $8.4783536\times10^{-30}$ \\
                     &                 & \code{'Debye'}    & $2.541747$ D \\
\code{'polar'}      & $e^2a_0^2/E_h$  & (default) & C$^2\cdot$m$^2\cdot$J$^{-1}$ \quad $1.6487773\times10^{-41}$ \\
                     &                 & \code{'Angstrom3'}& $0.148185$ \AA$^3$ (volume) \\
                     &                 & \code{'esu'}      & $1.48185\times10^{-25}$ cm$^3$ (esu) \\
\code{'beta'}       & $e^3a_0^3/E_h^2$& (default) & C$^3\cdot$m$^3\cdot$J$^{-2}$ \quad $3.2064\times10^{-53}$ \\
                     &                 & \code{'esu'}      & $8.6392\times10^{-33}$ esu \\
\code{'gamma'}      & $e^4a_0^4/E_h^3$& (default) & C$^4\cdot$m$^4\cdot$J$^{-3}$ \quad $6.2354\times10^{-65}$ \\
\code{'force'}      & $E_h/a_0$       & (default) & N \quad $8.2387\times10^{-8}$ \\
                     &                 & \code{'nN'}       & $82.387$ nN \\
\code{'frcconst'}   & $E_h/a_0^2$     & (default) & N/m \quad $1556.893$ \\
                     &                 & \code{'mDyne/A'}  & $15.5689$ mDyne/\AA \\
\code{'frequency'}  & $E_h/\hbar$     & (default) & rad/s \quad $4.13414\times10^{16}$ \\
                     &                 & \code{'Hz'}       & $6.57968\times10^{15}$ Hz \\
                     &                 & \code{'THz'}      & $6579.684$ THz \\
                     &                 & \code{'cm-1'}     & $219474.631$ cm$^{-1}$ \\
\code{'time'}       & $\hbar/E_h$     & (default) & s \quad $2.41888\times10^{-17}$ \\
                     &                 & \code{'fs'}       & $0.024189$ fs \\
                     &                 & \code{'as'}       & $24.1888$ as \\
\code{'pressure'}   & $E_h/a_0^3$     & (default) & Pa \quad $2.94210\times10^{13}$ \\
                     &                 & \code{'GPa'}      & $29421.0$ GPa \\
                     &                 & \code{'atm'}      & $2.90443\times10^{8}$ atm \\
\code{'efield'}     & $E_h/(e\,a_0)$  & (default) & V/m \quad $5.14221\times10^{11}$ \\
                     &                 & \code{'GV/m'}     & $514.221$ GV/m \\
                     &                 & \code{'V/A'}      & $51.4221$ V/\AA \\
\code{'mass'}       & $m_e$           & (default) & kg \quad $9.10938\times10^{-31}$ \\
                     &                 & \code{'Da'}       & $5.48580\times10^{-4}$ Da (AMU) \\
\code{'charge'}     & $e$             & (default) & C \quad $1.60218\times10^{-19}$ \\
\code{'magmom'}     & $\mu_B$         & (default) & J/T \quad $9.27401\times10^{-24}$ \\
\code{'temp'}       & $E_h \leftrightarrow$ K & (default) & K (via $k_BT$) \quad $315775.0$ \\
\bottomrule
\end{longtable}

\subsubsection*{Examples}

\begin{lstlisting}
% SCF energy: au -> kJ/mol and eV
AU_convert(-875.9322, 'energy', 'au2si', 'target', 'kJ/mol')
AU_convert(-875.9322, 'energy', 'au2si', 'target', 'eV')

% Relative Gibbs energy between two conformers
dG_au = en2.G - en1.G;
dG_kJ = AU_convert(dG_au, 'energy', 'au2si', 'target', 'kJ/mol');

% Vibrational zero-point energy (freq in cm-1 -> Eh)
zpe_au = AU_convert(1582.8/2, 'energy', 'si2au', 'target', 'cm-1');

% Dipole moment: au -> Debye (G16_dipole_polar with 'units','au')
dp   = G16_dipole_polar('V_E00t.out', 'units', 'au');
mu_D = AU_convert(dp.mu_tot, 'dipole', 'au2si', 'target', 'Debye');

% Polarisability tensor: au -> Angstrom^3 (volume polarisability)
dp       = G09_dipole_polar('INDACO_E025_ppDP.LOG');
alpha_A3 = AU_convert(dp.alpha_tensor, 'polar', 'au2si', 'target', 'Angstrom3');

% Force constants: au -> mDyne/Angstrom
nm       = G16_nmodes('V_E00t.out');
fc_mdyne = AU_convert(nm.frcconst, 'frcconst', 'au2si', 'target', 'mDyne/A');

% Geometry array: Angstrom -> Bohr
mol      = G16_structure('V_E00t.out');
xyz_bohr = AU_convert(mol.xyz, 'length', 'si2au', 'target', 'Angstrom');

% Thermal energy kBT at room temperature
kBT = AU_convert(298.15, 'temp', 'si2au');

% Hyperpolarisability: au -> esu
hp       = G16_hyperpolar('V_E00t.out');
beta_esu = AU_convert(hp.beta0.beta_vec, 'beta', 'au2si', 'target', 'esu');
\end{lstlisting}

% -----------------------------------------------------------------------------
\subsection{AU\_table}
\label{sec:gutil-au-table}

\begin{lstlisting}
AU_table                  % print complete table
AU_table('energy')        % print energy section only
AU_table('constants')     % print fundamental constants section
AU_table('polar')         % print polarisability section
\end{lstlisting}

Prints a formatted reference table of conversion factors and fundamental
physical constants to the MATLAB console. The filter argument is
case-insensitive and accepts the same quantity strings as
\code{AU\_convert} (\code{energy}, \code{length}, \code{dipole},
\code{polar}, \code{beta}, \code{gamma}, \code{force}, \code{frcconst},
\code{frequency}, \code{time}, \code{pressure}, \code{efield},
\code{mass}, \code{charge}, \code{magmom}, \code{temp}, plus
\code{constants} for the fundamental-constants section).

\begin{lstlisting}
>> AU_table('energy')

  -- ENERGY --
  1 Hartree (Eh)              = 4.3597447222e-18 J
                               = 2625.499639 kJ/mol
                               = 627.509474 kcal/mol
                               = 27.211386 eV
                               = 219474.6314 cm-1
                               = 27211.386 meV
  1 eV                        = 1.6021766340e-19 J
                               = 96.4853 kJ/mol
                               = 8065.5439 cm-1
  1 cm-1                      = 1.9864459e-23 J
                               = 4.55634e-06 Hartree
  1 kJ/mol                    = 3.80880e-04 Hartree
  1 kcal/mol                  = 1.59360e-03 Hartree
  1 kBT (298.15K)             = 9.44185e-04 Hartree  (= 2.4790 kJ/mol)
\end{lstlisting}

% -----------------------------------------------------------------------------
\subsection{G\_gaussian\_field\_convert}
\label{sec:gutil-gaussian-field-convert}

\begin{lstlisting}
out = G_gaussian_field_convert(N, 'g2phys')
out = G_gaussian_field_convert(N, 'g2phys', 'Unit', 'V/Ang')
out = G_gaussian_field_convert(value, 'phys2g', 'Unit', 'au')
\end{lstlisting}

Converts between Gaussian's \code{Field=X+N} route-keyword integer and a
physical electric-field strength. Gaussian's \code{Field} keyword (e.g.\
\code{Field=X+10}, \code{Field=X-5}) specifies a static electric field
whose magnitude, in atomic units, is the keyword's integer \(N\) times
\(0.0001\) --- i.e.\ \(N \times 10^{-4}\) a.u.

\begin{tbnote}
\textbf{Verified against the source, not assumed.} This convention was
confirmed directly from Gaussian's own keyword reference
(\url{https://gaussian.com/field/}): \emph{"N*0.0001 specifies the
magnitude of the field in atomic units"} (for the \(M{\pm}N\) format, the
one a plain dipole field like \code{Field=X+10} uses) --- quoted from the
source rather than assumed, since this exact convention is easy to get
wrong (a field-magnitude keyword invites exactly this kind of silent
off-by-a-power-of-ten mistake).
\end{tbnote}

\begin{argtabletwo}{Argument}{Description}{2.2cm}{X}
\code{value}     & Gaussian keyword integer \(N\) (direction \code{'g2phys'}) or a physical field strength (direction \code{'phys2g'}); may be a numeric array (applied elementwise) \\
\code{direction} & \code{'g2phys'} --- \(N \to\) physical field strength; \code{'phys2g'} --- physical field strength \(\to N\) \\
\end{argtabletwo}

\begin{argtable}{Option}{Default}{Description}{2.0cm}{2.0cm}{X}
\code{'Unit'}  & \code{'au'}  & \code{'au'} \textbar\ \code{'V/m'} \textbar\ \code{'V/cm'} \textbar\ \code{'V/Ang'} --- the physical unit on the non-Gaussian side, for either direction \\
\code{'Print'} & \code{true}  & Print a one-line summary \\
\end{argtable}

\begin{tbnote}
\textbf{Rounding on \code{'phys2g'}.} Gaussian's keyword can only express
multiples of \(0.0001\) a.u., so converting an arbitrary physical value
back to \(N\) necessarily rounds to the nearest integer. A warning is
printed whenever that rounding changes the value by more than 0.5\%, so a
value that does not correspond to a "clean" \(N\) is never silently
misrepresented.

\textbf{Consistency with \code{AU\_convert}.} The atomic unit of electric
field (\(E_h/(e\,a_0)\)) is computed from the same CODATA~2022 constants
\code{AU\_convert} itself uses (Section~\ref{sec:gutil-au-convert}), not a
separately rounded literal --- cross-checked to match
\code{AU\_convert(\ldots, 'efield', \ldots, 'target', 'V/A')} exactly.
\end{tbnote}

\begin{lstlisting}
G_gaussian_field_convert(10, 'g2phys')                    % N=10 -> 0.0010 au
G_gaussian_field_convert(10, 'g2phys', 'Unit', 'V/Ang')   % N=10 -> 0.0514 V/Ang
G_gaussian_field_convert(0.0025, 'phys2g', 'Unit', 'au')  % 0.0025 au -> N=25
\end{lstlisting}

See also \code{AU\_convert} (Section~\ref{sec:gutil-au-convert}).

% -----------------------------------------------------------------------------
\subsection{Integration with the G09/G16 Toolbox}

\code{AU\_convert} operates on plain numeric arrays and is fully
compatible with all structs returned by \code{G09\_*}/\code{G16\_*}
functions. The table below lists the most common conversions needed
after parsing a Gaussian output file.

\begin{argtable}{Field}{Quantity / target}{Typical use}{2.6cm}{2.6cm}{X}
\code{en.SCF} / \code{en.G} / \code{en.H} & \code{'energy'} $\to$ \code{'kJ/mol'} & Relative energies, reaction profiles \\
\code{en.SCF}                             & \code{'energy'} $\to$ \code{'eV'}     & Band gaps, ionisation potentials \\
\code{sp.freq}                            & already cm$^{-1}$                    & No conversion needed \\
\code{nm.frcconst}                        & \code{'frcconst'} $\to$ \code{'mDyne/A'} & Comparison with experimental force constants \\
\code{mol.xyz}                            & already \AA                          & Convert to Bohr (\code{'length'}, \code{si2au}) only for QM codes that need it \\
\code{dp.mu\_tot}                         & \code{'dipole'} $\to$ \code{'Debye'} & Dipole magnitude (au by default) \\
\code{dp.alpha\_tensor}                   & \code{'polar'} $\to$ \code{'Angstrom3'} & Volume polarisability for SERS/NLO analysis \\
\code{hp.beta0.beta\_vec}                 & \code{'beta'} $\to$ \code{'esu'}     & Hyperpolarisability for NLO comparison \\
$k_BT$ at $T$                             & \code{'temp'}, \code{si2au}           & Boltzmann weights in conformer analysis \\
\code{en.S\_JmolK}                        & already J~mol$^{-1}$~K$^{-1}$        & No conversion needed \\
\end{argtable}

\begin{tbnote}
Fields already returned in physical (non-au) units by G09/G16 Toolbox
functions: \code{sp.freq} (cm$^{-1}$), \code{en.S\_JmolK}
(J~mol$^{-1}$~K$^{-1}$), \code{en.ZPE\_kJ} (kJ/mol), \code{mol.xyz} (\AA),
\code{sp.IR} (KM/Mole), \code{sp.Raman} (\AA$^4$/AMU). These do not
require \code{AU\_convert}.
\end{tbnote}

\subsubsection*{Worked example: NLO property analysis}

\begin{lstlisting}
% Load properties from G16 output
dp = G16_dipole_polar('V_E00t.out');
hp = G16_hyperpolar('V_E00t.out');

% Convert to common NLO units
mu_D     = AU_convert(dp.mu_tot,         'dipole', 'au2si', 'target', 'Debye');
alpha_A3 = AU_convert(dp.alpha_iso,      'polar',  'au2si', 'target', 'Angstrom3');
beta_esu = AU_convert(hp.beta0.beta_vec, 'beta',   'au2si', 'target', 'esu');

fprintf('mu      = %.4f D\n',   mu_D)
fprintf('alpha   = %.3f A^3\n', alpha_A3)
fprintf('|beta|  = %.3e esu\n', beta_esu)
\end{lstlisting}

\subsubsection*{Worked example: Boltzmann-weighted conformer ensemble}

\begin{lstlisting}
T   = 298.15;                            % K
kBT = AU_convert(T, 'temp', 'si2au');    % Eh

files = {'conf1.out', 'conf2.out', 'conf3.out'};
G = zeros(1, numel(files));
for i = 1:numel(files)
    en   = G16_energy(files{i});
    G(i) = en.G;
end
dG    = G - min(G);                      % relative Gibbs energy (Eh)
w     = exp(-dG / kBT);                  % Boltzmann weights
w     = w / sum(w);                      % normalised
dG_kJ = AU_convert(dG, 'energy', 'au2si', 'target', 'kJ/mol');
\end{lstlisting}

% -----------------------------------------------------------------------------
\subsection{CODATA 2022 fundamental constants}

All conversion factors in \code{AU\_convert}/\code{AU\_table} are derived
from the CODATA 2022 values below. Values marked \emph{(exact)} are fixed
by definition in the 2019 SI redefinition.

\begin{argtable}{Constant}{Symbol}{Value}{3.2cm}{1.6cm}{X}
Bohr radius              & $a_0$    & $5.29177210544\times10^{-11}$ m \\
Hartree energy           & $E_h$    & $4.3597447222060\times10^{-18}$ J \\
Electron mass            & $m_e$    & $9.1093837139\times10^{-31}$ kg \\
Proton mass              & $m_p$    & $1.67262192595\times10^{-27}$ kg \\
Elementary charge        & $e$      & $1.602176634\times10^{-19}$ C (exact) \\
Planck constant          & $h$      & $6.62607015\times10^{-34}$ J$\cdot$s (exact) \\
Reduced Planck constant  & $\hbar$  & $1.054571817\times10^{-34}$ J$\cdot$s (exact) \\
Speed of light           & $c$      & $299\,792\,458$ m/s (exact) \\
Boltzmann constant       & $k_B$    & $1.380649\times10^{-23}$ J/K (exact) \\
Avogadro constant        & $N_A$    & $6.02214076\times10^{23}$ mol$^{-1}$ (exact) \\
Vacuum permittivity      & $\varepsilon_0$ & $8.8541878188\times10^{-12}$ F/m \\
Bohr magneton            & $\mu_B$  & $9.2740100657\times10^{-24}$ J/T \\
Nuclear magneton         & $\mu_N$  & $5.0507837393\times10^{-27}$ J/T \\
\end{argtable}

\begin{tbnote}
Source: NIST CODATA 2022, \url{https://physics.nist.gov/cuu/Constants/}
(published November 2024). The four exact constants ($e$, $h$, $k_B$,
$N_A$) have been fixed since the 2019 SI redefinition.
\end{tbnote}


% -----------------------------------------------------------------------------
\subsection{Python port (\code{g16\_au\_convert} / \code{g16\_au\_table} / \code{g16\_gaussian\_field\_convert})}
\label{sec:gutil-au2si-python}

Unlike \code{MOSurface}/\code{RamanBrowser} (\S\ref{sec:gutil-mo-surface-python},
\S\ref{sec:gutil-raman-browser-python}), au2SI's Python port is folded
directly into the \code{G16parser} package itself, not kept as a
separate companion: it has no dependency of its own beyond plain
\code{numpy}, and is exactly the kind of general, version-agnostic
utility (like \code{g16\_read\_input}) that belongs in the core Python
package rather than alongside it.

\begin{lstlisting}
import G16parser as g16

g16.g16_au_convert(-875.932, 'energy', 'au2si', target='kJ/mol')
g16.g16_au_convert(6.3347, 'dipole', 'si2au', target='Debye')
g16.g16_au_table('energy')
g16.g16_gaussian_field_convert(10, 'g2phys', unit='V/Ang')
\end{lstlisting}

\begin{argtabletwo}{Function}{Signature}{3.0cm}{X}
\code{g16\_au\_convert}  & \code{g16\_au\_convert(value, quantity, direction, target='', verbose=None)} -- same quantity strings as \code{AU\_convert} (\S\ref{sec:gutil-au-convert}); \code{quantity='help'} prints the reference table and returns \code{None} \\
\code{g16\_au\_table}    & \code{g16\_au\_table(section='')} -- same section filter as \code{AU\_table} (\S\ref{sec:gutil-au-table}), plus \code{'constants'} \\
\code{g16\_gaussian\_field\_convert} & \code{g16\_gaussian\_field\_convert(value, direction, unit='au', verbose=True)} \\
\end{argtabletwo}

\begin{tbnote}
\textbf{Validation.} Cross-checked directly against the MATLAB
original, not just internally self-consistent: 12 quantity/target
combinations spanning every category (energy, length, dipole, polar,
beta, force, efield, temp, mass, pressure, frequency) agree with
\code{AU\_convert.m} to $\sim$1e-13--1e-14 relative difference (limited
by \code{\%.12e} text-transfer precision, not a real discrepancy); both
\code{G\_gaussian\_field\_convert.m} directions (\code{g2phys}/
\code{phys2g}) match exactly. A 23-case \code{pytest} regression suite
(\code{tests/test\_au\_convert.py}) covers every function, including
array input, the \code{'help'} mode, error paths, and the
rounding-warning path on \code{g16\_gaussian\_field\_convert}.
\end{tbnote}

\newpage

\section{MOSurface: real-space orbital, density, and ESP isosurfaces}
\label{sec:gutil-mo-surface}
% =============================================================================
\code{MOSurface/} provides four GaussView-equivalent real-space
renderers, all built directly from the raw \code{.fchk} basis-set data
(or, for the fourth, directly from a saved \code{.cube} file) with no
toolbox or third-party dependency beyond MATLAB's own built-in
\code{isosurface}/\code{isonormals}/\code{patch}/\code{contourf}:
\code{G\_draw\_mo\_surface} (\S\ref{sec:gutil-mo-surface-fn}, a chosen
molecular orbital's $\pm$isosurface or 2D nodal map, GaussView's
``Molecular Orbitals''), \code{G\_draw\_density\_surface}
(\S\ref{sec:gutil-density-surface}, the total electron density, a
density \emph{difference} between two calculations, or the \emph{spin}
density of an unrestricted calculation, GaussView's ``Total
Density''/``Spin Density''), \code{G\_draw\_esp\_surface}
(\S\ref{sec:gutil-esp-surface}, the electrostatic potential mapped onto
the density envelope, GaussView's ``ESP mapped on electron density''),
and \code{G\_draw\_cube\_surface} (\S\ref{sec:gutil-cube-surface}, which
re-renders an isosurface straight from an already-saved \code{.cube}
file -- no \code{.fchk}, no basis-set re-evaluation). The first three
can each export the underlying volumetric grid to a Gaussian-format
\code{.cube} file via \code{'SaveCube'} (see each subsection);
\code{G\_draw\_cube\_surface} is the load-and-replot companion that
reads one back (or a real \code{cubegen}/GaussView-generated \code{.cube}
just as well). Every other \code{MOSurface/} file (\code{g\_}-prefixed,
lowercase) is an internal helper, never called directly by the user ---
only the four \code{G\_}-prefixed functions above are the public API. A
Python port of all four, \code{MOSurface/python/}, also exists (see
\S\ref{sec:gutil-mo-surface-python}).

% -----------------------------------------------------------------------------
\subsection{G\_draw\_mo\_surface}
\label{sec:gutil-mo-surface-fn}

\begin{lstlisting}
h = G_draw_mo_surface(data, mo_index)
h = G_draw_mo_surface(data, mo_index, 'Name', Value, ...)
\end{lstlisting}

\code{data} is the full struct returned by \code{G09\_fchk\_read}/
\code{G16\_fchk\_read} (must contain \code{.filename}, \code{.mol},
\code{.Nbasis}, \code{.alpha\_MO\_coeff}, \code{.xyz\_bohr}).
\code{mo\_index} selects which orbital to render: either a positive
integer (1-based column of
\code{reshape(data.alpha\_MO\_coeff, data.Nbasis, data.Nbasis)}), or the
string \code{'HOMO'}, \code{'LUMO'}, \code{'HOMO-n'}, \code{'LUMO+n'}
(resolved against \code{data.HOMO\_idx}).

\begin{argtable}{Option}{Default}{Description}{2.6cm}{2.4cm}{X}
\code{'Mode'}         & \code{'surface'}           & \code{'surface'}: 3D $\pm$isosurface (below). \code{'contour'}: a 2D filled-contour amplitude map on a single plane, with the nodal line drawn explicitly -- see \S\ref{sec:gutil-mo-surface-contour} \\
\code{'Plane'}        & \code{'auto'}              & (\code{'contour'} only) \code{'auto'}: best-fit plane through all atoms (PCA). \code{'xy'}/\code{'xz'}/\code{'yz'}: a coordinate plane through the centroid \\
\code{'PlaneOffset'}  & \code{0} \AA               & (\code{'contour'} only) shifts the evaluation plane along its own normal -- see \S\ref{sec:gutil-mo-surface-contour} \\
\code{'IsoValue'}     & \code{0.02}                & (\code{'surface'} only) Wavefunction-amplitude isosurface level (a.u.), GaussView's own default. Both $+$IsoValue and $-$IsoValue lobes are drawn. \\
\code{'GridSpacing'}  & \code{0.15} Bohr           & Real-space grid spacing \\
\code{'Padding'}      & \code{4.0} Bohr            & Grid/plane padding around the molecule's bounding box; auto-expanded if needed in \code{'surface'} mode (see below) \\
\code{'SaveCube'}     & \code{''} (off)            & (\code{'surface'} only) filename; if given, writes the signed MO-amplitude grid already computed for the isosurface (no extra evaluation) to a Gaussian-format MO \code{.cube} file (negative atom count plus the orbital-index line, matching real \code{cubegen} MO cubes) \\
\code{'PosColor'}     & blue                       & Positive-amplitude colour \\
\code{'NegColor'}     & red                        & Negative-amplitude colour \\
\code{'FaceAlpha'}    & \code{0.55}                & (\code{'surface'} only) Lobe transparency; used as-is with \code{'SurfaceStyle','transparent'} (the default), overridden to \code{1.0} with \code{'solid'} unless given explicitly (which always wins), unused with \code{'grid'} \\
\code{'SurfaceStyle'} & \code{'transparent'}       & (\code{'surface'} only) \code{'solid'} \textbar\ \code{'transparent'} \textbar\ \code{'grid'} -- as in GaussView. \code{'grid'} draws the isosurface as a wireframe (mesh edges only, no face fill) coloured by \code{PosColor}/\code{NegColor} instead of a filled patch. Neither this function nor \code{G\_draw\_density\_surface} decimates the \code{'surface'} mesh, so a fine \code{'GridSpacing'} can give a very dense wireframe -- coarsen it for a readable grid render \\
\code{'ShowMolecule'} & \code{true}                & \code{'surface'}: overlay the CPK model via whichever of \code{G09\_draw\_molecule}/\code{G16\_draw\_molecule} is on the path. \code{'contour'}: overlay a lightweight 2D atom/bond sketch instead \\
\code{'ShowLabels'}   & \code{false}               & (\code{'surface'} only) Atom labels, forwarded to the molecule drawer; always shown in \code{'contour'} mode \\
\code{'Ax'}           & \code{[]} (new figure)     & Existing axes handle \\
\code{'Verbose'}      & \code{true}                & Print grid size and timing \\
\end{argtable}

Returns a struct \code{h}. In \code{'surface'} mode: \code{h.pos}/\code{h.neg},
the positive/negative lobe \code{patch} handles (\code{[]} if that lobe was
not found at \code{IsoValue}). In \code{'contour'} mode: \code{h.contour}
(filled-contour handle) and \code{h.zero} (nodal-line handle, \code{[]} if
the amplitude never changes sign on the plane).

\subsubsection*{The physics}

Each contracted Gaussian-type basis function is evaluated directly on the
grid from its primitives,
\[
N(\alpha,l,m,n) = \left(\frac{2\alpha}{\pi}\right)^{3/4}
\sqrt{\frac{(4\alpha)^{\,l+m+n}}{(2l-1)!!\,(2m-1)!!\,(2n-1)!!}},
\]
the standard Cartesian Gaussian-type-orbital (GTO) normalization
constant~[\ref{ref:Boys1950}]. \textbf{Pure D/F/G shells} are built from the
individually-normalized Cartesian components combined into real solid
harmonics using the polynomial coefficients of
Ribaldone \& Desmarais~[\ref{ref:Ribaldone2024}] (a modern, explicit
re-derivation of the classic Schlegel \& Frisch
transformation~[\ref{ref:SchlegelFrisch1995}]); the overall normalization of
each pure component is \emph{computed at runtime}, not hard-coded, from
the exact combinatorial self-overlap of same-centre Cartesian Gaussians
--- deliberately avoiding hand-derived constants for the higher (F, G)
shells, after an arithmetic slip was caught while hand-checking the
simpler D case during development. \textbf{Cartesian D/F/G} shells use
the corresponding per-component normalization ratio, computed the same
way. \textbf{Shell-component ordering} (including the native Gaussian
\code{.fchk} reversal for Cartesian shells with $L\!\geq\!4$) follows the
convention documented by the IOData project's Gaussian \code{.fchk}
reader~[\ref{ref:IOData}]. All of this is summed directly into a single
\code{[Ngrid x 1]} accumulator, shell by shell, never forming the full
\code{[Ngrid x Nbasis]} basis matrix, so memory stays bounded regardless
of basis-set size.

\begin{tbnote}
\textbf{Supported shells: S, P, SP, D, F, G (pure or Cartesian).} A basis
set containing H (or higher) shells raises a clear, immediate error
rather than silently omitting or mis-rendering those basis functions ---
there is no partial/best-effort rendering mode. F/G support was validated
on a real Au-containing basis set (pure F and pure G functions on the Au
atom): every molecular orbital with significant weight on those shells
gave a self-normalization integral $\int|\psi|^2\,dV$ of 0.998--1.000.
\end{tbnote}

\begin{tbnote}
\textbf{Deep-core and some high-virtual orbitals can show grid
artifacts.} A uniform rectangular grid under-resolves very tight
(high-exponent) primitives: for such orbitals, the self-normalization
integral can be non-monotonic and fail to converge cleanly even as
\code{'GridSpacing'} is refined (root cause: a narrow, sharply-peaked
Gaussian is intrinsically hard for naive Riemann-sum quadrature to
capture, not a bug in the evaluator). Verified on real data: 10 of 12
ordinary valence/frontier orbitals sampled across a full basis set gave
0.993--1.000; the deepest core orbital and the highest virtual (both
dominated by very tight primitives) did not converge under grid
refinement. This does not affect the chemically relevant use case of
rendering valence/frontier (HOMO/LUMO-region) orbitals.
\end{tbnote}

\begin{tbnote}
\textbf{Automatic padding.} Some orbitals (typically diffuse high
virtuals) extend well beyond the default 4~Bohr padding and would
otherwise be visibly clipped --- a flat-cut isosurface face at the edge
of the plotted box. After each isosurface extraction, the vertices are
checked against the grid boundary; if any lie within 1.5 grid cells of
it, \code{'Padding'} is multiplied by 1.6 and the whole
grid/evaluation/extraction is redone, up to 3 times (subject to the same
30-million-point grid-size cap that also applies to a manually-set
\code{'Padding'}). If an explicit \code{'Padding'} was given, this
auto-expansion is skipped and a warning recommends increasing it instead.
\end{tbnote}

\begin{lstlisting}
data = G16_fchk_read('4-NTP.fchk');
G_draw_mo_surface(data, 'HOMO');
G_draw_mo_surface(data, 'LUMO+1', 'IsoValue', 0.03, 'GridSpacing', 0.12);
G_draw_mo_surface(data, 'HOMO', 'SaveCube', 'homo.cube');   % also export
G_draw_mo_surface(data, 'HOMO', 'SurfaceStyle', 'grid', 'GridSpacing', 0.3);
\end{lstlisting}

\subsubsection*{'Mode','contour': 2D nodal maps}
\label{sec:gutil-mo-surface-contour}

Instead of a 3D isosurface, \code{'Mode','contour'} evaluates the MO on a
single 2D plane and renders a filled amplitude map (diverging colour
scale built from \code{'NegColor'}/\code{'PosColor'}, pinned to white at
zero regardless of the actual data range) with the nodal
(zero-amplitude) line drawn explicitly in black, plus a lightweight 2D
atom/bond sketch for reference. It is much cheaper than a 3D grid, and
often clearer for reading off the nodal pattern of a planar/near-planar
$\pi$ system than a 3D isosurface viewed from an angle -- compare the
worked example below with the equivalent \code{'Mode','surface'} render
of the same orbital.

\begin{tbnote}
\textbf{A $\pi$-symmetry orbital is exactly zero on its own molecular
plane.} By symmetry, a $\pi$ MO's amplitude is antisymmetric about the
plane it is built perpendicular to, so \code{'Plane','auto'} (the
best-fit plane through the atoms) together with the default
\code{'PlaneOffset',0} shows essentially nothing for a typical aromatic
HOMO/LUMO (verified: max amplitude $\sim 10^{-9}$--$10^{-29}$~a.u., pure
numerical noise, on two different real aromatic test molecules). This is
a property of the orbital, not a bug. \code{'PlaneOffset'} shifts the
evaluation plane along its own normal by the given number of \AA{}
(typically 0.5--1.0 for a $\pi$ system), probing a plane parallel to the
molecular plane where the $\pi$ density actually peaks; the atom sketch
is always drawn at the true, unshifted positions, for reference. A
warning is printed whenever the resulting amplitude is essentially zero
everywhere, naming this as the likely cause.
\end{tbnote}

\begin{lstlisting}
% A HOMO with pi character: offset=0 shows (near) nothing by symmetry
G_draw_mo_surface(data, 'HOMO', 'Mode', 'contour');

% The pi density itself, 0.7 Angstrom above the ring
G_draw_mo_surface(data, 'HOMO', 'Mode', 'contour', 'PlaneOffset', 0.7);
\end{lstlisting}

See also \code{G09\_fchk\_read}\slash\code{G16\_fchk\_read},
\code{G09\_draw\_molecule}\slash\code{G16\_draw\_molecule} (main toolbox
manual).

% -----------------------------------------------------------------------------
\subsection{G\_draw\_density\_surface}
\label{sec:gutil-density-surface}

\begin{lstlisting}
h = G_draw_density_surface(data)
h = G_draw_density_surface(data, 'Name', Value, ...)
h = G_draw_density_surface(data, 'CompareTo', data2, ...)
h = G_draw_density_surface(data, 'SpinDensity', true, ...)
\end{lstlisting}

Renders the total electron density, or (given a second structure via
\code{'CompareTo'}) a density \emph{difference} map, or (for an
unrestricted calculation, via \code{'SpinDensity'}) the \emph{spin}
density $\rho_\alpha-\rho_\beta$ -- a real-space equivalent of
GaussView's ``Total Density''/``Spin Density'' surface types. It shares
its real-space evaluation engine with \code{G\_draw\_mo\_surface}
(\S\ref{sec:gutil-mo-surface-fn}): the density is electron-count-exact,
not a heuristic, built from
\[
\rho(\mathbf r) = \sum_i n_i\,|\psi_i(\mathbf r)|^2
\]
summed over all \emph{occupied} orbitals of the relevant spin channel
($n_i=2$ per orbital for the total/difference density, closed shell;
$n_i=1$ per orbital, alpha and beta separately, for
\code{'SpinDensity'}), evaluated in a \emph{single} pass through the
basis shells for every occupied orbital at once
(\code{g\_eval\_mo\_on\_grid} accepts an \code{[Nbasis x K]} coefficient
matrix, not just a single column, and returns an \code{[Ngrid x K]}
result) rather than looping the whole per-shell evaluation once per
orbital.

Most options (\code{'Mode'}, \code{'Plane'}, \code{'PlaneOffset'},
\code{'GridSpacing'}, \code{'Padding'} with the same auto-expansion,
\code{'FaceAlpha'}, \code{'SurfaceStyle'}, \code{'ShowMolecule'},
\code{'ShowLabels'}, \code{'AtomScale'}, \code{'Ax'}, \code{'Verbose'})
are identical to \code{G\_draw\_mo\_surface}. The differences:

\begin{argtabletwo}{Option}{Description}{2.4cm}{X}
\code{'CompareTo'} & A second data struct (same requirements as \code{data}). If given, renders $\rho(\text{CompareTo}) - \rho(\text{data})$ instead of the total density of \code{data} alone -- e.g. field-on minus field-off, to see where charge density increases/decreases. Both densities are evaluated on the same grid/plane, built from \code{data}'s geometry. Mutually exclusive with \code{'SpinDensity'} (default: \code{[]}, total density only) \\
\code{'SpinDensity'} & Logical; if \code{true}, renders $\rho_\alpha(\mathbf r)-\rho_\beta(\mathbf r)$ of \code{data} alone -- see the tbnote below for the open-shell requirement. Mutually exclusive with \code{'CompareTo'} (default: \code{false}) \\
\code{'IsoValue'} & Default \code{0.001} e/Bohr$^3$ -- the classic isodensity value conventionally used to define a molecule's outer envelope/``size'', appropriate for the \emph{total} density. A difference or spin density is usually one or more orders of magnitude smaller; if no isosurface is found, a warning suggests reducing \code{'IsoValue'} \\
\code{'SaveCube'} & (\code{'surface'} only) filename; if given, writes the density (total, \emph{difference} with \code{'CompareTo'}, or \emph{spin} with \code{'SpinDensity'}) grid already computed for the isosurface to a Gaussian-format \code{.cube} file, at no extra evaluation cost (default: \code{''}, off) \\
\code{'PosColor'} & The (only) density lobe colour without \code{'CompareTo'}/\code{'SpinDensity'}; the positive-difference or alpha-excess lobe colour with either \\
\code{'NegColor'} & Only used with \code{'CompareTo'} (negative-difference lobe) or \code{'SpinDensity'} (beta-excess lobe) \\
\end{argtabletwo}

Returns \code{h.pos}/\code{h.neg}/\code{h.contour}/\code{h.zero}, with
the same meaning as \code{G\_draw\_mo\_surface}; without
\code{'CompareTo'}/\code{'SpinDensity'}, \code{h.neg} and \code{h.zero}
are always \code{[]} (a total density is single-signed, so there is no
negative lobe and no zero/nodal line to trace).

\begin{tbnote}
\textbf{Closed-shell only for the total density/\code{'CompareTo'}; a
dedicated open-shell path for \code{'SpinDensity'}.} The total-density
and \code{'CompareTo'} paths require \code{data.Nalpha == data.Nbeta}
(and the same for \code{'CompareTo'}, if given) and raise a clear error
otherwise, rather than silently building a wrong electron count for an
open-shell (UHF/UKS) calculation. \code{'SpinDensity'} is the dedicated
open-shell path instead: it does \emph{not} gate on
\code{Nalpha==Nbeta} (not itself a restricted/unrestricted signal -- a
broken-symmetry biradical singlet, for instance, can have
\code{Nalpha==Nbeta} with genuinely different alpha/beta orbitals), but
instead re-reads the raw \code{"Beta MO coefficients"} \code{.fchk}
section directly (in the same self-contained spirit as
\code{g\_read\_aobasis\_from\_fchk}, no core-toolbox change), which
Gaussian only ever writes for a genuinely unrestricted calculation -- a
restricted (RHF/RKS) \code{.fchk} has identical alpha/beta orbitals by
construction (zero spin density everywhere) and this section is simply
absent, raising a clear error rather than a confusing empty or all-zero
plot. Each spin channel then uses occupation~1 per orbital (not~2), via
\code{g\_eval\_density\_on\_grid}'s optional \code{occ\_factor}
argument.

Validated first on synthetic test cases (no genuine UHF/UKS
\code{.fchk} was available in this repository at initial development
time): a beta channel copied byte-for-byte from the alpha channel gives
\emph{exactly} zero spin density at every sampled point (both a direct
pointwise check and the expected ``no isosurface found'' warning); a
beta channel using a shifted occupied window (still \code{Nbeta}
orthonormal columns of the same AO basis, but a genuinely different
occupied subspace) gives a clearly nonzero, spatially-varying spin
density at 78\% of sampled points.

Subsequently validated end-to-end against a real Gaussian UKS
calculation: the methyl radical (CH$_3^{\bullet}$, UB3LYP/6-31G(d), doublet,
$N_\alpha=5$, $N_\beta=4$). Integrating the rendered spin density over
the same real-space grid gives $\int(\rho_\alpha-\rho_\beta)\,dV =
1.0000$, matching $N_\alpha-N_\beta=1$ (the one unpaired electron)
essentially exactly -- a quantitative check, not just a qualitative
one. The isosurface shape matches the textbook picture for this
system's SOMO (a carbon $p_z$ orbital, perpendicular to the molecular
plane): the positive (alpha-excess) lobe spans $z\in[-1.14,+1.19]$~\AA{}
-- i.e.\ two separate lobes straddling the (planar) molecule, not a
single blob centred on it -- while a smaller negative (beta-excess)
lobe sits close to the C-H bonding region ($xy$-radius from the carbon
up to 1.44~\AA{}), consistent with the well-known spin polarization at
the hydrogens of CH$_3^{\bullet}$ (the same effect responsible for their
negative $^1$H hyperfine coupling constant in EPR).
\end{tbnote}

\begin{tbnote}
\textbf{A raw density plot on a plane through the nuclei is dominated by
sharp nuclear cusps.} The electron density peaks very sharply at each
nucleus -- one to several orders of magnitude larger than the
chemically-interesting bonding/valence-region values -- so a
\code{'Mode','contour'} colour scale naively set to the data's true
maximum saturates the whole map to near-white except at a few nuclear
pinpoints, hiding the region of actual interest. \code{'Mode','contour'}
therefore caps the colour scale at the 98th percentile of the plotted
values instead (values above the cap simply saturate to the
top-of-scale colour); \code{'Mode','surface'} is unaffected, since an
isosurface at a single, comparatively low \code{'IsoValue'} (e.g. the
default 0.001) sits well outside the nuclear cusp region already.
\end{tbnote}

\begin{lstlisting}
data = G16_fchk_read('4-NTP.fchk');
G_draw_density_surface(data);                    % total density, 0.001 envelope
G_draw_density_surface(data, 'Mode', 'contour');  % 2D map, capped colour scale

data1 = G16_fchk_read('nofield.fchk');
data2 = G16_fchk_read('field_x025.fchk');
G_draw_density_surface(data1, 'CompareTo', data2, 'IsoValue', 0.0005);
G_draw_density_surface(data, 'SaveCube', 'density.cube');   % also export
G_draw_density_surface(data, 'SurfaceStyle', 'grid', 'GridSpacing', 0.3);

radical = G16_fchk_read('radical_uks.fchk');   % unrestricted (UHF/UKS)
G_draw_density_surface(radical, 'SpinDensity', true, 'IsoValue', 0.002);
\end{lstlisting}

See also \code{G\_draw\_mo\_surface} (\S\ref{sec:gutil-mo-surface-fn}),
\code{G09\_fchk\_read}\slash\code{G16\_fchk\_read}.

% -----------------------------------------------------------------------------
\subsection{G\_draw\_esp\_surface}
\label{sec:gutil-esp-surface}

\begin{lstlisting}
h = G_draw_esp_surface(data)
h = G_draw_esp_surface(data, 'Name', Value, ...)
h = G_draw_esp_surface(data, 'CompareTo', data2, ...)
\end{lstlisting}

Renders the molecular electrostatic potential (ESP) mapped onto the
total-density isosurface (built via \code{G\_draw\_density\_surface}'s
machinery) -- a real-space equivalent of GaussView's classic ``ESP
mapped on electron density'' surface, showing electrophilic
(\code{PosColor}, blue, $V>0$) vs.\ nucleophilic (\code{NegColor}, red,
$V<0$) character:
\[
V(\mathbf r) = \sum_A \frac{Z_A}{|\mathbf r-\mathbf R_A|}
- \int \frac{\rho(\mathbf r')}{|\mathbf r-\mathbf r'|}\,d\mathbf r'
\]
$Z_A$ is read from the \code{.fchk} \texttt{"Nuclear charges"} section,
\emph{not} the atomic number: for an ECP (effective core potential) atom
this is the reduced effective charge seen by the explicit electrons --
verified on a real Au-ECP test file, \texttt{"Nuclear charges"} $=19$ for
Au vs.\ atomic number 79 (a 60-electron core ECP); using the atomic
number would silently give a badly wrong nuclear term for any such atom.

Unlike \code{G\_draw\_mo\_surface}/\code{G\_draw\_density\_surface}, the
electronic term is a genuine Coulomb integral, not a simple point
evaluation. By default it is evaluated \emph{exactly} (up to
double-precision/basis-set truncation), via the analytic
McMurchie-Davidson Hermite-Gaussian recursion~[\ref{ref:McMurchieDavidson1978},
\ref{ref:Helgaker2000}] -- treating each surface vertex as an arbitrary
field point in the same integral family as the nuclear-attraction term,
built from Hermite expansion coefficients $E_t^{ij}$, Hermite Coulomb
integrals $R_{tuv}$, and the Boys function $F_n(x)$ (evaluated in closed
form via the regularized incomplete gamma function). This only supports
S, P, SP, and pure D shells; for a basis with Cartesian D or F/G shells
(or if \code{'ESPMethod'} is explicitly \code{'numeric'}), it falls back
to the older \emph{numerical} Coulomb-grid-sum method (see the tbnote
below), with a printed notice if \code{'Verbose'}. The full derivation,
and the validation methodology and results summarized below, are in
\code{Theory\_Vibrations\_Polarizability.tex}, Part~III.

\begin{argtable}{Option}{Default}{Description}{2.8cm}{2.4cm}{X}
\code{'IsoValue'}       & \code{0.001}        & Density isovalue defining the surface shape (as in \code{G\_draw\_density\_surface}) \\
\code{'GridSpacing'} / \code{'Padding'} & \code{0.15} / \code{4.0} Bohr & Surface-\emph{shape} grid (auto-expanding padding, as in \code{G\_draw\_mo\_surface}) \\
\code{'MaxVertices'}    & \code{3000}         & The density mesh is decimated (\textsc{Matlab}'s built-in \code{reducepatch}) to at most this many vertices before the ESP evaluation, since ESP is smooth and a full-resolution mesh (often $>$10000 vertices) is unnecessary and would make it far slower than needed \\
\code{'CompareTo'}      & \code{[]}           & A second data struct (same requirements as \code{data}). If given, renders $V(\text{CompareTo}) - V(\text{data})$ instead of the absolute ESP of \code{data} alone -- e.g.\ field-on minus field-off. Both ESPs are evaluated independently (own analytic/numerical fallback, nuclear charges, density matrix) at the SAME vertices, built from \code{data}'s density-surface shape only, so the comparison is physically meaningful only if the two structures share the same atomic coordinate frame \\
\code{'ESPMethod'}      & \code{'auto'}       & \code{'auto'}: analytic when the basis is within scope, else numerical fallback (with notice). \code{'analytic'}: force analytic, erroring on an out-of-scope basis. \code{'numeric'}: force the numerical method regardless of basis \\
\code{'ESPGridSpacing'} & \code{0.20} Bohr    & \emph{Only used by the numerical method} (auto fallback or forced): spacing of the separate density grid used for the Coulomb sum -- see the accuracy note below \\
\code{'ESPPadding'}     & \code{6.0} Bohr     & \emph{Numerical method only}: padding for that grid; empirically far less important than \code{'ESPGridSpacing'} \\
\code{'DistanceFloor'}  & \code{0.5}$\times$\code{'ESPGridSpacing'} & \emph{Numerical method only}: minimum distance used in the Coulomb sum, avoiding a spurious spike when a grid point lands very close to a vertex \\
\code{'SaveCube'}       & \code{''} (off)     & Filename; if given, \emph{additionally} evaluates ESP on a separate, dedicated volumetric grid (see \code{'CubeSpacing'}/\code{'CubePadding'} -- NOT the surface-shape grid, NOT just the decimated vertices) and writes it to a Gaussian-format \code{.cube} file. This is a genuinely extra computation: ESP is otherwise only ever evaluated at the (few-thousand-at-most) decimated vertices \\
\code{'CubeSpacing'}    & \code{0.30} Bohr    & Grid spacing for \code{'SaveCube'}, coarser by default than \code{'GridSpacing'}/\code{'ESPGridSpacing'} -- with the analytic method, cost scales roughly with $N_\text{basis}^2$ per point, so a fine cube over many points can be genuinely slow (a 2-million-point safety cap raises a clear error rather than starting an impractically long run) \\
\code{'CubePadding'}    & \code{4.0} Bohr     & Padding for \code{'SaveCube'} \\
\code{'PosColor'} / \code{'NegColor'} & blue / red & Standard ESP convention (electrophilic / nucleophilic); with \code{'CompareTo'}, positive/negative \emph{difference} instead \\
\code{'FaceAlpha'}      & \code{1.0}          & Used as-is with \code{'SurfaceStyle','solid'} (the default here, unlike \code{G\_draw\_mo\_surface}/\code{G\_draw\_density\_surface}'s \code{'transparent'} default -- an ESP map is conventionally an opaque painted surface); overridden to \code{0.55} with \code{'transparent'} unless given explicitly; unused with \code{'grid'}. Also auto-lowered to \code{0.6} with \code{'ShowMolecule',true} + \code{'solid'} (see that entry) \\
\code{'SurfaceStyle'}   & \code{'solid'}      & \code{'solid'} \textbar\ \code{'transparent'} \textbar\ \code{'grid'} -- as in GaussView. \code{'grid'} draws a wireframe coloured by \emph{interpolating the ESP colormap along the edges} (\code{EdgeColor,'interp'}) rather than a flat colour, since ESP (unlike the density/MO lobes) has no single per-region colour to begin with. The decimated (\code{'MaxVertices'}) vertex set is unchanged by \code{'SurfaceStyle'} -- reduce \code{'MaxVertices'} for a more readable wireframe \\
\end{argtable}

Returns \code{h.surf}, the coloured surface patch handle.

\begin{tbnote}
\textbf{The analytic method is validated end-to-end against real
Gaussian ESP data; the numerical fallback is not publication-grade.} A
head-to-head comparison against a real Gaussian-generated ESP
\code{.cube} file (\code{cubegen}, Full Density Matrix option), sampled
at the \emph{actual} density-isosurface distance range used by this
function (2.4--4.4~Bohr from the nearest nucleus, the default
\code{IsoValue=0.001}): the default analytic method gave correlation
$1.0000000$ and an RMS error of $\sim\!0\%$ of the true $\sim\!0.14$
Hartree/e signal span; the numerical fallback, at its best-tested
setting, gave correlation $0.03$ (essentially none) and an RMS error
$54\%$ of that span. The numerical method's own internal aliasing issue
that drives this: a uniform grid systematically mis-estimates the raw
density right at each nuclear cusp (the same phenomenon documented for
individual MOs in \S\ref{sec:gutil-mo-surface-fn}) -- verified
separately, the total electron count captured by a naive grid sum
ranged from $-2\%$ to $+31\%$ depending on \code{'ESPGridSpacing'}
alone, \emph{non}-monotonically. The density grid is rescaled to the
exact electron count (\code{data.Nelec}) before the Coulomb sum, which
removes the dominant, roughly-uniform component of this bias but not
its remaining atom-to-atom spatial variation -- which is why, even with
this correction, the numerical method's correlation against real
Gaussian ESP data is essentially zero at realistic surface distances.
\textbf{Use the numerical fallback (\code{'ESPMethod','numeric'}, or the
automatic fallback for Cartesian-D/F/G bases) for a qualitative
electrophilic/nucleophilic pattern only, never for quantitative ESP
values.}
\end{tbnote}

\begin{lstlisting}
data = G16_fchk_read('4-NTP.fchk');
G_draw_esp_surface(data);                       % exact, analytic (default)
G_draw_esp_surface(data, 'MaxVertices', 1500);   % faster, same accuracy

data2 = G16_fchk_read('field_on.fchk');
G_draw_esp_surface(data, 'CompareTo', data2);    % ESP difference map

G_draw_esp_surface(data, 'SaveCube', 'esp.cube', 'CubeSpacing', 0.5);
G_draw_esp_surface(data, 'SurfaceStyle', 'grid', 'MaxVertices', 800);
\end{lstlisting}

See also \code{G\_draw\_density\_surface} (\S\ref{sec:gutil-density-surface}),
\code{G\_draw\_mo\_surface} (\S\ref{sec:gutil-mo-surface-fn}),
\code{G09\_fchk\_read}\slash\code{G16\_fchk\_read}.

% -----------------------------------------------------------------------------
\subsection{G\_draw\_cube\_surface}
\label{sec:gutil-cube-surface}

\begin{lstlisting}
h = G_draw_cube_surface(cubefile)
h = G_draw_cube_surface(cubefile, 'Name', Value, ...)
h = G_draw_cube_surface(cubefile, 'ColorBy', esp_cubefile, ...)
\end{lstlisting}

Renders an isosurface straight from a saved Gaussian-format \code{.cube}
file -- the load-and-replot companion to the other three functions' own
\code{'SaveCube'} option, and equally happy with a real \code{cubegen}/
GaussView-generated \code{.cube}. No \code{.fchk} is read and no
basis-set data is re-evaluated: the scalar field and the atoms are both
read directly from \code{cubefile} by \code{g\_read\_cube\_file}, and
\textsc{Matlab}'s own \code{isosurface} is applied to it exactly as the
other three functions apply it to a freshly-evaluated grid.

\begin{argtable}{Option}{Default}{Description}{2.8cm}{2.6cm}{X}
\code{'ColorBy'}      & \code{''}                  & Path to a SECOND \code{.cube} file. If given, the isosurface SHAPE still comes from \code{cubefile} (a single positive-value envelope, as for a density cube), but each vertex is coloured by trilinearly interpolating the SECOND cube's field there -- reproducing GaussView's classic ``ESP mapped on electron density'' picture from two independently-saved cube files (a density cube as \code{cubefile}, an ESP cube as \code{'ColorBy'}), with no \code{.fchk} needed at all. The two cubes must share an IDENTICAL grid (same origin, spacing, and point counts); a clear error is raised otherwise, rather than silently resampling -- see the tbnote below \\
\code{'IsoValue'}     & guessed                    & Guessed from \code{cubefile}'s title text: \code{0.001} for a density-like field (title contains ``density''), \code{0.02} for an MO-amplitude cube (negative-\code{Natoms} header, GaussView's own MO default) -- matching \code{G\_draw\_density\_surface}/\code{G\_draw\_mo\_surface}'s own defaults for a cube saved by this toolbox. For anything else (an ESP cube, or an unrecognized/external cube's title) there is no universally sensible default -- \code{'IsoValue'} must then be given explicitly, or a clear error explains why \\
\code{'SurfaceStyle'} & \code{'transparent'}       & \code{'solid'} \textbar\ \code{'transparent'} \textbar\ \code{'grid'} -- as in the other three functions. With \code{'ColorBy'}, \code{'grid'}'s edges are coloured by interpolating the second cube's field (\code{EdgeColor,'interp'}), like \code{G\_draw\_esp\_surface}; without it, by \code{'PosColor'}/\code{'NegColor'} like the others \\
\code{'FaceAlpha'}    & \code{0.55}                & As in \code{G\_draw\_density\_surface}/\code{G\_draw\_mo\_surface}: used as-is with \code{'transparent'}, overridden to \code{1.0} with \code{'solid'} unless given explicitly, unused with \code{'grid'} \\
\code{'PosColor'} / \code{'NegColor'} & blue / red & Lobe colours (no \code{'ColorBy'}), or the diverging-colormap endpoints (with \code{'ColorBy'}) \\
\code{'MaxVertices'}  & \code{5000}                & The isosurface mesh is decimated (\code{reducepatch}) to at most this many vertices, mainly to bound \code{'ColorBy'} interpolation cost \\
\code{'ShowMolecule'} & \code{true}                & Overlays a CPK model built directly from the atoms stored IN the cube file -- no \code{.fchk} needed \\
\code{'ShowLabels'} / \code{'AtomScale'} / \code{'Ax'} / \code{'Verbose'} & as elsewhere & Same meaning as in \code{G\_draw\_mo\_surface} \\
\end{argtable}

Returns \code{h.pos}/\code{h.neg} (positive/negative lobe patch handles;
\code{h.neg} is \code{[]} for an unsigned field, or whenever
\code{'ColorBy'} is given, since that mode always renders a single
shape-defining lobe).

\begin{tbnote}
\textbf{Signed-ness and the \code{'ColorBy'} grid requirement.} A field
is treated as SIGNED (both a $+$\code{IsoValue} and a
$-$\code{IsoValue} lobe, coloured \code{'PosColor'}/\code{'NegColor'})
if it is an MO cube, or if its values genuinely span both signs beyond
numerical noise (an ESP cube loaded alone, for instance, IS signed by
this rule); otherwise only the single $+$\code{IsoValue} lobe is drawn.
\code{'ColorBy'} overrides this entirely -- the shape is always the
single envelope of \code{cubefile}. Because \code{G\_draw\_density\_surface}/
\code{G\_draw\_mo\_surface} auto-expand \code{'Padding'} when the
isosurface would otherwise be clipped (\S\ref{sec:gutil-mo-surface-fn}),
two cubes saved with the same \emph{nominal} \code{'GridSpacing'}/
\code{'Padding'} can still end up on different grids if only one of them
triggered that auto-expansion -- verified directly: this is exactly what
happened on a first attempt at combining a density cube with an ESP
cube. Pass an \emph{explicit} \code{'Padding'} (which disables
auto-expansion) matching the ESP cube's \code{'CubePadding'} exactly to
guarantee identical grids.
\end{tbnote}

\begin{tbnote}
\textbf{\code{'ColorBy'}'s colour scale matches \code{G\_draw\_esp\_surface}'s
own convention, not a naive max.} The colour axis is set from the 99th
percentile of $|\text{colour field}|$ (\code{g\_pctile\_local}, the same
helper \code{G\_draw\_esp\_surface} itself uses), \emph{not} the raw
maximum -- verified to matter directly: an earlier version used
\code{max(abs(\ldots))}, and a single outlier vertex was enough to wash
out the entire colour scale compared to \code{G\_draw\_esp\_surface}'s
own rendering of the same data. With the percentile-based fix, the two
functions' colour axis limits on the same molecule agree to within
$\sim$0.3\% (the small remainder being ordinary trilinear-interpolation
error from reading the ESP field off a separately-saved cube grid,
rather than evaluating it exactly at each surface vertex the way
\code{G\_draw\_esp\_surface} does internally).
\end{tbnote}

\begin{lstlisting}
% Re-plot a density cube saved earlier -- no .fchk needed:
G_draw_cube_surface('density.cube');

% Re-plot a saved MO cube (sign auto-detected):
G_draw_cube_surface('homo.cube', 'SurfaceStyle', 'grid');

% Classic "ESP mapped on density" from two independently saved cubes,
% entirely without re-reading the .fchk:
G_draw_cube_surface('density.cube', 'ColorBy', 'esp.cube', 'IsoValue', 0.001);
\end{lstlisting}

See also \code{G\_draw\_density\_surface} (\S\ref{sec:gutil-density-surface}),
\code{G\_draw\_mo\_surface} (\S\ref{sec:gutil-mo-surface-fn}),
\code{G\_draw\_esp\_surface} (\S\ref{sec:gutil-esp-surface}).

% -----------------------------------------------------------------------------
\subsection{Python port (\code{MOSurface/python/})}
\label{sec:gutil-mo-surface-python}

All four functions above are also ported to Python, as a standalone
\code{mosurface} package in \code{MOSurface/python/} -- deliberately kept
as a separate companion rather than folded into the \code{G16parser}
package itself, so that this real-space isosurface engine (with its own
heavier dependencies, \code{scipy} and \code{scikit-image}, on top of
\code{G16parser}'s own \code{numpy}/\code{matplotlib}) does not become a
hard dependency of the core Python toolbox. Same naming philosophy as
\code{G16parser}: public functions are prefixed \code{g16\_}
(\code{g16\_draw\_mo\_surface}, \code{g16\_draw\_density\_surface},
\code{g16\_draw\_esp\_surface}, \code{g16\_draw\_cube\_surface}), every
other (leading-underscore) name is an internal helper. Options are the
same as their MATLAB counterparts, in \code{snake\_case}
(\code{'IsoValue'} $\to$ \code{isovalue}, \code{'ShowMolecule'} $\to$
\code{show\_molecule}, \ldots). \code{data} is the \code{Struct}
returned by \code{G16parser}'s own \code{g16\_fchk\_read} -- the one
real dependency on \code{G16parser} (for the input struct's shape and
for the \code{g16\_draw\_molecule} CPK overlay), mirroring
\code{MOSurface/}'s own dependency on \code{G09\_fchk\_read}/
\code{G16\_fchk\_read} and \code{G09\_draw\_molecule}/
\code{G16\_draw\_molecule}.

\begin{tbnote}
\textbf{Validation.} The numerically critical parts were cross-checked
directly against the MATLAB original on a real \code{.fchk}
(4-NTP, 218 basis functions), not just internally self-consistent:
\code{eval\_mo\_on\_grid}/\code{eval\_density\_on\_grid} (every S/P/SP/
D/F/G shell, pure and Cartesian) agree with
\code{g\_eval\_mo\_on\_grid.m}/\code{g\_eval\_density\_on\_grid.m} to
$\sim$1e-11 relative difference (HOMO/LUMO/HOMO-2/the highest-numbered
MO, and the total density); the analytic McMurchie-Davidson ESP
evaluator agrees with \code{g\_eval\_esp\_analytic.m} to $\sim$4e-13
relative difference; pure D/F/G self-normalization (synthetic
single-shell self-overlap, fine-grid Riemann sum) converges to
$1.0000$ for every component of every supported shell; \code{.cube}
read/write round-trips exactly (to \code{\%.5E} text precision) against
a brute-force nested-loop reference. All four functions were then run
end-to-end on real \code{.fchk} files (closed-shell; UKS spin density;
a field-on/field-off pair for density/ESP difference), producing
chemically sensible pictures -- e.g.\ the expected SOMO lobe
perpendicular to the molecular plane for a CH$_3$ radical's spin
density, and the expected heteroatom-localised electron-rich region in
ESP-on-density colouring.
\end{tbnote}

\begin{tbnote}
\textbf{Rendering differences from MATLAB (intentional, not a fidelity
concern).} Isosurface extraction uses
\code{skimage.measure.marching\_cubes} in place of \code{isosurface}/
\code{isonormals} (already returns per-vertex normals from the volume
gradient). Mesh decimation (\code{max\_vertices}) uses a simple
vertex-clustering simplification instead of \code{reducepatch} -- a
different, simpler algorithm, fine for a smooth ESP/\code{color\_by}
field. matplotlib's \code{Poly3DCollection} has no true per-vertex
Gouraud shading like \code{PATCH}: flat-coloured lobes use
\code{shade=True} (simple per-face directional shading), and ESP/
\code{color\_by} surfaces use a per-\emph{face} colour (averaged from
its 3 vertices) in place of \code{FaceColor,'interp'}.
\end{tbnote}

\begin{tbnote}
\textbf{Performance characteristics (read before using on a large
system).} Pure-Python/NumPy evaluation is markedly slower than
MATLAB's compiled engine in two specific ways. (1)~The default
auto-expanding \code{padding} re-evaluates the \emph{entire} grid on
every retry (up to 4 attempts, $\times1.6$ each); on a real
218-basis-function molecule, one retry sequence took several minutes
where an explicit, sufficient \code{padding}
(\code{padding\_explicit=True}) took $\sim$4~s -- pass an explicit
\code{padding} for routine use rather than relying on the auto-expand
default. (2)~\code{g16\_draw\_esp\_surface}'s analytic ESP evaluation
has a large \emph{fixed} per-call cost dominated by the number of shell
\emph{pairs} ($\sim$11~s fixed overhead measured for a 104-subshell
basis, largely independent of point count) -- fine for the normal case
(a few thousand \emph{decimated surface vertices}, via
\code{max\_vertices}), but this makes \code{save\_cube} with the
analytic method over a \emph{full volumetric grid} (potentially
hundreds of thousands of points) impractical at realistic resolution,
unlike MATLAB, where the same feature is merely slow (per
\code{G\_draw\_esp\_surface.m}'s own ``genuinely slow for a large
basis'' docstring note) rather than impractical. For \code{save\_cube}
on \code{g16\_draw\_esp\_surface}, use a deliberately coarse
\code{cube\_spacing}/small \code{cube\_padding}, or
\code{esp\_method='numeric'} (faster, less accurate -- see
G\_DRAW\_ESP\_SURFACE's own ``Accuracy'' notes above, which apply
unchanged here) if a full-resolution ESP cube is genuinely needed.
Vectorizing the analytic evaluator's inner shell-pair loop would
address this but was out of scope for this initial port, which
prioritized exact numerical fidelity over performance.
\end{tbnote}

\begin{lstlisting}
import G16parser as g16
from mosurface import g16_draw_mo_surface, g16_draw_esp_surface

data = g16.g16_fchk_read('4-NTP.fchk')
g16_draw_mo_surface(data, 'HOMO')
g16_draw_esp_surface(data, padding=4.0, padding_explicit=True, max_vertices=1500)
\end{lstlisting}

\newpage

\section{RamanBrowser: interactive Raman/IR stick-spectrum browser}
\label{sec:gutil-raman-browser}
% =============================================================================
\code{RamanBrowser/} provides a single click-to-select companion to the
core toolbox's own \code{G16\_modeViewer}/\code{G09\_modeViewer} (text
drop-down mode selectors): \code{G\_raman\_browser} plots Raman activity
(or IR intensity) against frequency as a stick spectrum and lets the
mode be picked by clicking directly on its stick -- often the more
natural way to pick out ``that strong peak at $\sim$1580~cm$^{-1}$''
than first reading its value off a list. Unlike the two
\code{modeViewer}s, which are each tied to one Gaussian major version,
\code{G\_raman\_browser} is a single function that works with \emph{both}:
it auto-detects Gaussian 09 vs.\ Gaussian 16 (\S\ref{sec:gutil-raman-browser-fn})
and dispatches internally to the matching
\code{G09\_}/\code{G16\_structure}/\code{nmodes}/\code{draw\_mode}
functions, unchanged, with no core-toolbox modification. A Python port
also exists (\S\ref{sec:gutil-raman-browser-python}), G16-only.

% -----------------------------------------------------------------------------
\subsection{G\_raman\_browser}
\label{sec:gutil-raman-browser-fn}

\begin{lstlisting}
G_raman_browser(filename)
G_raman_browser(filename, 'Name', Value, ...)
\end{lstlisting}

Opens a window with the requested stick spectrum. Clicking anywhere in
the axes selects the mode whose frequency is nearest the click (a red
marker highlights it, and an info line above the plot reports its mode
number, frequency, symmetry, and intensity). A ``Draw mode'' button then
calls \code{draw\_mode} for the selected mode, opening (or replacing) a
3D structure figure with its displacement arrows.

For \code{'Property','Raman'} on a file with \code{CPHF=RdFreq} data
(\code{nm.has\_RamanFr}, \S\ref{sec:gutil-raman-browser-fn}), the info
line reports \emph{both} the static and the dynamic Raman activity
instead of a single value, e.g.\ \code{Selected: mode 1, 17.4 cm-1 (A),
Raman static = 1.74, 785nm = 1.89} -- one \code{, Nnm = ...} term per
incident wavelength beyond the static one, converted from
\code{nm.IncidentLight}'s cm$^{-1}$ values. Without \code{CPHF=RdFreq}
data, or for \code{'Property','IR'}, the line keeps its original
single-value form (e.g.\ \code{ir = 4.28}).

\begin{argtable}{Option}{Default}{Description}{2.8cm}{2.6cm}{X}
\code{'Property'}   & \code{'Raman'}  & \code{'Raman'} \textbar\ \code{'IR'} -- which stick intensity to plot. \code{'Raman'} requires the file to actually contain Raman activities (a \code{Freq=Raman} job); a clear error suggests \code{'IR'} instead otherwise \\
\code{'Scale'}      & \code{1.5}      & Forwarded to \code{draw\_mode} \\
\code{'AtomScale'}  & \code{0.35}     & Forwarded to \code{draw\_mode} \\
\code{'ShowLabels'} & \code{false}    & Forwarded to \code{draw\_mode} \\
\end{argtable}

\begin{tbnote}
\textbf{Gaussian 09/16 auto-detection.} The Gaussian major version is
read from the ``Gaussian~NN, Revision~X.YY,'' citation line near the top
of \code{filename} via \code{G16\_gaussian\_version} (which, despite its
name, reads this line for any \code{NN}); \code{G09\_structure}/
\code{G09\_nmodes}/\code{G09\_draw\_mode} are used for a detected
Gaussian~09 file, \code{G16\_structure}/\code{G16\_nmodes}/
\code{G16\_draw\_mode} otherwise (including when the version cannot be
determined at all, matching this function's original, still most
common, use case). This is safe because the two triples' output
structs/call signatures are field-for-field and argument-for-argument
identical -- verified directly on a real Gaussian~09 file
(\code{CH4\_NBO.LOG}) and a real Gaussian~16 file
(\code{4-NTP\_nbo.out}) side by side, including a simulated stick click
on the Gaussian~09 file correctly reaching \code{G09\_draw\_mode}.
\end{tbnote}

\begin{tbnote}
\textbf{Static vs.\ dynamic Raman in the info line.} This relies on
\code{G09\_nmodes}/\code{G16\_nmodes}'s own \code{.RamanFr}/
\code{.IncidentLight}/\code{.has\_RamanFr} fields (\S\ref{sec:nmodes-ref})
-- \code{G\_raman\_browser} adds no
parsing of its own, it only formats what those functions already
expose. Verified directly on a real 785~nm job
(\code{MPBA\_SH.out}): a simulated click on mode~1 produced
\code{Raman static = 1.74, 785nm = 1.89}, matching \code{nm.Raman(1)}
and \code{nm.RamanFr(1,2)} exactly; the same click on a file with no
\code{CPHF=RdFreq} data produced the plain, unchanged
\code{raman = 0.87} form, with no spurious \code{Nnm = ...} term.
\end{tbnote}

\begin{tbnote}
\textbf{Window placement on multi-monitor systems.} A fixed pixel
\code{Position} for the browser window can land off-screen entirely on a
multi-monitor setup whose primary display does not start at virtual
desktop coordinate $[0,0]$ -- the window is technically created but never
visibly appears (confirmed directly: this happened on first use, and
restarting \textsc{Matlab} was the only workaround until fixed). The same
fix already used by \code{G16\_modeViewer} is applied here: a short-lived
invisible ordinary \code{figure} is queried for its default placement to
identify the active monitor from \code{MonitorPositions}, and the browser
window is then created directly on that monitor in a single atomic call.
\end{tbnote}

\begin{lstlisting}
% Browse the Raman spectrum, click a peak, press "Draw mode" (G16 file):
G_raman_browser('4-NTP.out');

% Browse IR intensities instead, with a larger displacement-arrow scale:
G_raman_browser('4-NTP.out', 'Property', 'IR', 'Scale', 2);

% Works identically on a Gaussian 09 file, auto-detected:
G_raman_browser('indaco.log');
\end{lstlisting}

See also \code{G16\_modeViewer}, \code{G16\_draw\_mode},
\code{G16\_nmodes}, \code{G16\_structure}, \code{G09\_modeViewer},
\code{G09\_draw\_mode}, \code{G09\_nmodes}, \code{G09\_structure},
\code{G16\_gaussian\_version}.

% -----------------------------------------------------------------------------
\subsection{Python port (\code{RamanBrowser/python/})}
\label{sec:gutil-raman-browser-python}

\code{G\_raman\_browser} is also ported to Python, as a standalone
\code{ramanbrowser} package in \code{RamanBrowser/python/} -- same
separate-companion-package status as \code{MOSurface/python/}'s own
\code{mosurface} package (kept out of \code{G16parser} itself, see
\S\ref{sec:gutil-mo-surface-python}), and the
same naming philosophy: one public function,
\code{g16\_raman\_browser}, everything else a leading-underscore
internal helper. \code{tkinter} (with an embedded matplotlib
\code{FigureCanvasTkAgg}) replaces \code{uifigure}/\code{uiaxes} for the
click-to-select stick spectrum window, and \code{g16\_draw\_mode}'s own
matplotlib window (as already used by \code{G16parser}'s
\code{g16\_mode\_viewer}) is opened by the ``Draw mode'' button, exactly
mirroring the MATLAB original's behaviour.

\begin{lstlisting}
from ramanbrowser import g16_raman_browser

g16_raman_browser('4-NTP.out')
g16_raman_browser('4-NTP.out', property='ir', scale=2)
\end{lstlisting}

\begin{tbnote}
\textbf{G16-only (a real scope difference from the MATLAB original).}
\code{G16parser} only targets Gaussian~16 output (a deliberate,
documented scope --- \S\ref{sec:gutil-mo-surface-python}'s own note on
why no \code{G09parser} exists), so this port cannot dispatch to a
Gaussian~09 parser the way \code{G\_raman\_browser.m} does via
\code{G09\_structure}/\code{G09\_nmodes}/\code{G09\_draw\_mode} --
there is nothing on the Python side to dispatch \emph{to}. Passing a
Gaussian~09 file is unsupported (not silently wrong -- \code{g16\_nmodes}
still expects Gaussian~16 formatting).
\end{tbnote}

\begin{tbnote}
\textbf{Static vs.\ dynamic Raman (\code{.RamanFr}), forward-compatible.}
The info line's static/dynamic Raman split (\S\ref{sec:gutil-raman-browser-fn})
is read via \code{getattr(nm, 'has\_RamanFr', False)}: since
\code{G16parser}'s own \code{g16\_nmodes} does not yet expose
\code{.RamanFr}/\code{.IncidentLight}/\code{.has\_RamanFr} (a separate,
already-known port gap), this currently always falls back to the plain
single-value form -- correctly, with no error or spurious output -- and
will pick up the enhanced display automatically, with no changes needed
here, once/if that field is ported to \code{G16parser}.
\end{tbnote}

\begin{tbnote}
\textbf{Validation.} Verified end-to-end with a simulated click
(a real \code{matplotlib.backend\_bases.MouseEvent}, not just a direct
function call) on \code{4-NTP\_nbo.out}: selecting the same stick
produced the identical result as the MATLAB original on the same file
--- \code{Selected: mode 10, 528.5 cm-1 (A), raman = 0.25} --- and the
``Draw mode'' button correctly invoked \code{g16\_draw\_mode} with mode
index 10.
\end{tbnote}

\newpage

% =============================================================================
\section{Orbital energy analysis}
% =============================================================================
\subsection{G09\_orbital\_energies / G16\_orbital\_energies}

\begin{lstlisting}
oe = G09_orbital_energies(filename)
oe = G09_orbital_energies(filename, 'step', 'last')
oe = G09_orbital_energies(filename, 'step', N)
\end{lstlisting}

Extracts molecular orbital energies (HOMO/LUMO and the full
occupied/virtual spectrum) from a Gaussian 09/16 output file, by parsing the
\code{"Alpha  occ. eigenvalues --"} / \code{"Alpha virt. eigenvalues --"}
lines printed by Gaussian's population analysis (and the corresponding
\code{"Beta"} lines for open-shell calculations). These blocks repeat once
per SCF calculation in the file (e.g.\ once per optimisation step); the
\code{'step'} option selects which block to report, consistently with
\code{G09\_energy}/\code{G09\_structure}.

\begin{argtable}{Option}{Default}{Description}{2.6cm}{2.4cm}{X}
\code{'step'} & \code{'last'} & Which eigenvalue block to use: \code{'first'}, \code{'last'}, or an integer index \\
\end{argtable}

\textbf{Output struct \code{oe}} (energies in Hartree unless noted):

\begin{argtable}{Field}{Type}{Description}{2.8cm}{2.2cm}{X}
\code{.alpha\_occ}  & vector & Occupied alpha orbital energies \\
\code{.alpha\_virt} & vector & Virtual alpha orbital energies \\
\code{.beta\_occ}   & vector & Occupied beta orbital energies (\code{[]} if closed-shell) \\
\code{.beta\_virt}  & vector & Virtual beta orbital energies (\code{[]} if closed-shell) \\
\code{.has\_beta}   & logical & \code{true} for open-shell (UHF/UKS) calculations \\
\code{.HOMO} / \code{.LUMO} & double & Highest occupied / lowest unoccupied orbital energy \\
\code{.gap}         & double & \code{LUMO -- HOMO} \\
\code{.HOMO\_alpha} / \code{.LUMO\_alpha} & double & Alpha HOMO/LUMO (\code{== .HOMO}/\code{.LUMO} for closed-shell) \\
\code{.HOMO\_beta} / \code{.LUMO\_beta} & double & Beta HOMO/LUMO (\code{[]} if closed-shell) \\
\code{.HOMO\_eV} / \code{.LUMO\_eV} / \code{.gap\_eV} & double & Same three quantities, in eV \\
\code{.step}    & int  & Block index actually used \\
\code{.Nsteps}  & int  & Number of eigenvalue blocks found in the file \\
\code{.filename} & char & Source file name \\
\end{argtable}

\begin{lstlisting}
oe = G09_orbital_energies('V_E00t.out');
fprintf('HOMO-LUMO gap = %.3f eV\n', oe.gap_eV);
\end{lstlisting}

\begin{tbnote}
For open-shell (UHF/UKS) calculations, \code{HOMO} and \code{LUMO} are
computed as $\max(\text{HOMO}_\alpha,\text{HOMO}_\beta)$ and
$\min(\text{LUMO}_\alpha,\text{LUMO}_\beta)$ respectively, so \code{.gap}
always refers to the smallest alpha/beta HOMO--LUMO gap. No differences
between the G09 and G16 versions beyond the parsed file's origin.
\end{tbnote}

% -----------------------------------------------------------------------------
\subsection{G09\_draw\_orbital / G16\_draw\_orbital}

\begin{lstlisting}
G09_draw_orbital(oe)
G09_draw_orbital(oe, 'Name', Value, ...)
\end{lstlisting}

Draws a molecular-orbital energy-level diagram from the struct returned by
\code{G09\_orbital\_energies}/\code{G16\_orbital\_energies}, showing a
user-selected number of occupied/virtual levels around the frontier
orbitals and highlighting the HOMO--LUMO transition with an arrow labelled
with the gap value.

\begin{argtable}{Option}{Default}{Description}{2.8cm}{3.0cm}{X}
\code{'NLevels'}     & \code{5}       & Number of occupied/virtual levels shown on each side of the frontier orbitals \\
\code{'Units'}       & \code{'eV'}    & Energy units: \code{'eV'} or \code{'Hartree'} \\
\code{'Ax'}          & (new figure)   & Draw into an existing axes handle \\
\code{'OccColor'}    & \code{[0.25\ 0.25\ 0.70]} & Colour for non-frontier occupied levels \\
\code{'VirtColor'}   & \code{[0.70\ 0.30\ 0.25]} & Colour for non-frontier virtual levels \\
\code{'HOMOColor'}   & \code{[0.00\ 0.45\ 0.85]} & Highlight colour for the HOMO \\
\code{'LUMOColor'}   & \code{[0.90\ 0.35\ 0.00]} & Highlight colour for the LUMO \\
\code{'ArrowColor'}  & \code{[0.15\ 0.60\ 0.15]} & HOMO$\to$LUMO transition arrow colour \\
\code{'ShowLabels'}  & \code{true}    & Annotate each level with its energy value \\
\code{'ShowSpins'}   & \code{true}    & Draw a paired-electron arrow glyph ($\uparrow\downarrow$) on occupied levels \\
\code{'FontSize'}    & \code{8}       & Font size for level energy-value labels \\
\code{'FrontierFontSize'} & \code{13} & Font size for the HOMO/LUMO tags and the gap annotation \\
\code{'ArrowHeadSize'}    & \code{0.2} & HOMO--LUMO arrow head size, as a fraction of the arrow length (passed to \texttt{quiver}) \\
\code{'Title'}       & filename or \code{'Orbital Energy Diagram'} & Figure title \\
\end{argtable}

\begin{lstlisting}
oe = G09_orbital_energies('a1.out');
G09_draw_orbital(oe)
G09_draw_orbital(oe, 'NLevels', 8, 'Units', 'Hartree')
\end{lstlisting}

\begin{tbnote}
Requires the \code{oe} struct returned by \code{G09\_orbital\_energies}/
\code{G16\_orbital\_energies} and errors if \code{alpha\_occ}/\code{alpha\_virt}
are missing or empty. No differences between the G09 and G16 versions
beyond the input struct's origin.
\end{tbnote}

\newpage

% =============================================================================
\section{Response properties (G16 only)}
% =============================================================================
\subsection{G16\_hyperpolar}

\begin{lstlisting}
hp = G16_hyperpolar(filename)
hp = G16_hyperpolar(filename, 'units', 'au')    % default
hp = G16_hyperpolar(filename, 'units', 'esu')   % 10^-30 esu
hp = G16_hyperpolar(filename, 'units', 'SI')    % 10^-50 C^3 m^3 J^-2
\end{lstlisting}

Extracts the dipole hyperpolarisability $\beta$ (static, dynamic, and
vibrational contributions) from a Gaussian~16 output file.

\begin{argtable}{Option}{Default}{Description}{2.6cm}{2.4cm}{X}
\code{'units'} & \code{'au'} & \code{'au'} \textbar\ \code{'esu'} ($10^{-30}$ esu) \textbar\ \code{'SI'} ($10^{-50}$ C$^3$m$^3$J$^{-2}$) \\
\end{argtable}

\begin{argtabletwo}{Field}{Description}{3.2cm}{X}
\code{.beta0.par\_z}, \code{.beta0.perp\_z} & Parallel / perpendicular component of static $\beta(0;0,0)$ along $z$ \\
\code{.beta0.vec\_x/y/z}, \code{.beta0.beta\_vec} & Vector components and magnitude $|\beta_\text{vec}|$ \\
\code{.beta0.tensor} & All Cartesian tensor components ($xxx,xxy,\ldots,zzz$) \\
\code{.beta\_dyn} & Struct array, one entry per laser frequency: \code{.lambda\_nm}, \code{.par\_z}, \code{.perp\_z}, \code{.vec\_x/y/z}, \code{.beta\_vec}, \code{.tensor} \\
\code{.N\_dyn} & Number of dynamic (laser) frequencies found \\
\code{.beta\_vib}, \code{.has\_vib} & Diagonal vibrational hyperpolarisability, logical \\
\code{.filename} & \\
\end{argtabletwo}

\begin{lstlisting}
hp = G16_hyperpolar('V_E00t.out');
hp.beta0.beta_vec          % static beta modulus, a.u.
hp.beta_dyn(1).lambda_nm   % wavelength of the first dynamic entry (nm)
hp.beta_dyn(1).beta_vec    % |beta| at that frequency
\end{lstlisting}

% -----------------------------------------------------------------------------
\subsection{G16\_tddft}

\begin{lstlisting}
td = G16_tddft(filename)
td = G16_tddft(filename, 'nstates', N)     % first N states only
td = G16_tddft(filename, 'plot', true)     % also generate UV-Vis plot
td = G16_tddft(filename, 'FWHM_eV', 0.3)   % broadening FWHM (eV)
\end{lstlisting}

Extracts TD-DFT excited states and generates a Gaussian-broadened UV-Vis
spectrum.

\begin{argtable}{Option}{Default}{Description}{2.6cm}{2.4cm}{X}
\code{'nstates'} & \code{Inf}  & Load only the first $N$ excited states \\
\code{'plot'}    & \code{false}& Generate the UV-Vis plot \\
\code{'FWHM\_eV'}& \code{0.30} & Gaussian broadening FWHM (eV) \\
\end{argtable}

\begin{argtabletwo}{Field}{Description}{3.2cm}{X}
\code{.n}, \code{.mult} & State index; multiplicity label (e.g.\ \code{'Singlet-A'}, \code{'Triplet-B'}) \\
\code{.eV}, \code{.nm}, \code{.f} & Excitation energy (eV), wavelength (nm), oscillator strength \\
\code{.S2}, \code{.has\_S2} & $\langle S^2\rangle$ expectation value (0 for pure singlets), logical \\
\code{.trans} & Cell array, one $[N_\text{contrib}\times3]$ matrix per state: columns \code{[MO\_from, MO\_to, coeff]} (\code{MO\_to}$<0$ denotes a de-excitation) \\
\code{.Nstates} & Total number of excited states loaded \\
\code{.x\_nm}, \code{.eps\_cont} & Wavelength grid (100--1000~nm), Gaussian-broadened UV-Vis spectrum (arb.\ units) \\
\code{.FWHM\_eV}, \code{.filename} & \\
\end{argtabletwo}

\begin{lstlisting}
td = G16_tddft('violacein_td.out');
plot(td.x_nm, td.eps_cont)     % UV-Vis spectrum
stem(td.nm, td.f)              % transition stem plot
td.trans{2}                    % transitions contributing to the 2nd state
\end{lstlisting}

\begin{tbnote}
No G09 equivalent is provided: TD-DFT excited-state analysis and
hyperpolarisability response properties in this toolbox are currently
supported only for Gaussian~16 output files.
\end{tbnote}

\newpage

% =============================================================================
\section{Structural analysis utilities}
% =============================================================================
\subsection{G09\_get\_bond\_length / G16\_get\_bond\_length}

\begin{lstlisting}
bondTable = G09_get_bond_length(mol)
bondTable = G09_get_bond_length(mol, 'Name', Value, ...)
\end{lstlisting}

Builds the bond-length table of a molecule from a \code{mol} struct (fields
\code{.symbols}, \code{.xyz}), detecting bonds via a covalent-radius
criterion, and returns a MATLAB \code{table}.

\begin{argtable}{Option}{Default}{Description}{2.6cm}{2.4cm}{X}
\code{'Tolerance'} & \code{1.15} & Multiplicative factor on the sum of covalent radii for the bonding criterion \\
\code{'IncludeH'}  & \code{true} & Include bonds involving hydrogen atoms \\
\code{'SortBy'}    & \code{'distance'} & Sort by \code{'distance'} (ascending) or \code{'atom'} (atom index) \\
\code{'SaveAs'}    & \code{''}  & Path (\code{.xlsx} or \code{.csv}) to save the table to; not saved if omitted \\
\end{argtable}

Output columns: \code{Atom1}, \code{Sym1}, \code{Atom2}, \code{Sym2},
\code{Distance\_Ang}.

\begin{lstlisting}
T = G09_get_bond_length(mol, 'SortBy', 'atom');
T = G16_get_bond_length(mol, 'IncludeH', false, 'SaveAs', 'bonds.xlsx');
\end{lstlisting}

\begin{tbnote}
The bonding criterion is purely geometric (covalent radii), not derived
from Gaussian's own connectivity/topology matrix --- suitable for quick
structural analysis, not as the "official" G09/G16 connectivity.
\end{tbnote}

% -----------------------------------------------------------------------------
\subsection{G09\_nbo\_bonds / G16\_nbo\_bonds}
\label{sec:nbo-bonds-ref}

\begin{lstlisting}
bt = G09_nbo_bonds(filename)
bt = G09_nbo_bonds(filename, 'Name', Value, ...)
\end{lstlisting}

Determines bond order (single/double/triple) from an actual Gaussian NBO
(Natural Bond Orbital) analysis, as opposed to \code{get\_bond\_length}'s
geometric covalent-radius criterion or \code{draw\_molecule}'s bond-length
heuristic. Reads the "Natural Bond Orbitals (Summary)" table and counts,
for every atom pair, how many bonding (\code{BD}) natural bond orbitals
NBO assigned to it: one \code{BD} orbital means a single bond (\(\sigma\)
only), two mean a double bond (\(\sigma+\pi\)), three mean a triple bond
(\(\sigma+2\pi\)). Antibonding (\code{BD*}) orbitals are excluded.

\begin{argtable}{Option}{Default}{Description}{2.6cm}{2.4cm}{X}
\code{'section'} & \code{'last'} & Which "Natural Bond Orbitals (Summary)" block to read, for files with more than one (e.g.\ multi-step opt+freq+polar jobs that repeat the NBO analysis at each step); also accepts \code{'first'} \\
\code{'Lines'}   & \code{\{\}}   & Pre-read cell array of file lines, to skip re-reading the file \\
\end{argtable}

Output columns: \code{Atom1}, \code{Sym1}, \code{Atom2}, \code{Sym2},
\code{BondOrder}, \code{Occupancy} (summed NBO occupancy, \(\approx 2\)
per \code{BD} orbital).

\begin{lstlisting}
bt = G09_nbo_bonds('CH4_NBO.LOG');
\end{lstlisting}

\begin{tbnote}
Requires the source file to have been computed with the \code{pop=nbo}
Gaussian keyword (or similar, e.g.\ \code{pop=(nbo,savenbo)}); without it,
the output file simply does not contain NBO data and the function raises
an error --- bond order cannot be recovered after the fact from
geometry/energy alone. This reads the "Natural Bond Orbitals (Summary)"
table that \code{pop=nbo} always prints, not the separate "Wiberg bond
index matrix" (which additionally requires \code{pop=(nbo,bndidx)} and is
not parsed by this function).
\end{tbnote}

% -----------------------------------------------------------------------------
\subsection{G09\_gaussian\_version / G16\_gaussian\_version}

\begin{lstlisting}
gv = G09_gaussian_version(filename)
gv = G16_gaussian_version(filename)
\end{lstlisting}

Detects which Gaussian version/revision produced a
\code{.out}/\code{.log}/\code{.fchk} file.

\begin{tbnote}
For \code{.out}/\code{.log} files, reads the \code{"Gaussian NN, Revision
X.YY,"} citation line printed near the top of the file (works for any
\code{NN}: 09, 16, \ldots). \code{.fchk} files do not store this
information at all; in that case the function looks for a sibling
\code{.log}/\code{.out} file with the same base name in the same folder
and reads the version from there. If no sibling file is found,
\code{.major}/\code{.revision}/\code{.full} are returned empty and a
warning is issued.
\end{tbnote}

\begin{argtabletwo}{Field}{Description}{2.8cm}{X}
\code{.major}    & Gaussian major version, e.g.\ 9, 16 (\code{[]} if unknown) \\
\code{.revision} & Revision string, e.g.\ \code{'A.02'}, \code{'C.01'} (\code{''} if unknown) \\
\code{.full}     & Citation line as printed, e.g.\ \code{'Gaussian 09, Revision A.02'} \\
\code{.source}   & \code{'out/log'} \textbar\ \code{'fchk-sibling:<path>'} \textbar\ \code{'unknown'} \\
\code{.filename} & \\
\end{argtabletwo}

\begin{lstlisting}
gv = G09_gaussian_version('a1.out');        % gv.major = 16, gv.revision = 'C.01'
gv = G16_gaussian_version('molecule.fchk'); % falls back to molecule.out/.log if present
\end{lstlisting}

% -----------------------------------------------------------------------------
\subsection{G09\_restart / G16\_restart}

\begin{lstlisting}
gjf_file = G09_restart(filename)
gjf_file = G09_restart(filename, 'Name', Value, ...)

gjf_file = G16_restart(filename)
gjf_file = G16_restart(filename, 'Name', Value, ...)
\end{lstlisting}

Generates a Gaussian 09/16 input file (\code{.gjf}) to restart a
calculation from a geometry found in an existing \code{.log}/\code{.out}
file. \code{G16\_restart} reads geometry via \code{G16\_structure} (so it
follows that function's Standard/Input orientation handling); otherwise
the two functions are identical.

\begin{argtable}{Option}{Default}{Description}{2.6cm}{3.0cm}{X}
\code{'step'}        & \code{'last'} & \code{'last'} \textbar\ \code{'first'} \textbar\ integer $N$ --- which geometry step to restart from \\
\code{'output'}      & \code{<base>\_restarted.gjf} & Output \code{.gjf} path \\
\code{'extra\_route'}& \code{''} & Additional keywords appended to the route section \\
\code{'nproc'}       & \code{0}  & \code{\%nprocshared} override (\code{0} = keep/omit) \\
\code{'mem'}         & \code{''} & \code{\%mem} override (\code{''} = keep/omit) \\
\end{argtable}

\begin{tbnote}
If the source output file has no \code{\%chk} line (common in Gaussian~09
Windows output), the checkpoint name is derived from the source file name
instead. The generated \code{.gjf} is valid for restarting under either
Gaussian 09 or Gaussian 16, and can be read back with \code{G\_read\_input}.
\end{tbnote}

Returns the path of the generated \code{.gjf} file.

% -----------------------------------------------------------------------------
\subsection{G09\_available\_data / G16\_available\_data}
\label{sec:available-data-ref}

\begin{lstlisting}
T = G09_available_data(filename)
T = G16_available_data(filename)
\end{lstlisting}

Reads the route section (see \code{G09\_route}/\code{G16\_route}) and
checks it against the keyword requirements documented in
Section~\ref{sec:keyword-section-map} ("Which \code{.in} keywords
generate which \code{.out} section"), so you can tell up front which
data will actually be present in a file --- e.g.\ a TD-DFT-only single
point job (\code{td=(singlets,nstates=100)}, no \code{freq}) has no
vibrational normal modes, no IR/Raman spectra, and no thermochemistry
(ZPE/H/G/S), so calling \code{nmodes}/\code{spectra} on it raises an
error and \code{energy}'s \code{*\_corr}/\code{E0}/\code{U}/\code{H}/
\code{G}/\code{T}/\code{P}/\code{S} fields come back \code{NaN} ---
none of that is a bug, it is simply not in the file. This function lets
you check that in one call, rather than discovering it one error/NaN at
a time.

\begin{argtabletwo}{Field}{Description}{2.4cm}{X}
\code{.Function}  & Toolbox function/quantity name \\
\code{.Available} & \code{logical} --- true if the route contains the keyword(s) needed \\
\code{.Requires}  & Human-readable keyword requirement (e.g.\ \code{'freq'}, \code{'polar + cphf=rdfreq'}) \\
\code{.Notes}     & Short clarifying note \\
\end{argtabletwo}

\begin{lstlisting}
G16_available_data('V_E00_R_24_TD.out');
T = G16_available_data('molecule.out');
T(~T.Available, :)   % see what's NOT available in this file
\end{lstlisting}

\begin{tbnote}
This predicts availability from the route keywords alone; it does not
itself read the data sections, so an incomplete/crashed job can still
show a keyword as present without the corresponding output actually
having been printed (the extraction function will still raise its own
clear error in that case). \code{G09\_available\_data} omits the
\code{hyperpolar}/\code{tddft} rows entirely, since those functions do
not exist in the G09 toolbox.
\end{tbnote}

% -----------------------------------------------------------------------------
\subsection{G\_read\_input \normalfont{(shared, no G09\_/G16\_ prefix)}}

\begin{lstlisting}
ginp = G_read_input(filename)
\end{lstlisting}

Reads a Gaussian \emph{input} file (\code{.gjf}, \code{.com}, or
\code{.in}) and extracts the link0 commands, route section, title,
charge/multiplicity, and starting Cartesian geometry.

\begin{tbnote}
The Gaussian input file format does not differ between Gaussian~09 and
Gaussian~16, so this is the toolbox's only function without a
\code{G09\_}/\code{G16\_} prefix pair: an identical copy is kept in both
the \code{G09/} and \code{G16/} folders, so it is available regardless of
which toolbox you have on the MATLAB path. Only Cartesian-coordinate
geometry blocks are supported (not Z-matrix input).
\end{tbnote}

\begin{argtable}{Field}{Type}{Description}{2.8cm}{2.2cm}{X}
\code{.symbols}  & cell   & Atomic symbols \\
\code{.xyz}      & double & Starting coordinates in \AA\ (X Y Z) \\
\code{.Natoms}   & int    & Number of atoms \\
\code{.charge}   & int    & Total molecular charge \\
\code{.mult}     & int    & Spin multiplicity \\
\code{.title}    & char   & Title/comment line(s), joined with a space \\
\code{.route}    & char   & Full route section, single line \\
\code{.method}   & char   & Method parsed from the route (e.g.\ \code{'b3lyp'}), \code{''} if not found \\
\code{.basis}    & char   & Basis set parsed from the route (e.g.\ \code{'6-311+g(d,p)'}), \code{''} if not found \\
\code{.chk}      & char   & \code{\%chk} link0 value, \code{''} if absent \\
\code{.mem}      & char   & \code{\%mem} link0 value, \code{''} if absent \\
\code{.nproc}    & char   & \code{\%nprocshared}/\code{\%nproc} link0 value, \code{''} if absent \\
\code{.filename} & char   & Source file path \\
\end{argtable}

\begin{tbnote}
\code{ginp} has the same \code{.symbols}/\code{.xyz}/\code{.Natoms}/
\code{.filename} fields as the struct returned by
\code{G09\_structure}/\code{G16\_structure}, so it can be passed directly
to \code{G09\_draw\_molecule}/\code{G16\_draw\_molecule} or
\code{G09\_get\_bond\_length}/\code{G16\_get\_bond\_length} without any
conversion.
\end{tbnote}

\begin{lstlisting}
ginp = G_read_input('molecule.gjf');
G16_draw_molecule(ginp, 'ShowAxes', true);
fprintf('%s/%s, charge %d, mult %d\n', ginp.method, ginp.basis, ...
    ginp.charge, ginp.mult);
\end{lstlisting}

% -----------------------------------------------------------------------------
\subsection{G09\_list / G16\_list}

\begin{lstlisting}
G09_list()
T = G09_list()

G16_list()
T = G16_list()
\end{lstlisting}

Lists every \code{G09\_*.m}/\code{G16\_*.m} function found in the toolbox
folder, printing the function name and its H1 description line (sorted
alphabetically). Also returns the same information as a \code{table} with
columns \code{.Name}, \code{.Description}, \code{.File}.

\begin{lstlisting}
G09_list();
T = G16_list();
T(contains(T.Description, 'dipole', 'IgnoreCase', true), :)
\end{lstlisting}

\newpage

% =============================================================================
\section{One-call data extraction}
% =============================================================================
\subsection{G09\_read\_all / G16\_read\_all}

\begin{lstlisting}
T = G09_read_all(filename)
T = G16_read_all(filename)
\end{lstlisting}

Runs the full set of \code{G09\_XXX}/\code{G16\_XXX} parsers on a single
output file and collects the results into one struct \code{T}. The file is
read from disk once and the resulting lines are passed to every
sub-parser via their \code{'Lines'} option, instead of each sub-parser
independently re-reading and re-splitting the file.

\begin{argtable}{Field}{Source function}{Description}{2.8cm}{3.2cm}{X}
\code{.charge}    & \code{G09/G16\_charges}       & Mulliken and APT atomic charges (plotting disabled) \\
\code{.energy}    & \code{G09/G16\_energy}        & SCF energy and thermochemistry \\
\code{.structure} & \code{G09/G16\_structure}     & Optimized molecular structure \\
\code{.dipolar}   & \code{G09/G16\_dipole\_polar} & Dipole moment and polarisability \\
\code{.nmodes}    & \code{G09/G16\_nmodes}        & IR/Raman vibrational normal modes; omitted (with a non-blocking warning) if the file has no frequency calculation --- check with \code{isfield(T, 'nmodes')} \\
\code{.spectra}   & \code{G09/G16\_spectra}       & Frequencies, IR/Raman intensities, simulated spectra; omitted under the same condition as \code{.nmodes} \\
\code{.route}     & \code{G16\_route} (G16 only)  & Gaussian route section details \\
\code{.chargemol.charge}, \code{.chargemol.mol} & \code{G16\_charge\_mult} (G16 only) & Total molecular charge, spin multiplicity \\
\end{argtable}

\begin{lstlisting}
T = G16_read_all('zeatin.out');
disp(T.energy)
G16_draw_molecule(T.structure);
\end{lstlisting}

\begin{tbnote}
\textbf{Known discrepancy:} the current \code{G09\_read\_all} does
\emph{not} populate \code{.route} or \code{.chargemol}, unlike
\code{G16\_read\_all}, which retains both. If parity with \code{G16\_read\_all}
is required, add
\begin{lstlisting}
T.route = G09_route(filename, 'Lines', lines);
[c, m]  = G09_charge_mult(filename, 'Lines', lines);
T.chargemol.charge = c;
T.chargemol.mol    = m;
\end{lstlisting}
to \code{G09\_read\_all.m}.
\end{tbnote}

\begin{tbnote}
\textbf{Bug fixed (2026-08-05):} \code{G09\_read\_all}/\code{G16\_read\_all}
(and the Python \code{g16\_read\_all}) used to call \code{nmodes}/\code{spectra}
unconditionally, so a job with no frequency calculation at all (e.g.\ a
TD-DFT-only single point job, \code{td=(singlets,nstates=100)} with no
\code{freq} keyword) crashed the \emph{entire} call, even though every
other field (charges, energy, structure, dipole, route, charge/multiplicity)
would have extracted perfectly fine. This directly contradicted
\code{G09\_write\_report.m}/\code{G16\_write\_report.m} (and their Python
counterpart), which already used \code{isfield}/\code{getattr} checks
expecting \code{.nmodes}/\code{.spectra} to be \emph{optionally absent}.
Fixed by wrapping just those two calls in \code{try}/\code{catch}: on
failure the field is omitted entirely (not left as an empty placeholder)
and a single non-blocking \code{warning(...)} explains why, restoring the
graceful-degradation behaviour the rest of the toolbox already assumed.
\end{tbnote}

% -----------------------------------------------------------------------------
\subsection{G09\_batch\_read\_all / G16\_batch\_read\_all}

\begin{lstlisting}
T = G09_batch_read_all(folder)
T = G09_batch_read_all(folder, 'Recursive', true)
T = G09_batch_read_all(folder, 'WriteReports', true)
T = G09_batch_read_all(folder, 'SaveAs', 'summary.xlsx')

T = G16_batch_read_all(folder)
\end{lstlisting}

Scans \code{folder} for \code{.log}/\code{.out} files (case-insensitive),
runs \code{G09\_read\_all}/\code{G16\_read\_all} (plus
\code{G09\_orbital\_energies}/\code{G16\_orbital\_energies}, for the
HOMO-LUMO gap) on each, and returns one row per file --- useful for
screening many calculations at once (conformers, scans, batches of jobs)
without calling \code{read\_all} file-by-file.

\begin{tbnote}
A file that fails to parse (e.g.\ an incomplete/crashed job, or a file
from the other Gaussian version run through the wrong toolbox) does
\emph{not} stop the batch: its row is filled with \code{NaN} and the
caught error message is recorded in the \code{.Status} column instead.
\end{tbnote}

\begin{argtable}{Option}{Default}{Description}{2.8cm}{2.6cm}{X}
\code{'Recursive'}    & \code{false} & Also scan subfolders \\
\code{'WriteReports'} & \code{false} & Write a \code{G09\_write\_report}/\code{G16\_write\_report} \code{.txt} next to each source file \\
\code{'SaveAs'}       & \code{''}    & Path to save the summary table (\code{.csv} or \code{.xlsx}, inferred from the extension) \\
\end{argtable}

\begin{argtabletwo}{Field}{Description}{3.0cm}{X}
\code{.File}        & Source file path \\
\code{.Natoms}      & Atom count (from \code{.structure.Natoms}) \\
\code{.SCF\_Hartree}& SCF energy \\
\code{.E0\_Hartree} & SCF + ZPE \\
\code{.G\_kJmol}    & Gibbs free energy, kJ/mol \\
\code{.mu\_tot}, \code{.mu\_units} & Dipole magnitude and its units \\
\code{.HOMO\_eV}, \code{.LUMO\_eV}, \code{.Gap\_eV} & Frontier orbital energies and gap, eV \\
\code{.Status}      & \code{'ok'}, or the caught error message if parsing failed \\
\end{argtabletwo}

\begin{lstlisting}
T = G16_batch_read_all('results/', 'WriteReports', true);
T(T.Gap_eV == min(T.Gap_eV), :)   % smallest HOMO-LUMO gap
\end{lstlisting}

\newpage

% =============================================================================
\section{Report generation}
% =============================================================================
\subsection{G09\_write\_report / G16\_write\_report}

\begin{lstlisting}
outfile = G09_write_report(T)
outfile = G09_write_report(T, outfile)

outfile = G16_write_report(T)
outfile = G16_write_report(T, outfile)
\end{lstlisting}

Writes a human-readable text report from a \code{G09\_read\_all}/
\code{G16\_read\_all} struct \code{T}, formatting a summary of every field
present in \code{T} into a plain \code{.txt} file. If \code{outfile} is
omitted, the file is named after the source Gaussian file (e.g.\
\code{'zeatin.out'} $\to$ \code{'zeatin\_report.txt'}) in the current
folder.

\begin{argtable}{Argument}{Default}{Description}{2.4cm}{2.8cm}{X}
\code{T}       & (required)   & Struct returned by \code{G09\_read\_all}/\code{G16\_read\_all} \\
\code{outfile} & \code{<source>\_report.txt} & Output \code{.txt} path \\
\end{argtable}

Returns the path written to. Sections are included only if the
corresponding field is present in \code{T}, in this order: route,
charge/multiplicity, molecular structure (full Cartesian coordinate
table), energetics (with thermochemistry if available), dipole moment and
polarisability (including the dynamic polarisability table, if present),
atomic charges (full per-atom table, sum of charges, dipole-from-charges
overlay if computed), vibrational normal modes (frequency/IR/Raman/reduced
mass/symmetry table), and a summary of the simulated IR/Raman spectra.

\begin{tbnote}
The full broadened-spectrum continuum arrays (\code{T.spectra.x},
\code{.IR\_cont}, \code{.Raman\_cont}) are deliberately \emph{not} dumped
into the report --- only the grid range, FWHM, and whether Raman data is
present are summarised. Access those arrays directly from \code{T.spectra}
for plotting or export.
\end{tbnote}

\begin{lstlisting}
T = G16_read_all('zeatin.out');
G16_write_report(T);                        % -> zeatin_report.txt
G16_write_report(T, 'zeatin_summary.txt');   % explicit output path
\end{lstlisting}

No differences between the G09 and G16 versions beyond the input struct's
origin and the report header ("Gaussian 09"/"Gaussian 16 Calculation
Report").

See also \code{G09\_read\_all}/\code{G16\_read\_all}.

\code{G16parser} is a Python~3 port of the G16 MATLAB toolbox, covering the
same parsing, checkpoint-analysis, and static-visualisation functionality
described in the previous sections. It is installable as a standard
editable package and exposes one \code{g16\_xxx} function per MATLAB
\code{G16\_xxx} counterpart, each returning a \code{Struct} object
(attribute access, e.g.\ \code{mol.xyz}, \code{ch.dipole\_Debye}) instead
of a MATLAB struct.

\begin{lstlisting}
cd G16parser
pip install -e .
\end{lstlisting}

\begin{lstlisting}
import G16parser as g16

mol = g16.g16_structure('molecule.out')
g16.g16_draw_molecule(mol, show_axes=True)

ch = g16.g16_charges('molecule.out', show_dipole=True)

oe = g16.g16_orbital_energies('molecule.out')
g16.g16_draw_orbital(oe)

T = g16.g16_read_all('molecule.out')   # everything in one call, single file read
\end{lstlisting}

\begin{tbnote}
\textbf{Requirements:} Python~$\geq$3.9, \code{numpy}, \code{pandas},
\code{matplotlib} (installed automatically via \code{pip install -e .}).
No Gaussian installation or compiled dependency is needed.
\end{tbnote}

\subsection{Function reference}

Every function mirrors its MATLAB \code{G16\_xxx} counterpart in name,
parameters (in \code{snake\_case}), defaults, and returned fields --- see
the corresponding MATLAB subsection earlier in this manual for full
parameter/field documentation. All extraction functions accept an optional
\code{lines=} parameter (pre-read file lines) mirroring the MATLAB
\code{'Lines'} option; \code{g16\_read\_all} uses this internally to read
the file exactly once.

\begin{argtabletwo}{Function}{Description}{3.4cm}{X}
\code{g16\_au\_convert}       & Converts a physical quantity between atomic units and SI/another common unit (\S\ref{sec:gutil-au2si-python}) \\
\code{g16\_au\_table}         & Prints a formatted reference table of atomic-unit conversion factors (CODATA 2022) \\
\code{g16\_gaussian\_\discretionary{}{}{}field\_\discretionary{}{}{}convert} & Converts between Gaussian's \code{Field=X+N} route keyword and a physical field strength \\
\code{g16\_read\_input}       & Reads a Gaussian \emph{input} file (\code{.gjf}/\code{.com}/\code{.in}): link0 commands, route, title, charge/multiplicity, and starting geometry --- output \code{Struct} is compatible with \code{g16\_draw\_molecule}/\code{g16\_get\_bond\_length} \\
\code{g16\_structure}         & Molecular geometry (\code{symbols}, \code{xyz} as \code{np.ndarray}, atom count) \\
\code{g16\_energy}            & SCF energy and thermochemistry \\
\code{g16\_convergence}       & Geometry-optimisation convergence criteria per step (\code{plot=True} for a matplotlib figure) \\
\code{g16\_charges}           & Mulliken/APT atomic charges, optional dipole overlay and 3D plot \\
\code{g16\_dipole\_polar}     & Dipole moment and (static/dynamic) polarisability \\
\code{g16\_nmodes}            & Vibrational normal-mode displacement vectors \\
\code{g16\_spectra}           & IR/Raman spectra (stick + Lorentzian-broadened continuum) \\
\code{g16\_spectra\_nm}       & Same Lorentzian-broadened IR/Raman spectra as \code{g16\_spectra}, built directly from an already-parsed \code{nm} struct (from \code{g16\_nmodes}, or the \code{.nm} field of \code{g16\_fchk\_read}) --- no \code{.log}/\code{.out} file needed \\
\code{g16\_orbital\_\discretionary{}{}{}energies} & HOMO/LUMO and the full occupied/virtual MO spectrum \\
\code{g16\_charge\_mult}      & Molecular charge and spin multiplicity \\
\code{g16\_route}             & Route section string \\
\code{g16\_available\_data}   & Lists which toolbox quantities are expected to be extractable from a file, by checking its route keywords (e.g.\ ``no \code{freq} keyword $\to$ no nmodes/spectra/thermochemistry'') --- catches this before an error or a NaN \\
\code{g16\_get\_\discretionary{}{}{}bond\_\discretionary{}{}{}length} & Bond-length table (pandas DataFrame) from covalent radii \\
\code{g16\_nbo\_bonds}        & Bond order (single/double/triple) from an actual Gaussian NBO analysis (\code{pop=nbo}), by counting bonding (BD) natural bond orbitals per atom pair \\
\code{g16\_fchk\_read}        & Reads a Gaussian formatted checkpoint (\code{.fchk}) file --- works on both G09 and G16 \code{.fchk} files, since the format does not depend on the Gaussian version \\
\code{g16\_charges\_fchk}     & Visualises charges from a \code{g16\_fchk\_read} struct, without needing the corresponding \code{.log}/\code{.out} file \\
\code{g16\_gaussian\_\discretionary{}{}{}version} & Detects the Gaussian version/revision (works on \code{.fchk} via a sibling \code{.log}/\code{.out}) \\
\code{g16\_hyperpolar}        & Dipole hyperpolarisability ($\beta$) \\
\code{g16\_tddft}             & TD-DFT excited states \\
\code{g16\_read\_all}         & Runs the full extraction set in one call, reading the file only once \\
\code{g16\_restart}           & Generates a \code{.gjf} restart input file from an existing output file \\
\code{g16\_batch\_read\_all}  & Runs \code{g16\_read\_all} (+ \code{g16\_orbital\_energies}) over every \code{.log}/\code{.out} file in a folder and aggregates the key results into one summary DataFrame; per-file failures are recorded, not fatal \\
\code{g16\_draw\_molecule}    & 3D CPK ball-and-stick render (matplotlib \code{mplot3d}) \\
\code{g16\_draw\_mode}        & 3D structure with a vibrational mode's displacement arrows \\
\code{g16\_draw\_deformation} & 3D structure with displacement arrows between two related geometries of the same molecule (e.g.\ field on/off, two conformers, before/after optimisation) \\
\code{g16\_draw\_orbital}     & Orbital energy-level diagram with HOMO--LUMO gap arrow \\
\code{g16\_animate\_mode}     & Exports an MP4 animation of a vibrational mode (requires \code{ffmpeg} on \code{PATH}) \\
\code{g16\_mode\_viewer}      & Interactive Tkinter window to browse/render vibrational modes, sortable by mode number, IR, or Raman intensity, with an "Animate mode (MP4)..." button \\
\code{g16\_list}              & Lists every function in the toolbox with its one-line description (pandas DataFrame) \\
\code{g16\_write\_report}     & Writes a human-readable \code{.txt} summary report from a \code{g16\_read\_all} Struct \\
\end{argtabletwo}

\subsection{Selected signatures}

\begin{lstlisting}
g16_au_convert(value, quantity, direction, target="", verbose=None)
g16_au_table(section="")
g16_gaussian_field_convert(value, direction, unit="au", verbose=True)
g16_read_input(filename)
g16_available_data(filename)
g16_structure(filename, orientation="auto", step="last", lines=None)
g16_energy(filename, step="last", lines=None)
g16_convergence(filename, plot=False, lines=None)
g16_nmodes(filename, section="last", modes=None, lines=None)
g16_spectra(filename, FWHM=10, xmin=0, xmax=4000, dx=1, normalize=False,
            plot=False, section="last", lines=None)
g16_spectra_nm(nm, FWHM=10, xmin=0, xmax=4000, dx=1, normalize=False, plot=False)
g16_charges(filename, type="Mulliken", mode="atom", plot=True,
            atom_scale=0.35, bond_tol=1.30, bond_list=None, font_size=8,
            color_scale="RdBu",
            threshold=0.0, show_dipole=False, dipole_origin="negcharge",
            dipole_scale=1.0, dipole_color=(0, 0.6, 0), dipole_line_width=2.5,
            show_dipole_label=True, dipole_font_size=11, dipole_units="Debye",
            lines=None)
g16_draw_molecule(mol, atom_scale=0.35, bond_tol=1.30, show_labels=True,
                  show_legend=True, title="", bg_color=(0.95, 0.95, 0.95),
                  ax=None, show_axes=False, axes_length=None, bond_list=None)
g16_draw_deformation(mol1, mol2, overlay=False, scale=1, arrow_color=(1.0, 0.4, 0.1),
                     overlay2_color=(0.75, 0.1, 0.1), atom_scale=0.35, bond_tol=1.30,
                     show_labels=False, arrow_threshold=0.05)
g16_animate_mode(mol, nm, mode_idx, filename=None, scale=1.5, flip_sign=False,
                 atom_scale=0.35, bond_tol=1.30, show_labels=False,
                 frames_per_cycle=30, n_cycles=2, fps=20, view=None,
                 dpi=150, progress_callback=None,
                 crop=True, crop_margin=12, show_title=False)
g16_mode_viewer(filename, scale=1.5, atom_scale=0.35, bond_tol=1.30,
                show_labels=False, flip_sign=False, arrow_color=None,
                font_size=10, **extra_opts)
g16_get_bond_length(mol, tolerance=1.15, include_h=True,
                    sort_by="distance", save_as="")
g16_nbo_bonds(filename, section="last", lines=None)
g16_fchk_read(filename, verbose=True)
g16_charges_fchk(mol, ch, mode="atom", plot=True, atom_scale=0.35,
                 bond_tol=1.30, font_size=8, color_scale="RdBu", threshold=0.0)
g16_hyperpolar(filename, units="au", lines=None)
g16_tddft(filename, nstates=math.inf, plot=False, fwhm_ev=0.30, lines=None)
g16_restart(filename, step="last", output="", extra_route="", nproc=0, mem="")
g16_batch_read_all(folder, recursive=False, write_reports=False, save_as="")
g16_write_report(T, outfile=None)
\end{lstlisting}

\subsection{Notes on the port}

\begin{tbnote}
\textbf{Test suite.} A \code{pytest} suite lives in \code{G16parser/tests/}
(\code{pip install -e ".[test]"} then \code{pytest} from the
\code{G16parser/} folder). Most tests need one real Gaussian~16
\code{.out}/\code{.log} file dropped into \code{tests/fixtures/} (any
filename); without one, those tests are skipped rather than failed, so
the suite stays runnable out of the box. \code{tests/fixtures/water.gjf}
is a small hand-written synthetic input file (not from a real
calculation) used for the \code{g16\_read\_input} tests, which do not
depend on the optional real-output fixture.
\end{tbnote}

\begin{tbnote}
\textbf{Continuous integration.} The test suite runs automatically on
every push/pull request to \code{main} via GitHub Actions
(\code{.github/workflows/python-tests.yml}), against Python 3.9 and
3.12 on Ubuntu. Coverage is Python-only: running the MATLAB
\code{G09\_*.m}/\code{G16\_*.m} functions in CI would additionally
require a valid MATLAB licence (MathWorks' \code{matlab-actions/setup-matlab}
GitHub Action needs one even for public repositories), so the MATLAB
side of the toolbox continues to be verified manually.
\end{tbnote}

\begin{tbnote}
\textbf{3D interactivity.} matplotlib's \code{mplot3d} has no MATLAB-style
\code{rotate3d} widget; use the interactive backend's normal mouse/toolbar
controls instead.
\end{tbnote}

\begin{tbnote}
\textbf{Naming exception: \code{g16\_read\_input}.} On the MATLAB side this
function has no version prefix at all (\code{G\_read\_input}, identical
copies in both the \code{G09/} and \code{G16/} folders) because the
Gaussian input file format does not differ between versions. Since only a
\code{G16parser} Python package exists (there is no separate
\code{G09parser}), the port simply follows this package's \code{g16\_}
naming convention like every other function here.
\end{tbnote}

\begin{tbnote}
\textbf{g16\_restart} and \textbf{g16\_batch\_read\_all} port
\code{G16\_restart.m} and \code{G16\_batch\_read\_all.m} only --- there is
no Python equivalent of \code{G09\_restart.m}/\code{G09\_batch\_read\_all.m},
consistent with every other G09-only MATLAB function (e.g.\
\code{G09\_read\_lines}, an internal helper with no G16 counterpart since
\code{G16\_*.m} functions inline the equivalent file-reading logic
instead), since this package only targets Gaussian 16 output. Note that
\code{G09\_fchk\_read}/\code{G09\_charges\_fchk} are \emph{not} G09-only
despite the prefix (identical \code{G16\_} counterparts exist, since the
\code{.fchk} format itself does not depend on the Gaussian version) ---
both were ported to Python (\code{g16\_fchk\_read}/\code{g16\_charges\_fchk})
alongside their \code{G16\_} counterparts.
\end{tbnote}

\begin{tbnote}
\textbf{\code{g16\_animate\_mode}'s frame cropping is implemented
differently from \code{G09\_animate\_mode.m}/\code{G16\_animate\_mode.m}.}
MATLAB captures each rendered frame's pixels directly (\code{getframe})
and crops that array in memory before writing it to \code{VideoWriter}.
matplotlib's \code{FFMpegWriter} offers no equivalent per-frame pixel
hook, so the Python port instead (1) renders the two oscillation-extreme
probe frames to compute the same tight content bounding box (padded by
\code{crop\_margin}, evened for \code{yuv420p}), (2) writes the full,
uncropped video to a temporary file exactly as without \code{crop}, then
(3) runs a single \code{ffmpeg -vf crop=W:H:X:Y} pass over it to produce
the final, cropped file. This relies on nothing beyond \code{ffmpeg}
itself (already a hard requirement for MP4 output at all), avoiding both
a new dependency and any private/internal matplotlib API. Verified
directly: cropped output is consistently smaller (in both dimensions,
always even) than the uncropped video of the same mode, and enabling
\code{show\_title} measurably enlarges the crop's height, matching the
MATLAB behaviour.
\end{tbnote}

\begin{tbnote}
\textbf{Backend selection.} The package forces the \code{TkAgg} matplotlib
backend (if Tk is available and \code{pyplot} has not been imported yet)
instead of the native macOS backend: on some older matplotlib releases,
interactively rotating a 3D plot with the native Cocoa backend can segfault
the whole Python process. Call \code{matplotlib.use(...)} yourself
\emph{before} importing \code{G16parser} if a different backend is needed.
\end{tbnote}

\begin{tbnote}
\textbf{Bug found during porting, MATLAB-side only (fixed 2026-07-13):}
in an early version of \code{G16\_tddft.m}, MATLAB's \code{regexp} silently
dropped the optional trailing \code{<S**2>=} capture group even when it
matched (a MATLAB-specific regex-engine quirk), so \code{td.S2} was always
\code{NaN} in MATLAB even when the tag was present in the file. This
Python port's \code{re}-based parser never shared that bug. Rather than
leaving this as a permanent discrepancy between the two implementations,
the MATLAB-side parser was fixed to extract \code{<S**2>=} with its own
independent \code{regexp} call, instead of a trailing optional group;
both implementations now agree.
\end{tbnote}

\begin{tbnote}
\textbf{Bug found during porting, MATLAB-side only (fixed 2026-08-03):}
the IR/Raman intensity calculation in \code{G09\_fchk\_read.m}/
\code{G16\_fchk\_read.m} re-reshaped the already-correctly-shaped dipole/
polarisability derivative matrices a second time
(\code{reshape(dip\_deriv, N3, 3)'}) before use. Since MATLAB's
\code{reshape} reinterprets the same column-major buffer under the new
shape rather than transposing back to the original layout, this silently
scrambled which derivative value belonged to which atom/displacement,
producing plausible-looking but wrong \code{nm.IR}/\code{nm.Raman}
values --- caught by comparing a \code{.fchk}-derived spectrum against
the same molecule's IR/Raman as printed directly by Gaussian in the
corresponding \code{.out} file. Fixed in both MATLAB files by using the
matrices directly (no second reshape needed). This Python port's
\code{g16\_fchk\_read} never had the bug (written directly with the
correct indexing) --- verified to match the \code{.out}-printed values,
and a regression test (\code{test\_fchk\_ir\_raman\_matches\_out\_ground\_truth})
now guards against it ever being reintroduced.
\end{tbnote}

\begin{tbnote}
\textbf{g16\_spectra\_nm} ports \code{G\_spectra\_nm.m}
(Section~\ref{sec:spectra-nm-ref}), a genuinely version-agnostic MATLAB
function (no \code{G09\_}/\code{G16\_} prefix, like \code{G\_read\_input}):
it builds the same Lorentzian-broadened IR/Raman continuum as
\code{g16\_spectra}, but directly from an already-parsed \code{nm} Struct
(from \code{g16\_nmodes}, or the \code{.nm} field of
\code{g16\_fchk\_read}) instead of reading a \code{.log}/\code{.out}
file --- useful when only a \code{.fchk} file is available. Verified to
match \code{g16\_spectra} called directly on the corresponding
\code{.out} file for the same molecule.
\end{tbnote}

\begin{tbnote}
\textbf{g16\_available\_data} ports \code{G09\_available\_data.m}/
\code{G16\_available\_data.m} (Section~\ref{sec:available-data-ref}).
Unlike the MATLAB \code{G09\_} version, it never omits the
\code{hyperpolar}/\code{tddft} rows, since this package only targets
Gaussian 16 output and both functions always exist here.
\end{tbnote}

\begin{tbnote}
\textbf{Bug fixed (2026-08-05), both MATLAB and Python:} \code{g16\_read\_all}
(and \code{G09\_read\_all.m}/\code{G16\_read\_all.m}) used to call
\code{g16\_nmodes}/\code{g16\_spectra} unconditionally, so a file with no
frequency calculation at all (e.g.\ a TD-DFT-only single point job) raised
\code{ValueError} and crashed the entire call, even though every other
field would have extracted fine. Fixed by wrapping just those two calls
in \code{try}/\code{except}: on failure, \code{nmodes}/\code{spectra} are
set to \code{None} (checked via \code{getattr(T, "nmodes", None)}, matching
how \code{g16\_write\_report} already consumed them) and a single
\code{warnings.warn(...)} explains why.
\end{tbnote}

\begin{tbnote}
\textbf{g16\_mode\_viewer} is a Tkinter rewrite of \code{G16\_modeViewer.m}'s
\code{uifigure} GUI, with the same controls (mode selector, order-by,
per-option redraw, title, save-as PDF/EPS/JPEG, and an "Animate mode
(MP4)..." button wired to \code{g16\_animate\_mode}, using the current
option values and the displayed mode figure's camera orientation). One
simplification: ``target figure'' always means the most recently drawn
mode window, not whichever window the user last clicked on (MATLAB's
\code{groot().CurrentFigure} has no simple Tkinter equivalent). In an
IPython/Spyder console with the Tk graphics backend active, the viewer
detects the running event loop and returns to the prompt immediately
instead of blocking; run as a plain script, it blocks until the viewer
window is closed.

\textbf{Known issue:} \code{g16\_mode\_viewer} can segfault when run from
\emph{inside} Spyder's IPython console (a crash inside ipykernel's own
periodic Tk event-loop pump, not in this toolbox's code) --- confirmed
working correctly when run from a plain terminal/script instead:
\begin{lstlisting}
python3 -c "import G16parser as g16; g16.g16_mode_viewer('file.out')"
\end{lstlisting}
All other functions, including static plots, work fine inside Spyder.
\end{tbnote}

\begin{tbnote}
\textbf{g16\_animate\_mode} mirrors \code{G09\_animate\_mode.m}/
\code{G16\_animate\_mode.m}: it oscillates the molecule along the mode's
displacement vector and saves an MP4 via
\code{matplotlib.animation.FuncAnimation} with \code{FFMpegWriter}.
\textbf{Requires \code{ffmpeg} installed and on \code{PATH}}
(\code{brew install ffmpeg} on macOS, \code{sudo apt install ffmpeg} on
Ubuntu/Debian) --- matplotlib does not bundle a video encoder itself;
without it the call fails with \code{FileNotFoundError: ... 'ffmpeg'} at
the final encoding step (frame generation itself does not need ffmpeg).
Bonds are computed once from the equilibrium geometry via
\code{g16\_get\_bond\_length} and passed to \code{g16\_draw\_molecule}
through its \code{bond\_list} parameter, so they stay fixed across all
frames instead of flickering as instantaneous distances cross
\code{bond\_tol}. By default the animation uses matplotlib's default 3D
view; pass \code{view=(ax.azim, ax.elev)} to start from an
already-rotated figure's orientation instead.
\end{tbnote}

\begin{tbnote}
\textbf{g16\_write\_report} takes the \code{Struct} returned by
\code{g16\_read\_all} and writes the same section-by-section plain-text
report as \code{G09\_write\_report.m}/\code{G16\_write\_report.m}
(route, charge/multiplicity, structure, energetics, dipole/polarisability,
atomic charges, normal modes, spectra summary), using Python
\code{getattr}-based field checks so that sections are included only when
present, exactly mirroring the MATLAB \code{isfield} logic.
\end{tbnote}

\begin{tbnote}
\textbf{Bond order (single/double/triple).} \code{g16\_draw\_molecule}
mirrors \code{G09\_draw\_molecule.m}/\code{G16\_draw\_molecule.m}: for
C-C and C-O pairs, bond order is estimated purely from bond length (any
other element pair, including C-N, is always drawn as a single bond)
and rendered as 1/2/3 parallel lines. This is a geometric estimate, not
Gaussian's own bond-order analysis (e.g.\ Wiberg/NBO indices). For an
actual bond order derived from Gaussian's own NBO analysis, use
\code{g16\_nbo\_bonds} (Section~\ref{sec:nbo-bonds-ref}). The \code{bond\_list}
parameter accepts an optional 3rd column with a pre-computed order, used
by \code{g16\_animate\_mode} (via \code{g16\_get\_bond\_length}) to keep
both bond topology and bond order fixed across all animation frames:
\begin{lstlisting}
mol = g16.g16_structure('molecule_nbo.out')
bt  = g16.g16_nbo_bonds('molecule_nbo.out')   # requires 'pop=nbo' in the source job

bond_list = bt[['Atom1', 'Atom2', 'BondOrder']].to_numpy()
bond_list[:, :2] -= 1   # g16_nbo_bonds uses 1-based Gaussian atom numbers;
                        # g16_draw_molecule's bond_list is 0-based
g16.g16_draw_molecule(mol, bond_list=bond_list)
\end{lstlisting}
\code{g16\_charges} accepts the same \code{bond\_list=} parameter,
passed straight through to its internal \code{g16\_draw\_molecule} call:
\begin{lstlisting}
g16.g16_charges('molecule_nbo.out', bond_list=bond_list, show_dipole=True)
\end{lstlisting}
The C-C double/single threshold is set to 1.36~\AA\ rather than the generic
reference-length midpoint, so symmetric aromatic rings (C-C
$\approx$1.39--1.40~\AA, no real length alternation) render as
all-single instead of all-double. C-N is deliberately excluded from
length-based classification: on asymmetric pyridine-like rings a single
threshold produced one ring C-N bond rendered double and its
chemically-equivalent neighbour rendered single --- see the MATLAB
\code{G09\_draw\_molecule}/\code{G16\_draw\_molecule} subsection earlier
in this manual for the full rationale.
\end{tbnote}

\newpage

% =============================================================================
\section{Worked examples}
% =============================================================================
Four self-contained scripts in \code{G16/} (\code{example\_1.m},
\code{example\_2.m}, \code{example\_3.m}, \code{example\_4.m})
demonstrate the toolbox end to end on a single real DFT calculation:
\code{test\_2.out}, a geometry optimisation of Violacein (G.~Cassone et
al., \emph{Phys. Chem. Chem. Phys.}, doi: 10.1039/d6cp01164k). Each
script only needs \code{G16/} on the MATLAB path and \code{test\_2.out}
in the current folder.

\subsection{example\_1.m --- Structure, vibrational mode, and charges}

Covers \code{G16\_structure}/\code{G16\_draw\_molecule} (3D structure
rendering), \code{G16\_nmodes}/\code{G16\_draw\_mode} (a single
vibrational mode, \#91, the main C=C stretch), and \code{G16\_charges}
(Mulliken charges with a dipole moment overlay), all rendered from the
same camera angle and each exported to its own PDF.

\begin{lstlisting}
%% G16 Toolbox -- Example #1: structure, vibrational mode, and charges
%
%   Demonstrates three core G16_*.m functions on a real DFT calculation:
%     1) G16_structure  + G16_draw_molecule  -- 3D structure rendering
%     2) G16_nmodes     + G16_draw_mode      -- a single vibrational mode
%     3) G16_charges                         -- Mulliken charges + dipole
%
%   Data file: test_2.out -- DFT geometry optimisation of Violacein.
%   Reference: G. Cassone et al., Phys. Chem. Chem. Phys.,
%              doi: 10.1039/d6cp01164k
%
%   Requirements: G16/ on the MATLAB path, test_2.out in the current
%   folder. Running this script creates three PDFs in the current
%   folder: test_2_mol.pdf, test_2_mode_91.pdf, test_2_charges.pdf.

clear
close all
clc

filename = 'test_2.out';

% Camera orientation shared by all three figures below, so the molecule
% is viewed from the same angle in every plot.
view_az = -16.4417;
view_el =  44.0197;

%% 1) Structure -- load the optimised geometry and render it in 3D
mol = G16_structure(filename);

% Create the axes up front so it can be tweaked (view angle, export)
% after G16_draw_molecule has finished drawing into it.
ax1 = axes;
set(ax1, 'Visible', 'off');
G16_draw_molecule(mol, 'Title', 'Violacein', 'Ax', ax1);
set(ax1, 'View', [view_az, view_el]);

% ContentType 'vector' gives a crisp publication figure but is slow for
% a molecule this size; use 'image' instead for a quick low-effort preview.
exportgraphics(ax1, 'test_2_mol.pdf', 'ContentType', 'vector');

%% 2) Vibrational mode -- mode #91 is the main C=C stretch
nm = G16_nmodes(filename);

ax2 = axes;
set(ax2, 'Visible', 'off');
G16_draw_mode(mol, nm, 91);
set(ax2, 'View', [view_az, view_el]);

% Replace only the molecule-name line of the auto-generated two-line
% title; the second line (mode number, frequency, IR/Raman intensity)
% is left untouched.
t = get(gca, 'Title');
t.String(1) = {'Violacein'};
set(gca, 'Title', t);

exportgraphics(gca, 'test_2_mode_91.pdf', 'ContentType', 'vector');

%% 3) Atomic charges -- Mulliken charges with a dipole moment overlay
G16_charges(filename, 'Plot', true, 'ShowDipole', true);
set(gca, 'View', [view_az, view_el]);

t = get(gca, 'Title');
t.String = {'Violacein - Mulliken Charges (atom)'};
set(gca, 'Title', t);

exportgraphics(gca, 'test_2_charges.pdf', 'ContentType', 'vector');

fprintf('\nDone. Figures: structure, mode 91, Mulliken charges.\n');
fprintf('Saved: test_2_mol.pdf, test_2_mode_91.pdf, test_2_charges.pdf\n');
\end{lstlisting}

\subsection{example\_2.m --- Infrared and Raman spectra}

Demonstrates \code{G16\_spectra} two ways: a manual plot built directly
from its raw \code{sp.x}/\code{sp.Raman\_cont} output fields, and the
function's own built-in two-panel Raman+IR auto-plot (\code{'plot',
true}), restyled afterwards with larger, bold axis labels for a
print-ready figure. Both are exported to PDF.

\begin{lstlisting}
%% G16 Toolbox -- Example #2: Infrared and Raman spectra
%
%   Demonstrates G16_spectra on a real DFT calculation, two ways:
%     1) Manual plot -- build a single Raman continuum plot directly
%        from the raw sp.x/sp.Raman_cont fields returned by G16_spectra.
%     2) Built-in plot -- let G16_spectra draw its own two-panel
%        Raman+IR figure ('plot', true), then restyle the axes/labels
%        afterwards (larger, bold text) for a print-ready figure.
%
%   Data file: test_2.out -- DFT geometry optimisation of Violacein.
%   Reference: G. Cassone et al., Phys. Chem. Chem. Phys.,
%              doi: 10.1039/d6cp01164k
%
%   Requirements: G16/ on the MATLAB path, test_2.out in the current
%   folder. Running this script creates two PDFs in the current
%   folder: test_2_raman_manual.pdf, test_2_ir_raman.pdf.

clear; close all; clc
filename = 'test_2.out';
%% 1) Manual plot -- Raman continuum only, from the raw sp fields
sp = G16_spectra(filename);
fig1 = figure('Color', 'white');
plot(sp.x, sp.Raman_cont, 'LineWidth', 2);
xlabel('Wavenumber (cm^{-1})', 'FontWeight', 'bold');
ylabel('Raman activity (A^4 AMU^{-1})', 'FontWeight', 'bold');
set(gca, 'FontSize', 12);
exportgraphics(fig1, 'test_2_raman_manual.pdf', 'ContentType', 'vector');
%% 2) Built-in two-panel plot (Raman + IR), then restyled
% 'FWHM' broadens the Lorentzian lines and 'xmin' crops the noisy
% low-frequency tail; 'plot', true triggers G16_spectra's own
% two-subplot Raman/IR figure (see G16_plot_spectra, local to
% G16_spectra.m).
G16_spectra(filename, 'FWHM', 15, 'xmin', 400, 'plot', true);
fig2 = gcf;
% G16_spectra creates the Raman axes first, then the IR axes below it;
% findobj returns them in reverse creation order, so flip to restore
% [Raman; IR].
axs      = flipud(findobj(fig2, 'Type', 'axes'));
ax_raman = axs(1);
ax_ir    = axs(2);
% Upsize the auto-generated labels (FontSize 10, regular weight) to a
% bolder, larger style for a print-ready figure.
set([ax_raman, ax_ir], 'FontSize', 12);
xlabel(ax_raman, 'Wavenumber (cm^{-1})',         'FontWeight', 'bold', 'FontSize', 12);
ylabel(ax_raman, 'Raman activity (A^4 AMU^{-1})', 'FontWeight', 'bold', 'FontSize', 12);
xlabel(ax_ir,    'Wavenumber (cm^{-1})',         'FontWeight', 'bold', 'FontSize', 12);
ylabel(ax_ir,    'IR intensity (KM mol^{-1})',    'FontWeight', 'bold', 'FontSize', 12);
exportgraphics(fig2, 'test_2_ir_raman.pdf', 'ContentType', 'vector');
fprintf('\nDone. Figures: manual Raman plot, built-in IR+Raman plot.\n');
fprintf('Saved: test_2_raman_manual.pdf, test_2_ir_raman.pdf\n');
\end{lstlisting}

\subsection{example\_3.m --- Energetics}

Demonstrates \code{G16\_orbital\_energies}, \code{G16\_energy}, and
\code{G16\_convergence}, renders the HOMO-LUMO diagram via
\code{G16\_draw\_orbital}, and combines all three structs into a
one-line summary --- a small-scale taste of what
\code{G16\_read\_all}/\code{G16\_write\_report} do at full scale across
every extraction function.

\begin{lstlisting}
%% G16 Toolbox -- Example #3: Energetics
%
%   Demonstrates G16_orbital_energies, G16_energy, and G16_convergence
%   on a real DFT calculation, then combines their results into a
%   one-line summary -- a small taste of what G16_read_all/
%   G16_write_report do at full scale across every extraction function.
%   Also renders the HOMO-LUMO diagram via G16_draw_orbital.
%
%   Data file: test_2.out -- DFT geometry optimisation of Violacein.
%   Reference: G. Cassone et al., Phys. Chem. Chem. Phys.,
%              doi: 10.1039/d6cp01164k
%
%   Requirements: G16/ on the MATLAB path, test_2.out in the current
%   folder. Each function already prints its own formatted summary to
%   the command window (shown below as a reference for what to expect
%   on this file); this script does not repeat those printouts. Running
%   it creates test_2_orbital_diagram.pdf in the current folder.

clear; close all; clc
filename = 'test_2.out';
%% 1) HOMO-LUMO gap -- detailed data in the oe structure
oe = G16_orbital_energies(filename);
% Console output for this file:
%   -- G16_orbital_energies (step 3/3): test_2.out --
%     HOMO = -0.190210 Ha  (-5.1759 eV)
%     LUMO = -0.096890 Ha  (-2.6365 eV)
%     Gap  =  0.093320 Ha  ( 2.5394 eV)

% Orbital energy-level diagram (title defaults to the source filename)
G16_draw_orbital(oe);
exportgraphics(gcf, 'test_2_orbital_diagram.pdf', 'ContentType', 'vector');
%% 2) Energy and thermochemistry -- detailed data in the en structure
en = G16_energy(filename);
% Console output for this file:
%   -- G16_energy: test_2.out --
%     Method : RB3LYP
%     SCF    : -1160.20276011  Ha
%     ZPE    : +0.29182000  Ha   (766.17 kJ/mol)
%     E0+ZPE : -1159.91094000  Ha
%     H      : -1159.88988100  Ha
%     G      : -1159.96045200  Ha
%     S      :  621.4460  J/(mol.K)
%     T = 298.15 K,  P = 1.00000 atm

%% 3) Geometry-optimisation convergence -- detailed data in the cv structure
cv = G16_convergence(filename);
% Console output for this file:
%   G16_convergence: test_2.out
%     Steps read  : 11
%     Converged   : YES  (step 10)
%     Last step   : MaxF=1.00e-06 (thr 1.50e-05)  RMSF=0.00e+00 (thr 1.00e-05)
%                   MaxD=1.24e-04 (thr 6.00e-05)  RMSD=2.50e-05 (thr 4.00e-05)

%% 4) Combine all three into a one-line summary
if cv.converged
    conv_label = sprintf('YES (step %d/%d)', cv.conv_step, cv.Nsteps);
else
    conv_label = sprintf('NO (%d steps read)', cv.Nsteps);
end
fprintf('\nSummary for %s:\n', filename);
fprintf('  HOMO-LUMO gap : %.3f eV\n', oe.gap_eV);
if en.has_thermo
    fprintf('  Gibbs energy  : %.2f kJ/mol\n', en.G_kJ);
else
    fprintf('  Gibbs energy  : n/a (no frequency/thermochemistry job)\n');
end
fprintf('  Converged     : %s\n', conv_label);
\end{lstlisting}

\subsection{example\_4.m --- Recovering and restyling graphics handles}

Shows how to reach back into a figure already drawn by
\code{G16\_draw\_molecule} and restyle every part of it after the fact
(title, axis labels, legend, atom spheres, bond lines, atom-label text),
rather than only being able to set style options at call time.

\begin{tbnote}
\textbf{The key gotcha this example exists to document:} most of the
graphics objects \code{G16\_draw\_molecule} creates (all bond lines, and
every atom sphere after the first drawn for a given element) are created
with \code{'HandleVisibility', 'off'}, so they do not clutter the
legend or turn up in a plain \code{findobj} search. Reaching them
requires \code{findall} instead of \code{findobj}. Title, axis labels,
and the legend itself keep the default (\code{'on'}) visibility, so
either function finds those. The script also re-enables the axes
(\code{axis(ax, 'on')}) to reveal the X/Y/Z labels, which
\code{G16\_draw\_molecule} always creates but hides by calling
\code{axis(ax, 'off')} internally; and separates the atom-label text
objects from the title/axis-label text objects (all four share
\code{Type = 'text'}) with a \code{setdiff}.
\end{tbnote}

\begin{lstlisting}
%% Example 4
%% Recover graphics handles from G16_draw_molecule and restyle them
mol = G16_structure('test_2.out');
G16_draw_molecule(mol, 'Title', 'Violacein');

ax  = gca;
fig = gcf;

%% 1) Title -- always visible (HandleVisibility default 'on')
title(ax, 'Violacein (DFT-optimised)', 'FontSize', 13, ...
      'FontWeight', 'bold', 'Interpreter', 'tex');

%% 2) Axis labels -- G16_draw_molecule calls axis(ax,'off'), so labels
%    exist but are not shown until the axis is switched back on
axis(ax, 'on');
xlabel(ax, 'X (\AA)', 'FontSize', 11, 'Interpreter', 'latex');
ylabel(ax, 'Y (\AA)', 'FontSize', 11, 'Interpreter', 'latex');
zlabel(ax, 'Z (\AA)', 'FontSize', 11, 'Interpreter', 'latex');

%% 3) Legend -- created with default HandleVisibility, findobj works fine
hleg = findobj(fig, 'Type', 'Legend');
set(hleg, 'FontSize', 10, 'Location', 'northeast', 'Box', 'on');

%% 4) Atom spheres (Surface objects) -- use findall: most have
%    HandleVisibility 'off' (only the first atom per element is 'on')
hatoms = findall(ax, 'Type', 'Surface');
set(hatoms, 'FaceAlpha', 0.9, 'EdgeColor', 'none');

%% 5) Bond lines -- all created with HandleVisibility 'off', findall required
hbonds = findall(ax, 'Type', 'Line');
set(hbonds, 'LineWidth', 2.5, 'Color', [0.35 0.35 0.35]);

%% 6) Atom-label text objects only, excluding Title/XLabel/YLabel/ZLabel
%    (these are also Type='text', so must be explicitly subtracted out)
htitle_h = get(ax, 'Title');
hxlab    = get(ax, 'XLabel');
hylab    = get(ax, 'YLabel');
hzlab    = get(ax, 'ZLabel');
halltext = findall(ax, 'Type', 'Text');
hatomlbl = setdiff(halltext, [htitle_h; hxlab; hylab; hzlab]);
set(hatomlbl, 'FontSize', 8, 'FontWeight', 'bold');

%% 7) Background / figure colour
set(ax,  'Color', [1 1 1]);
set(fig, 'Color', [1 1 1]);

exportgraphics(fig, 'test_2_mol_restyled.pdf', 'ContentType', 'vector');
\end{lstlisting}

\newpage

% =============================================================================
\section{References}
% =============================================================================
Background reading for the physics/mathematics behind
\code{MOSurface}'s real-space evaluation: Gaussian-type-orbital point
evaluation (\S\ref{sec:gutil-mo-surface-fn}) and the analytic
McMurchie-Davidson ESP integrals (\S\ref{sec:gutil-esp-surface}).

\begin{enumerate}[nosep,leftmargin=2.2em]
\item \label{ref:Boys1950} S.~F.~Boys, ``Electronic Wave Functions I. A
      General Method of Calculation for the Stationary States of Any
      Molecular System,'' \emph{Proc.\ R.\ Soc.\ Lond.\ A} \textbf{200},
      542--554 (1950). --- Introduces Gaussian-type orbitals for
      molecular electronic structure; origin of the GTO normalization
      constant used throughout \code{G\_draw\_mo\_surface}.
\item \label{ref:Helgaker2000} T.~Helgaker, P.~J\o rgensen, J.~Olsen,
      \emph{Molecular Electronic-Structure Theory}, Wiley (2000). ---
      Standard graduate reference for Gaussian-type-orbital integrals,
      the Gaussian product theorem, and solid harmonics; useful
      background for \S\ref{sec:gutil-mo-surface-fn} and (Ch.~9) the
      McMurchie-Davidson Hermite-Gaussian integral scheme used by
      \code{G\_draw\_esp\_surface} (\S\ref{sec:gutil-esp-surface}).
\item \label{ref:McMurchieDavidson1978} L.~E.~McMurchie, E.~R.~Davidson,
      ``One- and Two-Electron Integrals over Cartesian Gaussian
      Functions,'' \emph{J.\ Comp.\ Phys.}\ \textbf{26}, 218--231 (1978).
      \url{https://doi.org/10.1016/0021-9991(78)90092-X} --- The
      original Hermite-Gaussian expansion/recursion scheme for analytic
      Gaussian integrals; the direct source for the exact electronic-ESP
      evaluation in \code{G\_draw\_esp\_surface}
      (\S\ref{sec:gutil-esp-surface}).
\item \label{ref:SchlegelFrisch1995} H.~B.~Schlegel, M.~J.~Frisch,
      ``Transformation Between Cartesian and Pure Spherical Harmonic
      Gaussians,'' \emph{Int.\ J.\ Quantum Chem.}\ \textbf{54}, 83--87
      (1995). \url{https://doi.org/10.1002/qua.560540202} --- The
      original general formula for Cartesian~$\leftrightarrow$~pure
      spherical-harmonic Gaussian transformation coefficients, from one
      of Gaussian's own lead developers.
\item \label{ref:Ribaldone2024} C.~Ribaldone, J.~K.~Desmarais,
      ``Spherical to Cartesian Coordinates Transformation for Solid
      Harmonics Revisited: Construction of the Hartree Potential,''
      \emph{arXiv:2412.16733} (2024).
      \url{https://arxiv.org/abs/2412.16733} --- Explicit, tabulated
      real-solid-harmonic polynomial coefficients up to $\ell=10$; the
      direct source for the pure D/F/G formulas implemented in
      \code{G\_draw\_mo\_surface}.
\item \label{ref:IOData} T.~Verstraelen \emph{et al.}, \emph{IOData}
      (theochem/iodata), \url{https://github.com/theochem/iodata},
      \url{https://iodata.readthedocs.io}. --- Reference implementation
      and documentation of Gaussian's native \code{.fchk} basis-function
      ordering convention (including the Cartesian $L\!\geq\!4$ shell
      reversal), used to verify the shell-component ordering in
      \code{G\_draw\_mo\_surface}.
\end{enumerate}

\newpage

% =============================================================================
\section{Publication and licensing}
% =============================================================================
\subsection{Declaration of Generative AI and AI-assisted technologies in the writing process}

During the preparation of this work, the author(s) used Claude (Anthropic)
to assist with code drafting, refactoring, and documentation of the
G09/G16 toolbox and its associated user manual. After using this tool, the
author(s) reviewed and edited the content as needed and take full
responsibility for the content of the publication.

\subsection{Note on the Use of AI-Assisted Development Tools}

Portions of the codebase and documentation described in this manual ---
including the \code{G09\_}/\code{G16\_} MATLAB function pairs, the
\code{g16toolbox} Python package, the \code{G09\_modeViewer.m}/
\code{G16\_modeViewer.m} interactive GUIs, and this reference manual ---
were developed with the assistance of Claude (Anthropic), an AI-based
coding tool, used interactively under the direct supervision of the
author. All AI-generated code was tested against reference Gaussian
09/16 output files and reviewed by the author prior to inclusion.
Algorithm design, scientific validation, and correctness of results
remain the sole responsibility of the author.

\subsection{Scope and disclaimer}
\label{sec:scope-disclaimer}

This toolbox is provided as a research and productivity aid for parsing,
analysing, and visualising Gaussian~09/16 output files --- it is not a
validated or certified computational-chemistry package, and its numerical
output has not been benchmarked beyond the specific test cases described
in this manual and its companion \code{G\_Utility} manual (which carries
its own, matching ``Scope and disclaimer'' note). A number of functions
rely on geometric or heuristic approximations rather than a direct read
of a corresponding Gaussian-computed quantity: for instance, bond order in
\code{G09\_draw\_molecule}/\code{G16\_draw\_molecule}
(\S\ref{sec:draw-molecule}) is estimated purely from bond length, not
from an actual Gaussian Wiberg/NBO bond-order analysis, and C--N bonds
are deliberately never classified as double even when geometrically
short, since a single length threshold cannot separate an aromatic C--N
from a genuine C=N. Every such approximation is documented, together
with its known limitations, at the point in this manual (or the
\code{G\_Utility} manual) where the relevant function is described.
Users remain responsible for independently verifying any result before
relying on it for a scientific conclusion or publication --- the same
standard that would be applied to the output of any third-party analysis
tool. This note is a scientific/methodological caveat and does not
replace the legal warranty disclaimer already given in the
\code{LICENSE} file distributed with the source code.

\subsection{Licence}

This software is released under the \textbf{BSD 3-Clause License}, an
OSI-approved open source license accepted by SoftwareX. See the
\code{LICENSE.txt} file distributed with the source code for the full
licence text.

\subsection{Citation}

If this toolbox contributes to published research, please acknowledge it
by citing the source repository:

\begin{quote}
S.~Trusso, \emph{G09/G16 Toolbox: MATLAB and Python libraries for parsing,
analysing and visualising Gaussian~09/16 output files},
\url{https://github.com/nello62/Gaussian09G16_output_Matlab_parser}
(accessed \today).
\end{quote}

\end{document}